% !TeX program = lualatex
\documentclass[12pt,oneside]{book}
\usepackage[utf8]{inputenc}
\usepackage[T1]{fontenc}
\usepackage{geometry}
\geometry{margin=1in}
\usepackage{amsmath,amssymb}
\usepackage{graphicx}
\usepackage{xcolor}
\usepackage{tcolorbox}
\tcbuselibrary{skins,breakable}

% Suit symbols – use Unicode
\newcommand{\spadesuit}{♠}
\newcommand{\heartsuit}{♥}
\newcommand{\clubsuit}{♣}
\newcommand{\diamondsuit}{♦}

% Almanac environment
\newenvironment{almanac}
  {\begin{quote}\small\itshape}
  {\end{quote}}

\begin{document}

% ----------------------------------------------------------------------
% FRONT MATTER
% ----------------------------------------------------------------------
\frontmatter

\title{\Huge Fate's Edge\\[0.3\baselineskip]
\Large The Southern Realms (Akilan)}
\author{Kofi of the Salt Syndicate \\ \small (stamped, witnessed, and slightly stained)}
\date{\today}
\maketitle

\tableofcontents

% ----------------------------------------------------------------------
% PREFACE – Framing device (unreliable trader)
% ----------------------------------------------------------------------
\chapter*{Preface: The Salt Trader’s Ledger}
\addcontentsline{toc}{chapter}{Preface}

\noindent
{\Large
\textit{I am a son of Oshiira. My grandfather weighed grain for the Widow; my mother balanced the canal ledgers. I inherited a brass scale and a healthy suspicion of anyone who does not stamp their receipts three times.}
}

The south is not a mystery. It is a network of trade routes, water rights, and harvest schedules – if you know how to read them. The Oshiiran Ledger Empire keeps the Amaranthine fed. Taharka’s passes guard the monsoon gold. Ngomebe’s walking cities are a wonder, yes, but also a contract: their iron for our grain, renewed every nine years.

And then there are the others. Sekogo’s jungle folk whisper to trees. The Crimson Valley’s silk‑wardens tie promises into cords that tighten when you lie. And Ashaan… Ashaan is a tragedy. A once‑great empire that sold its soul to keep its throne. I do not envy them. I do not forgive them. I simply note their tariffs and move on.

This book is not a tourist’s guide. It is a trader’s inventory. I have marked the places where a misplaced stamp can halt a caravan, where a forgotten water right can start a war, and where a single true story can buy safe passage through a hundred miles of hostile hills.

You want to trade the southern routes? Learn the customs. Respect the spirits. Never refuse the third cup of coffee in Taharka – and never, ever ask a Sidhi freedman about their ancestry.

The ledger is open. The scales are balanced. Let us begin.

\vspace{1em}
\noindent\textit{— Kofi of the Salt Syndicate, Mbaro weigh‑house, stamped and witnessed}

% ----------------------------------------------------------------------
% SAFETY & SENSIBILITY
% ----------------------------------------------------------------------
\chapter*{Safety \& Sensibility}
\addcontentsline{toc}{chapter}{Safety \& Sensibility}

Before play, establish how the table handles sensitive content. Fate's Edge uses three standard tools:

\begin{itemize}
  \item \textbf{Lines}: Content that will not appear at all (e.g., “no on‑screen torture of children”).
  \item \textbf{Veils}: Content that happens off‑screen or fades to black (e.g., “slavery is referenced but not depicted in graphic detail”).
  \item \textbf{X‑Card}: Any player may stop a scene at any time by drawing an X on a card or saying “X.” No questions, no explanations.
\end{itemize}

\subsection*{A Note on Ashaan}
Ashaan is a setting of systemic cruelty: chattel slavery, state‑sanctioned violence, child sacrifice (secret cults), and the erasure of identity. The Veiled Sisters rule through terror; the Black Hand erases dissent; the Breath‑less resistance struggles in shadows.

These themes are presented as \textbf{horror and injustice}, never glamorised. Before playing in Ashaan, use the standard safety tools with particular care. The GM should adjust or omit elements as needed for the table’s comfort. Your safety matters more than canon.

\textit{“The Veil is not a choice. It is survival.”} – Breath‑less agent Tamar

% ----------------------------------------------------------------------
% BEYOND THE DARK CONTINENT – Essay on tropes
% ----------------------------------------------------------------------
\chapter*{Beyond the Dark Continent: A Note on Inspiration}
\addcontentsline{toc}{chapter}{Beyond the Dark Continent}

The Southern Realms (Akilan) draw inspiration from the rich and varied cultures of Saharan, Sahelian, West African, East African, and Southern Africa – but they are not a pastiche of any single real‑world culture. This essay explains how Fate’s Edge approaches these inspirations and deliberately avoids harmful tropes.

\subsection*{No Monolithic “Africa”}
The southern continent is vast, with distinct regions that have little in common. Oshiira’s grain‑ledger bureaucracy owes more to the medieval empires of the Sahel (Ghana, Mali, Songhai) than to any “tribal” stereotype. Ngomebe’s walking cities are technological marvels, not “primitive” structures. Taharka’s highland fortresses echo Ethiopia’s rock‑hewn churches and mountain kingdoms. Sekogo’s jungle societies take cues from the forest kingdoms of West Africa (Benin, Ifẹ, Kongo). Each region has its own economy, religion, and internal conflicts.

\subsection*{Agency, Not Exoticism}
The people of the Southern Realms are not passive recipients of outside influence. They trade, they raid, they negotiate, they resist. The Oshiiran Vermilion Corps audits foreign merchants; Taharkan passes have never been conquered; Ngomebe’s mason‑kin maintain their own engineering partnerships with Aeler dwarves. When the Kahfagian Admiralty encroaches on the Brass Coast, the Sidhi diaspora and native rajas fight back.

\subsection*{No “Noble Savages”}
The beast‑kin of the Valewood are refugees of magical experiments, not “one with nature” caricatures. The Lethai‑ar are sophisticated legal scholars, not mystical forest dwellers. The Aelaerem of the Southern Hills are not “pygmies” – they are a distinct culture with their own coffee rites and knot‑magic, and they trade as equals with Taharkan lords.

\subsection*{Avoiding the “Dark Continent” Label}
The book explicitly names the region “The Southern Realms (Akilan).” \textit{Akilan} is an in‑world term from the old Akilan Empire that preceded the Utaran conquest; it is not a real‑world exonym. No character refers to the south as “dark” or “uncivilised” without that bias being challenged by the text.

\subsection*{Colonialism as Villain}
The Kahfagian Admiralty’s presence on the Brass Coast is depicted as exploitative. The Ashaani slave trade is unequivocally evil. The Fenwood sack of Lence is a war crime. The party is never encouraged to side with oppressors without consequences.

\subsection*{GM Guidance}
\begin{itemize}
  \item Encourage players to create characters from these cultures without relying on stereotypes. Use the almanacs and NPC tables for grounded details.
  \item When a player repeats a harmful trope (e.g., “the jungle is savage”), pause and ask: “Is that what your character believes, or is that the truth of the world?”
  \item Use the safety tools liberally, especially in Ashaan.
\end{itemize}

The Southern Realms are not a dark continent. They are a bright, complicated, and fully human place – and they are ready for your stories.

% ----------------------------------------------------------------------
% HOW TO USE THIS BOOK
% ----------------------------------------------------------------------
\chapter*{How to Use This Book}
\addcontentsline{toc}{chapter}{How to Use This Book}

This book is a companion to the \textit{Fate’s Edge Core Amaranthine Worldbook} and the \textit{System Reference Document (SRD)}. It provides regional content for the Southern Realms (Akilan) using the same card‑driven generator system.

Each regional chapter contains:
\begin{itemize}
  \item A short atmosphere paragraph (in the voice of Kofi, your biased trader – take his opinions with a grain of salt).
  \item A local quote.
  \item An \emph{almanac} – a quick Q\&A about daily life.
  \item Then a full set of generator tables (\spadesuit\ places, \heartsuit\ people, \clubsuit\ complications, \diamondsuit\ rewards), imported from separate files. To use them, draw cards and consult the appropriate rank as described in the Core book’s *Cartomancer’s Compass*.
\end{itemize}

For the core rules (attributes, skills, actions, magic, summoning, travel, etc.), refer to the \textit{Fate’s Edge SRD}. The generators assume you are familiar with the Position, Effect, and Clock mechanics.

\vspace{1em}
\noindent\textit{“The road is a ledger. Every oasis is a credit, every broken wheel a debt.”} – Dusana of the Raven Road

% ----------------------------------------------------------------------
% REGIONAL CHAPTERS
% ----------------------------------------------------------------------
\mainmatter

\chapter{Gateways \& Travel}
\emph{The ledger is open. The scales are balanced. The ninth page is waiting.}

\section*{Travel Modes of the Southern Realms}

\subsection*{River Legs (Oshiira, Crimson Valley)}
Canals, weigh‑houses, and canal mothers. Travel clocks advance when barges queue at locks. A \textbf{Water‑Share Tablet} or a \textbf{Vermilion Corps Seal} buys priority passage.

\subsection*{Desert Legs (Ashaan, Fhara)}
Water is the primary resource. A dry well hazard can end a journey. Travel requires a \textbf{Water‑Skin} and a \textbf{Spirit‑Lamp} to ward off the restless dead. The \textbf{Brass Gate} (Ashaan‑Oshiira border) is fortified with spirit‑buoys that scream when crossed.

\subsection*{Jungle Legs (Sekogo, Crimson Valley)}
Navigation by smell and sound; hidden Leopard Society patrols. A \textbf{Honey‑Jar} or an \textbf{Ukwe Leaf} (witness token) pacifies the spirits of the grove. Never cut a marked tree – the wood remembers.

\subsection*{Highland Legs (Taharka)}
Thin air; breath‑counting; passes that close for storms. A \textbf{Pass‑Stone} or a \textbf{Horse Bond} grants passage. Never refuse the second cup of coffee – it is a test of honour.

\subsection*{Stone Legs (Ngomebe)}
Travel with a walking city requires a \textbf{Keystone Charter} and a \textbf{Stone‑Song}. The city’s Migration Clock advances each day; if it stops, the Hollow walks.

\section*{Gateways to Other Expansions}

\begin{itemize}
  \item \textbf{Amaranthine Core (Kahfagia, Ecktoria, etc.)}: From Oshiira, sail the \textbf{Narrow of Sorrows} north to Kahfagia. From the Brass Coast, the \textbf{Black River} leads to Thepyrgos.
  \item \textbf{Way of Silk (Fhara, Sidhi, Pereshi, Kuvani, Tulkani)}: From Oshiira, the \textbf{Salt Road} crosses the Sahel to Fhara. From Ashaan, the \textbf{Red Desert} leads to Pereshi highlands.
  \item \textbf{Eastern Realms (Sihai, Nihori, Ayokha)}: From Ngomebe, the \textbf{Salt Scar} connects to the eastern passes. From Taharka, the monsoon winds carry ships to Ayokha.
  \item \textbf{Dhahara (Brass Coast, Nine Kingdoms, Himadri)}: The Brass Coast is part of Dhahara – see the Dhahara expansion for details.
\end{itemize}
\newpage
% -------------------- OSHIIRA --------------------
\chapter{Oshiira}
\emph{The Ledger Empire.}

\noindent
Kofi says: \textit{“The Widow’s health is a state secret, and the Vermilion Corps has ears in every weigh‑house. Speak carelessly, and your name may be the next one forgotten – or, worse, audited.”}

\noindent
Oshiira is a confederacy of hundreds of tribes and city‑states, united by hunger – not famine control, but the Grain Covenant: contribute to the granaries, and you will never starve. The capital, Mbaro, is a city of canals, weigh‑houses, and the great Granary Tower. The Widow rules, but the Vermilion Corps guards the ledgers, and the Narrow Captains provoke Kahfagia just to remind the Admiralty who really controls the straits.

\vspace{0.5em}
\noindent\textit{“Never speak of succession before naming the Widow. The empress’s health is a state secret, and the Vermilion Corps has ears in every weigh‑house.”} \\
— Vermilion Corps Auditor

\vspace{1em}
\begin{almanac}
\textbf{What do people eat for breakfast?} Millet porridge with dried fish – the millet from the Sahel flood plains, the fish smoked over canal‑barges. The porridge is always the same. The fish varies by season.

\textbf{What does a child learn before they can walk?} To weigh a grain. Every Oshiiran child is given a balance scale and a single grain of millet. By age five, they know the difference between a northern harvest and a southern one – the difference is the tax.

\textbf{What is the first thing you notice about a stranger?} Their hands. Merchants have ink‑stained fingers – the ink from the Ledger. Clean fingers have never signed a contract. Red stains mean a death warrant.

\textbf{What is the most common crime?} Ledger fraud – but fraud is a career. The Vermilion Corps audits the auditors. The land does not lie.

\textbf{What does no one talk about but everyone knows?} The Widow has not been seen in public for seven years. Her voice is a recording. Her seal is a stamp. The Vermilion Corps is the Widow. The Widow is dead. The empire continues.

\textbf{Where do people go when they die?} To the Granary Tower. The dead are ground into meal and mixed with millet. The grain is shipped across the Amaranthine. The empire does not ask.

\textbf{How do you apologize for a serious wrong?} You submit a correction to the Ledger, witnessed by a Vermilion Corps auditor. The seal glows under moonlight – the Widow’s eye.

\textbf{What is the most important festival?} The Harvest of the Ledger (autumn). The names of the year’s debtors are read aloud. Those who cannot pay are added to the Ledger’s ninth page, sealed with black wax.

\textbf{What is the secret that could destroy a powerful person?} The Ledger’s missing page was burned by the Widow herself. The ash was mixed into the mortar of the Granary Tower. It still glows under moonlight – it spells the name of the heir she killed.
\end{almanac}

\newpage
\section{Oshiira --- ``The Ledger Empire''}
\label{chap:oshiira}

\subsection*{Quick Reference}
\textbf{Tagline:} \textit{The ledger does not forgive. The Widow does not forget. The canals remember every promise.}

\begin{quote}
  \textbf{Vermilion Corps Auditor (Elite):}  
  ``The ledger does not forgive arithmetic errors. A single miscounted bushel can unmake a house, a tribe, an empire. I am not here to be kind. I am here to be correct.''
\end{quote}

\begin{quote}
  \textbf{Canal Mother (Commoner):}  
  ``Water is memory. Every canal remembers the hands that dug it, the barges that passed, the widows who wept at its edge. Control the water, and you control what the empire remembers. The Vermilion Corps controls the grain. We control the thirst.''
\end{quote}

\textbf{Genre / Mood:} Sahelian bureaucracy, grain-house noir, the weight of succession.

\textbf{Starting Location:} A grain-loading pier on the Shoreless Bay, where Oshiiran wheat is being loaded onto Ecktorian barges. The air smells of chaff and opportunity. A Narrow Captain with a scarred throat approaches, her hand resting on a ledger wrapped in oilcloth. ``Kahfagian `pirates' were sighted at dawn. The Widow wants deniable assets. Congratulations — you're now deniable. Find out who's leaking our sailing schedules before the next shipment sails. Or don't. I'm sure the canals will remember you fondly.''

\vspace{2mm}
\begin{tcolorbox}[colback=black!3,colframe=blue!60!black,title={Theme \& Atmosphere}]
Oshiira is a confederacy of hundreds of tribes and city-states, united by hunger. Not famine—\textit{control}. The Grain Ledger is not a book. It is a weapon. The capital, Mbaro, is a city of canals, weigh-houses, and the great Granary Tower — a stone spine that casts a shadow over every table from the Shoreless Bay to the Red Desert's edge. Oshiira is the breadbasket of the Amaranthine; its wheat fills the holds of Ecktorian, Vhasian, and Silkstrand ships. But the breadbasket has a lock. The Widow holds the key.

For centuries, the Widow — the empress of Oshiira — has held the confederacy together through marriage alliances, grain quotas, and the careful, brutal distribution of water rights. But age is catching up. Paranoia is replacing wisdom. The Vermilion Corps audits not only grain but loyalty. They have ears in every weigh-house, eyes in every canal. And somewhere in the Granary Tower, a page of the Ledger has been torn out — a page that names the true heir.

The Vermilion Corps has been searching for that page for nine years. They have executed nine scribes. The ninth scribe's ghost still walks the tower, counting bushels that no longer exist.

\vspace{2mm}
\textbf{What you notice first:} The smell of grain dust and canal water — sweet and rotten at once. The constant scratch of reed pens on parchment, like insects in the walls. The way everyone — from merchant to slave — touches their receipt before speaking of trade. The way the silence falls when a Vermilion Corps auditor enters a room.

\textbf{The one rule that kills newcomers:} Never speak of succession before naming the Widow. The empress's health is a state secret, and the Vermilion Corps has ears in every weigh-house. Speak carelessly, and your name may be the next one erased from the Ledger.
\end{tcolorbox}

\subsection*{Regional Mood: The Weight of Grain}

Power here is measured in bushels, not blood. A tribe's standing depends on its harvest records; a merchant's fortune on the accuracy of his scales. The Vermilion Corps audits without mercy, and the Canal Mothers control the water that turns the mill-wheels — and also the water that drowns. Oshiira fought four wars with Ecktoria to a standstill, though the empire suffered worse. The memory of those wars lingers in every weigh-house and every tariff. But the true tension is internal: the Widow has no clear successor, and her four children — each backed by a different faction — circle the throne like vultures around a dying lion. The Vermilion Corps maintains order, but no one knows which heir the Corps truly serves. Not even the Corps itself. The Ninth Heir — the one who was never supposed to exist — watches from the shadows, waiting for the ledger to crack.

\subsection*{Prominent NPCs of Oshiira}
\label{sec:oshiira-npcs}

\begin{longtable}{|p{0.2\textwidth}|p{0.25\textwidth}|p{0.3\textwidth}|p{0.15\textwidth}|}
\hline
\textbf{Name} & \textbf{Role/Title} & \textbf{Motivation} & \textbf{Rumoured Quirk} \\
\hline
Mothers of the Vermilion Eye & Vermilion Corps auditor & Find every miscount, collect every debt & She counts her own breaths to ensure she is not lying. Her breath count never matches mine. \\
\hline
Lady Tano & Canal Mother of the eastern reach & Keep the water flowing and her family fed & She has never seen the sea. She refuses to start. She says the sea is just a canal that lost its nerve. \\
\hline
Captain Sefu & Narrow Captain, Oshiiran navy & Provoke Kahfagia into breaking the peace & He names his cannons after his ex-wives. The newest one is named ``Regret.'' \\
\hline
Merchant Kofi & Salt Syndicate factor & Control the desert caravan routes & He tastes salt before every deal. The salt tells him if the seller is lying. He has not been wrong in thirty years. \\
\hline
Widow's Scribe & Granary Tower archivist & Protect the Ledger at any cost & He has not left the tower in forty years. The Ledger is his family. His shadow has a different posture. \\
\hline
The Vermilion Hand & Internal enforcement (rumoured) & A ghost who audits oath-breakers & Said to appear when a debt is nine years overdue. Those who see her never owe again. They also never speak again. \\
\hline
\end{longtable}

\subsection*{Names of Oshiira}
\label{sec:oshiira-names}

Inspired by the Sahelian and West African traditions: melodic, with repeated syllables and clan prefixes. Surnames often denote a trade, a water source, or an ancestor's deed.

\begin{longtable}{|p{0.25\textwidth}|p{0.25\textwidth}|p{0.2\textwidth}|p{0.2\textwidth}|}
\hline
\textbf{First Names (Male)} & \textbf{First Names (Female)} & \textbf{Surnames} & \textbf{Epithets} \\
\hline
Kofi & Adwoa & Keita & \textit{the Ledger-Keeper} \\
\hline
Kwame & Amara & Diallo & \textit{Canal-Born} \\
\hline
Sefu & Nia & Toure & \textit{Weigh-Master} \\
\hline
Chidi & Zuri & Sow & \textit{Ninth Daughter} \\
\hline
Obi & Imani & Bangura & \textit{the Widow's Eye} \\
\hline
Idriss & Fatou & Kamara & \textit{Salt-Tongue} \\
\hline
\end{longtable}

\paragraph*{(Canal/Weigh-House/Granary)} Mudflat weigh-station where grain sacks pile under reed awnings and the scales squeal like guilty witnesses. Canal junction where three barges queue for a single pilot's writ and the pilots queue for a single bribe.

\section*{Spades --- Places}
\label{sec:oshiira-places}
\begin{enumerate}
\setcounter{enumi}{1}
\item \textbf{Mudflat Weigh-Station} --- Grain sacks piled under reed awnings.
  \hfill \textit{The scales are bronze. They squeal when a false weight is placed. The squeal is a confession.} `[TRADE]`
\item \textbf{Granary Tower (Cracked Ledger Door)} --- A door sealed with wax and warning.
  \hfill \textit{The wax is the Widow's own seal. Break it, and the Vermilion Corps will know. The door has been broken nine times. The wax is always fresh.} `[AUTHORITY]`
\item \textbf{Canal Junction} --- Three barges, one pilot, one writ, many arguments.
  \hfill \textit{The pilot's writ is stamped with a watermark that shifts in sunlight. The watermark is a face. The face changes with the tide.} `[TRAVEL]`
\item \textbf{The Widow's Weigh-House} --- Bronze scales, silent clerks, no windows.
  \hfill \textit{The clerks do not speak. They tap their fingers. Each tap is a number. The numbers spell a name. The name is yours.} `[JUSTICE]`
\item \textbf{Grain Exchange Floor} --- Traders shout in couplets; bids are poems.
  \hfill \textit{A bad rhyme loses your entire stake. A good one doubles it. The best poem buys the poet a death warrant. The Widow hates competition.} `[TRADE]`
\item \textbf{Canal Mother's Court} --- A shaded bench where water rights are assigned.
  \hfill \textit{The bench is carved from a single log. The log was a bridge she ordered burned. The bridge was full of children. She says it was necessary.} `[JUSTICE]`
\item \textbf{Salt Syndicate Counting House} --- Desert caravans pay their tolls here.
  \hfill \textit{The floor is salt. Sweep it, and you are cleaning a debt. The salt is grey. It tastes of old tears and older bargains.} `[TRADE]`
\item \textbf{Shoreless Bay Grain Piers} --- Wooden docks stretching into shallow water.
  \hfill \textit{At low tide, you can see the bones of ships that tried to bypass the tariffs. The bones are not all wood.} `[TRAVEL]`
\item \textbf{Narrow Captains' Mess} --- A longhouse where naval officers drink and scheme.
  \hfill \textit{The toasts are in code. ``To the Widow'' means ``prepare for war.'' ``To the canals'' means ``hide the contraband.'' No one toasts to the Ninth Heir. The silence is loud.} `[POLITICS]`
\item[J] \textbf{Kahfagian Cold War Station} --- A lighthouse that is also a spy post.
  \hfill \textit{The lamp flashes in a sequence that changes weekly. The codebook is in cipher. The cipher is a poem about a drowned sailor. The sailor was a spy. The spy was never caught.} `[AUTHORITY]`
\item[Q] \textbf{Mothers' Stepwell} --- A deep well with stone steps leading down.
  \hfill \textit{Each step is a different mother's name. The ninth step is empty. The empty step is for the mother who has not yet lost a child. She will.} `[RITUAL]`
\item[K] \textbf{Granary Tower Top} --- A wind-swept platform with a single brass bell.
  \hfill \textit{The bell rings only when the Widow dies. It has not rung in forty years. The bell-ringer has not slept in forty years. He is waiting.} `[OMEN]`
\item[A] \textbf{The Ledger Vault} --- A subterranean archive beneath the Granary.
  \hfill \textit{The air is dry. The ledgers are on chains. The chains have teeth. The teeth are for anyone who reads without permission. The ledgers bite back.} `[LEGEND]`
\end{enumerate}

\paragraph*{(Clerk/Mother/Captain)} Grain inspector with a brass stamp and a grudge; his stamp leaves a mark that glows under moonlight. He knows where you slept last night. Canal Mother bargaining over water rights; her sons stand behind her, counting your breaths. Vermilion Corps auditor with ink on fingers and a ledger that never balances.

\section*{Hearts --- People \& Factions}
\label{sec:oshiira-people}
\begin{enumerate}
\setcounter{enumi}{1}
\item \textbf{Grain Inspector} --- Brass stamp, iron will, no sense of humour.
  \hfill \textit{His stamp leaves a mark that glows under moonlight. He knows. He also knows why you flinched.} `[AUTHORITY]`
\item \textbf{Canal Mother} --- Matriarch who controls the sluice gates.
  \hfill \textit{She has three sons. Each runs a different canal. She plays them against each other. The loser gets the dry canal. The dry canal is a grave.} `[COMMUNITY]`
\item \textbf{Vermilion Corps Auditor} --- Ink-stained fingers, a ledger that never closes.
  \hfill \textit{She audits the auditors. Her own books are forged — perfection. She has not slept in nine years. Sleep is a miscount.} `[JUSTICE]`
\item \textbf{Narrow Captain} --- Chafes at the cold war pact. Wants blood.
  \hfill \textit{He keeps a map of Kahfagian ports. Each has a red X. He adds one per season. The newest X is over a lighthouse. The lighthouse keeper was his brother.} `[POLITICS]`
\item \textbf{Salt Syndicate Merchant} --- Offers loans with interest paid in favours.
  \hfill \textit{His interest rate is written in disappearing ink. You never know the true cost. The cost is always higher than you think.} `[TRADE]`
\item \textbf{Weigh-House Clerk} --- Silent, watchful, never blinks.
  \hfill \textit{He has not spoken in ten years. He communicates by stamping receipts. The stamps have a rhythm. The rhythm is a prayer. The prayer is for death.} `[JUSTICE]`
\item \textbf{Tribal Chief} --- Caught between the Widow's law and ancestral custom.
  \hfill \textit{He keeps two ledgers: one for Mbaro, one for his ancestors. The ancestors are winning.} `[COMMUNITY]`
\item \textbf{Barge Pilot} --- Knows every canal, every shortcut, every bribe.
  \hfill \textit{His boat is faster than it should be. The hull is lined with contraband. The contraband is not for sale. It is for leverage.} `[TRAVEL]`
\item \textbf{Spice Merchant} --- Sells cardamom and secrets in equal measure.
  \hfill \textit{Her spice boxes have false bottoms. The bottoms contain intelligence. The intelligence is always true. The truth is always painful.} `[TRADE]`
\item[J] \textbf{Widow's Scribe} --- The empress's personal archivist.
  \hfill \textit{He writes in a script that only the Widow can read. He has forgotten his own name. The script remembers it for him.} `[AUTHORITY]`
\item[Q] \textbf{Vermilion Hand} --- A rumoured figure in red gloves.
  \hfill \textit{She appears when a debt is nine years overdue. Her arrival is a death sentence. Her gloves are red because they never dry.} `[MYSTERY]`
\item[K] \textbf{The Widow (Empress)} --- Dying, paranoid, still counting.
  \hfill \textit{She has not slept in three years. The Ledger keeps her awake. The Ledger is also the only thing keeping her alive.} `[AUTHORITY]`
\item[A] \textbf{The Ninth Heir} --- A child no one acknowledges, raised in secret.
  \hfill \textit{She was taught to read ledgers before she could walk. She is the Widow's only hope. She is also the Vermilion Corps' only fear.} `[LEGEND]`
\end{enumerate}

\paragraph*{(Dust/Locust/Succession)} Dust from the Ashaani desert coats the grain; inspector halts the entire caravan; the dust moves against the wind. Locust swarm devours exposed supplies; the locusts spell a name in the sky. The name is yours.

\section*{Clubs --- Complications/Threats}
\label{sec:oshiira-complications}
\begin{enumerate}
\setcounter{enumi}{1}
\item \textbf{Desert Dust} --- Red sand from Ashaan coats the grain; inspector halts caravan. `[WEATHER]`
  \hfill \textit{The dust smells of old magic. Touch it, and a spirit may notice you. The spirit is not from Ashaan. The spirit is from the grain's ghost.}
\item \textbf{Locust Swarm} --- A chittering cloud descends; Supply –1. `[DISASTER]`
  \hfill \textit{The locusts leave behind a single word written in their bodies: ``Pay.'' The word is in the Widow's own handwriting.}
\item \textbf{Cold War Provocation} --- A Kahfagian ``pirate'' is sighted in Oshiiran waters. `[POLITICS]`
  \hfill \textit{The pirate flies a flag of truce. The truce is a lie. The flag is stitched from a Vermilion Corps uniform.}
\item \textbf{Widow's Health Fails} --- Succession rumours spike; trade freezes. `[POLITICS]`
  \hfill \textit{The Grain Exchange floor goes silent. No one wants to be on the wrong side. The wrong side is the side that bet first.}
\item \textbf{Audit In Progress} --- Vermilion Corps seals the weigh-house; no one leaves. `[JUSTICE]`
  \hfill \textit{The auditor asks one question: ``Where is the ninth bushel?'' There is no ninth bushel. The question is a test.}
\item \textbf{Canal Blockage} --- A barge sinks at a junction; water backs up. `[TRAVEL]`
  \hfill \textit{The barge was sunk deliberately. A Canal Mother is making a point. The point is drowning.}
\item \textbf{False Receipt} --- Your grain receipt is a forgery. Prove otherwise. `[TRADE]`
  \hfill \textit{The forger used real paper. The ink is the only difference. The ink is the forger's blood. The forger is a ghost.}
\item \textbf{Salt Syndicate Call} --- Your loan is due. In favours. Now. `[DEBT]`
  \hfill \textit{The favour is to poison a rival's well. Refuse, and your own well is next. The rival is your ally. The Syndicate knows.}
\item \textbf{Narrow Captain's Gambit} --- A captain offers you a bribe to provoke Kahfagia. `[POLITICS]`
  \hfill \textit{The bribe is a ship. The ship is stolen. You are the decoy. The captain's real target is the Widow's own fleet.}
\item[J] \textbf{Tribal Feud} --- Two tribes claim the same harvest; you are the witness. `[COMMUNITY]`
  \hfill \textit{Both ledgers are correct. The error is in the Widow's own records. The Widow does not make errors. The error is a trap.}
\item[Q] \textbf{Vermilion Hand Appears} --- A figure in red gloves sits at your table. `[MYSTERY]`
  \hfill \textit{She asks for a name. Give a name, or she will take yours. The name she wants is the Ninth Heir's. She knows you don't know it.}
\item[K] \textbf{Widow's Summons} --- The empress demands your presence at the Granary. `[AUTHORITY]`
  \hfill \textit{The summons is written in disappearing ink. It has already begun to fade. The ink is the Widow's blood. She is not well.}
\item[A] \textbf{The Ninth Bushel} --- A grain shipment is short by one bushel. The empire demands blood. `[OMEN]`
  \hfill \textit{The missing bushel is a metaphor. The empire wants a scapegoat. The scapegoat is you. The metaphor is a knife.}
\end{enumerate}

\paragraph*{(Stamp/Token/Priority)} Receipt for stolen grain (blackmail material — the grain was stolen from the Widow's own stores). Barge priority token (skip the queue — the token is a dried finger). Weigh-house seal (legal immunity for one audit — the seal is warm, and it has a pulse).

\section*{Diamonds --- Rewards/Leverage}
\label{sec:oshiira-rewards}
\begin{enumerate}
\setcounter{enumi}{1}
\item \textbf{Grain Receipt (Stolen)} --- Blackmail material against a merchant. `[SOCIAL]`
  \hfill \textit{The receipt is stained with wine. The merchant's daughter was drunk. The daughter is the Ninth Heir.}
\item \textbf{Barge Priority Token} --- Skip the queue at any canal junction (once). `[TRAVEL]`
  \hfill \textit{The token is a clay disc. It breaks after one use. Intentionally. The break reveals a map. The map is to a hidden canal.}
\item \textbf{Weigh-House Seal} --- Legal immunity from one audit. `[JUSTICE]`
  \hfill \textit{The seal is warm. It cools when an auditor is near. It burns when the auditor is lying. It has burned nine times today.}
\item \textbf{Treaty Seal (Cold War)} --- One use to demand a parley with Kahfagia. `[POLITICS]`
  \hfill \textit{The seal is half of a pair. The other half is held by a Kahfagian admiral. The admiral is dead. His ghost holds the other half.}
\item \textbf{Canal Mother's Favour} --- One water-right reassignment. `[COMMUNITY]`
  \hfill \textit{The favour is a key. The key opens a sluice gate. The gate controls a district's water. The district is your enemy's.}
\item \textbf{Vermilion Corps Pass} --- Safe passage through any audit checkpoint. `[AUTHORITY]`
  \hfill \textit{The pass is a single red feather. No one questions it. The feather came from a bird that only eats the eyes of the dead.}
\item \textbf{Salt Syndicate Loan Waiver} --- One debt forgiven, no questions asked. `[TRADE]`
  \hfill \textit{The waiver is a pinch of salt in a leather pouch. Swallow it to activate. The salt tastes of tears. Whose tears?}
\item \textbf{Narrow Captain's Letter of Marque} --- Legal piracy against Kahfagia. `[COMBAT]`
  \hfill \textit{The letter is signed by the Widow. Her signature is a single dot. The dot is a period. The period ends a sentence. The sentence is your life.}
\item \textbf{Tribal Chief's Witness} --- One oath witnessed by a neutral chief. `[SOCIAL]`
  \hfill \textit{The witness is a scar on your arm. The scar is the chief's signature. The signature reads ``wait.''}
\item[J] \textbf{Widow's Cipher} --- Read one intercepted Kahfagian dispatch. `[POLITICS]`
  \hfill \textit{The cipher is a poem. The poem is about a drowned sailor. The sailor was a spy. The spy was you, in a past life you don't remember.}
\item[Q] \textbf{Vermilion Hand's Silence} --- One crime erased from the Ledger. `[MYSTERY]`
  \hfill \textit{The hand leaves a red thumbprint on the page. The page then burns. The smoke smells of cinnamon. The cinnamon is a warning.}
\item[K] \textbf{Heir's Trust} --- The ninth heir owes you a secret. `[LEGEND]`
  \hfill \textit{The secret is a name. The name is the Widow's true successor. The name is yours.}
\item[A] \textbf{The Ledger Page} --- A single page from the Granary's master ledger. `[LEGEND]`
  \hfill \textit{The page lists every debt the empire has forgotten. The debts are still valid. The interest is compounding. The interest is measured in lives.}
\end{enumerate}

\section*{Quick use notes}
\label{sec:oshiira-quick-use}
\begin{itemize}
\item Draw until you have all four suits: Spade = place, Heart = actor, Club = pressure, Diamond = leverage. Highest rank sets the main Clock (2–5 → 4, 6–10 → 6, J/Q/K → 8, A → 10).
\item Diamonds are codified outcomes (receipts/tokens/seals) that change position rather than call for a roll.
\item If any A appears, echo \textbf{grain-omens}—scales that tip without weight, ledgers that rustle in still air, the Widow's cough echoing from the tower, a single grain of wheat that weeps.
\end{itemize}

\subsection*{Additional Features (Cultural Mechanics)}
\label{sec:oshiira-mechanics}

\begin{itemize}
\item \textbf{The Grain Ledger:} Every harvest, every debt, every tariff is recorded. A single page can make or break a family. The Ledger is kept in the Granary Tower, guarded by the Vermilion Corps and the Widow's own magic — and by something else, something that lives in the walls and counts in the dark. Once per session, a character may spend a Boon to consult a copy of the Ledger; the GM must answer one question about grain shipments, water rights, or tax records truthfully (though the answer may be a riddle, and the riddle may have teeth).

\item \textbf{The Widow's Peace:} A fragile truce with Kahfagia, maintained by tariffs, spies, and the occasional deniable raid. Break it, and you break the empire. The party may be asked to perform a deniable raid — success earns a Letter of Marque; failure means the Widow will deny all knowledge. The Widow's denial is written in disappearing ink. Your name will fade with it.

\item \textbf{Canal Law:} Water rights are the second currency. The Canal Mothers decide who drinks, who farms, and who drowns. A character who holds a water right (a clay tablet) may reroute a canal for one scene, flooding or draining a district. Using this right ticks the \textbf{Resentment Clock [4]} for that district; when full, the locals rise up. The rising is polite. The drowning is not.

\item \textbf{Succession Clocks:} The Widow's four children each have a \textbf{Claim Clock [6]}. Ticks when the party aids that heir, spreads a rumour, or assassinates a rival. When any heir's clock fills, they make a public bid for the throne. The party can accelerate or decelerate these clocks through their actions. The Widow herself has a \textbf{Paranoia Clock [8]}; when it fills, she orders a purge — the Vermilion Corps arrests anyone suspected of disloyalty. The purge is always correct. The Vermilion Corps makes sure of it.

\item \textbf{Local Patrons (Syncretic Names):}
  \begin{itemize}
    \item \textbf{The Mother of Scales} — Mykkiel. \textit{The ledger as scripture; a miscount is a sin. The penance is a death.}
    \item \textbf{The Green Keeper} — The Witness. \textit{The flame that shows the truth in grain dust. The truth is always a debt.}
    \item \textbf{The Canal Father} — The Traveler. \textit{The water that connects all tribes; the right of passage. The passage is never free.}
    \item \textbf{The Ninth Hand} — Khemesh. \textit{The pressure of debt; the weight of the Ledger. The weight is crushing.}
  \end{itemize}
\end{itemize}

\subsection*{Lore Echoes (Cross-Expansion Hooks)}
\begin{itemize}
  \item \textbf{The Widow's Ninth Succession (Unfinished Ledger):} The Widow is dying. Her four children vie for the throne. One is a half-breed, one is a spirit, one is a puppet, and one is already dead. The Vermilion Corps has a secret ledger listing which heir each faction supports. The party has been offered a barony to escort the heir of their choice to Mbaro. The barony is a swamp. The heir does not know this. The swamp knows. The swamp is hungry.

  \item \textbf{The Salt Syndicate's Ashaani Contacts (The Veil):} The Salt Syndicate trades with Ashaan across the Red Desert. They smuggle more than salt — they smuggle spirit-bottles, forbidden texts, and escaped slaves. The Veiled Sisters tolerate the trade because the Syndicate also carries their spies. The party may be hired as unwitting couriers; the cargo is a living memory. The memory is of a debt the Widow has not yet collected.

  \item \textbf{The Narrow Captain's Kahfagian Lover (The Flag of Storms):} Captain Sefu has a secret: his ex-wife is a Kahfagian admiral. They communicate via coded lantern signals across the Narrow of Sorrows. Their correspondence is a map to a peace treaty — or a declaration of war. The party is asked to deliver a letter. The letter is sealed with a kiss. The kiss is poisoned. The poison is a metaphor for trust.

  \item \textbf{The Granary Tower's Gnome Engineers (Aeler Influence):} The Granary Tower's counterweights were designed by Aeler vent-priors. The breath-bonds that keep the ledgers safe are of Aeler make. If the Aeler ever withdrew their engineers, the tower would collapse within a season. A faction within the Vermilion Corps wants to replace the Aeler with native engineers. The party is hired to steal the blueprints. The blueprints are written in breath. The breath is the engineers' own. They will want it back.
\end{itemize}

\subsection*{Oshiira Superstitions (burn for +1 die once)}
\begin{itemize}
  \item \textbf{Bread to the Ledger:} Crumble a crust onto a weigh-house scale. Ignore the first \emph{Desert Dust} penalty. `[SUPERSTITION]` The scale will dip. The dip is a blessing. The blessing is measured in crumbs.

  \item \textbf{Knot the Quill:} Tie a single knot in an audit quill. Cancel \emph{False Receipt} or take –1 die on your next negotiation (your choice). `[SUPERSTITION]` The knot will bleed. The blood is ink. The ink is your signature.

  \item \textbf{Name to the Canal:} Whisper a dead Canal Mother's name into the water. Gain +1 Effect when bargaining for water rights, but mark \emph{Obligation} to the Widow. `[SUPERSTITION]` The water will ripple. The ripple will spell a warning. The warning is always the same: \textit{``She is watching.''}
\end{itemize}

\subsection*{Thematic SB Spend Table}
\label{sec:oshiira-sb}

\paragraph{Minor Complications (1 SB)}
\begin{itemize}
\item \textbf{Exposure:} Your actions draw unwanted attention from \textbf{Grain Inspectors or Canal Mothers}. The attention is polite. The politeness is a threat.
\item \textbf{Noise:} Sounds of your actions alert nearby \textbf{weigh-house clerks or barge pilots}. The clerks will not move. The pilots will.
\item \textbf{Trace:} Evidence of your passage marks your route for \textbf{auditors or Syndicate debt-collectors}. The evidence is a single grain of wheat. The grain is counting.
\item \textbf{Delay:} A brief but meaningful setback costs you \textbf{time or favourable tariffs}. The setback is a locked canal gate. The key is in Mbaro. You are not.
\item \textbf{Supply Strain:} Mark +1 segment on a relevant \textbf{resource clock}. The resource is grain. The grain is rotting. The rot is spreading.
\end{itemize}

\paragraph{Moderate Setbacks (2 SB)}
\begin{itemize}
\item \textbf{Alarm Raised:} \textbf{Local Vermilion Corps or Narrow Captain} becomes aware and begins responding. The response is a single question. The question is always the same: ``Where is the ninth bushel?''
\item \textbf{Position Lost:} You lose advantageous ground/cover/stealth due to \textbf{canal blockage or dust storm}. The dust is red. The red is blood. The blood is not yours. Not yet.
\item \textbf{Foe Appears:} A \textbf{rival merchant or tribal chief} arrives on scene. They are smiling. The smile is a receipt. The receipt is for your life.
\item \textbf{Gear Trouble:} A piece of equipment becomes \textbf{Compromised/Neglected}. The equipment is a scale. The scale now lies. The lie is the truth.
\item \textbf{Lock/Barrier:} A simple obstacle now requires a test to overcome. The obstacle is a ledger page. The test is arithmetic. The arithmetic is rigged.
\end{itemize}

\paragraph{Serious Trouble (3 SB)}
\begin{itemize}
\item \textbf{Reinforcements:} Additional \textbf{auditors, guards, or Syndicate enforcers} arrive. They are silent. Their silence is a verdict.
\item \textbf{Key Gear Breaks:} A crucial tool/weapon becomes temporarily unusable. The tool was a quill. The quill was a knife. The knife was a promise.
\item \textbf{Major Twist:} The situation fundamentally changes—\textbf{Widow's health fails/oath invoked/audit declared}. The twist is a page torn from the Ledger. The page is blank. The blank is a threat.
\item \textbf{Rail Tick:} Advance a relevant campaign/front clock by 1 segment. The clock is the Widow's heartbeat. It is slowing.
\item \textbf{Condition Applied:} Mark \textbf{Fatigue 1/Harm 1/Condition} appropriate to fiction. The condition is \textit{Indebted}. The debt is unspecified. The creditor is watching.
\end{itemize}

\paragraph{Major Turns (4+ SB)}
\begin{itemize}
\item \textbf{Trap Springs:} A prepared danger activates with full effect. The trap is a weighing scale. The scale is balanced. The balance is a lie.
\item \textbf{Authority Arrival:} \textbf{Widow, Vermilion Hand, or Ninth Heir} intervenes. The intervention is a single word. The word is ``count.''
\item \textbf{Scene Shift:} The environment changes dramatically—\textbf{canal breaches/ledger catches fire/succession crisis begins}. The fire is green. The green is the Widow's signature colour. The colour is a curse.
\item \textbf{Patron Omen:} Divine/arcane forces take notice—\textbf{omen appears/blessing lost/curse manifests}. The omen is a bushel of wheat. The wheat is screaming. The scream is a number.
\item \textbf{Narrative Pivot:} The story takes an unexpected turn that reframes objectives. The turn is a page. The page is the one torn from the Ledger. The page was never missing. It was waiting.
\end{itemize}

\paragraph{Region-Specific SB Options}
\begin{itemize}
\item \textbf{Oshiira (Grain Law):} Scales tip without cause, ledgers rustle, the Widow's cough echoes from the tower. The cough is a code. The code is a warning.
\item \textbf{Oshiira (Cold War):} Kahfagian flags appear on the horizon, cipher books fall open to the wrong page, spies exchange glances. The glances are receipts. The receipts are for betrayals not yet committed.
\item \textbf{Oshiira (Canal Law):} Sluice gates open unbidden, water rises, a Canal Mother's son draws a knife. The knife is for the heir who tried to buy his mother's vote. The water will wash away the blood. The water always washes away the blood.
\end{itemize}

\subsection*{Pursuit Ladder (Near / Mid / Far)}
\begin{itemize}
\item \textbf{Advance/Withdraw (Wits + Stealth):} On a hit, shift one band. On a strong hit, also impose \emph{False Trail}.
\item \textbf{Cut the Bribe (Presence + Sway):} On a hit, freeze range for one exchange; if a weigh-house is under your seal, +1 die.
\item \textbf{Weather the Audit (Resolve + Lore):} On a hit in \emph{Vermilion Inspection}, you pick who testifies; on a miss, the auditor does.
\end{itemize}

\subsection*{Boss Seeds (Grain \& Ledger)}
\begin{itemize}
\item \textbf{The Ledger-Keeper (Nine Volumes)} — \emph{Phases}: Audit Thresholds $\rightarrow$ Orchestrated Famine $\rightarrow$ Succession Coup. \emph{Cracks}: cut a lineage page; reveal swapped seal; survive a staged grain shortage.
\item \textbf{The Salt Mother (Brine-Syndicate)} — \emph{Phases}: Loans Called Mercy $\rightarrow$ Procession at Weigh-House $\rightarrow$ Salt-Baptism of a Captain. \emph{Cracks}: relic debt-waiver, survivor testimony, salt turning against them.
\item \textbf{The Vermilion Hand (Ghost Auditor)} — \emph{Phases}: Debt Called $\rightarrow$ Witness Summoned $\rightarrow$ Ledger Sealed. \emph{Cracks}: pay a forgotten wergild, speak the true name of the Ninth Heir, survive a night in the Ledger Vault.
\end{itemize}

\begin{tcolorbox}[colback=black!3,colframe=black!40!white,title={Patronage \& Power}]
Among the Oshiira, power flows through grain, water, and the written word. The Widow holds ultimate authority, but the Vermilion Corps audits without mercy, the Canal Mothers control the waterways, and the Salt Syndicate finances the caravans. A tribe's standing depends on its harvest records; a merchant's fortune on the accuracy of his scales. But with the Widow's health failing, loyalty is no longer measured in bushels — it is measured in heirs.

\textbf{For the GM:}  
Power in Oshiiran society is bureaucratic and paper-based, but the succession crisis adds a layer of paranoia and betrayal. Patronage often takes the form of sealed receipts, priority tokens, or debt waivers — each recorded in the Ledger and subject to audit. To emphasize the Widow's decline:
\begin{itemize}
\item Tie rewards to visible documents (receipts, passes, seals) that can be challenged, stolen, or voided.
\item Let rival heirs issue conflicting claims, forcing players to choose whose ledger matters more.
\item Use the Succession Clocks as a visible campaign front. Each session, tick one heir's clock based on party actions (or inaction).
\item The Widow's Paranoia Clock can be lowered by proving loyalty (e.g., exposing a spy, recovering a stolen page) or raised by rumours, failures, or the appearance of the Ninth Heir.
\end{itemize}
In Oshiiran society, your receipt is your bond, and your grain is your greatest treasure. But when the Widow dies, the Ledger may be rewritten. And the Ledger always wins.
\end{tcolorbox}

\clearpage
\subsection*{The Four Heirs: A Game of Succession}
\label{sec:oshiira-heirs}

\begin{tcolorbox}[colback=black!3,colframe=red!60!black,title={The Widow's Children — Four Claims, One Throne}]
The Widow has four living children. Each believes they are the rightful heir. Each is wrong — or right — depending on which page of the Ledger you read. The Vermilion Corps has a secret fifth ledger that names the true successor. That page was torn out forty years ago. The Widow tore it herself. She has not told anyone what it said.

\vspace{2mm}
\begin{longtable}{|p{0.18\textwidth}|p{0.28\textwidth}|p{0.28\textwidth}|p{0.16\textwidth}|}
\hline
\textbf{Heir} & \textbf{Faction} & \textbf{Claim} & \textbf{Clock} \\
\hline
\textbf{Prince Kwame} & Old Blood (Traditionalist) & Firstborn. The law is clear. & \texttt{[6]} \\
\hline
\textbf{Princess Adwoa} & Vermilion Corps (Reformer) & The Ledger chose her. She has the stamp. & \texttt{[6]} \\
\hline
\textbf{Prince Sefu} & Narrow Captains (Militarist) & The canals obey him. The navy follows. & \texttt{[6]} \\
\hline
\textbf{Princess Nia} & Salt Syndicate (Merchant) & She holds the debts of the other three. & \texttt{[6]} \\
\hline
\end{longtable}

\vspace{2mm}
\textbf{The Ninth Heir (Secret):} A child no one acknowledges, raised in the Granary Tower's lowest vault. She was taught to read ledgers before she could walk. She is the Widow's only hope — and the Vermilion Corps' only fear. Her clock is hidden \texttt{[8]}, and it ticks only when the party learns something the Corps wants buried.
\end{tcolorbox}

\begin{tcolorbox}[colback=red!5!white,colframe=red!60!black,title={Mini-Game: The Succession Gambit}]

\textbf{Setup:} The party arrives in Mbaro during the Month of Weights, when all ledgers are audited and the canals run low. The Widow has announced she will name her successor before the month ends — or die trying. The Vermilion Corps has sealed the city gates. No one leaves until the Ledger speaks.

\vspace{2mm}
\textbf{The Objective:} Secure the succession for one heir — or broker a truce, or steal the throne for the Ninth Heir, or burn the whole system down. The party's choice determines Oshiira's fate for a generation.

\vspace{2mm}
\textbf{The Mechanics:}
\begin{itemize}
\item Each heir has a \textbf{Claim Clock [6]}. The party can tick a clock by:
  \begin{itemize}
    \item Delivering a rival's secret to that heir (+1)
    \item Assassinating a rival's key ally (+2, but marks \textbf{Heat})
    \item Forging a ledger page that discredits a rival (+1, requires a \textit{Forger's Kit})
    \item Securing a foreign alliance (Ecktoria, Kahfagia, or the Covenant) (+2)
  \end{itemize}
\item The party can \textbf{tick down} a clock by:
  \begin{itemize}
    \item Exposing the heir's own secret (–1)
    \item Convincing their key ally to defect (–1, requires a successful social test)
    \item Intercepting their grain shipment (–2, requires a raid)
  \end{itemize}
\item When any heir's clock reaches 6, they make a public bid for the throne. The Widow hears their claim. Then she asks one question: \textit{``Where is the ninth page?''}
  \begin{itemize}
    \item If the heir answers correctly, they become regent. The party gains that heir's \textbf{Boon}.
    \item If they answer incorrectly, the Widow coughs once. The Vermilion Corps arrests everyone in the room. The party has one scene to escape.
  \end{itemize}
\end{itemize}

\vspace{2mm}
\textbf{Heir-Specific Secrets (What the Vermilion Corps Knows):}

\begin{longtable}{|p{0.2\textwidth}|p{0.5\textwidth}|p{0.2\textwidth}|}
\hline
\textbf{Heir} & \textbf{Secret} & \textbf{Location of Evidence} \\
\hline
Kwame & Not the firstborn. His mother was a concubine. The Widow adopted him. & Granary Tower, Vault 9, sealed with black wax. \\
\hline
Adwoa & She ordered the assassination of a rival auditor — and pinned it on a tribal chief. & Vermilion Corps archive, hidden behind a false wall. \\
\hline
Sefu & He sold naval supply routes to Kahfagian pirates to fund his campaign. & Narrow Captains' Mess, inside a hollow cannon named ``Regret.'' \\
\hline
Nia & She owns the debt on the Granary Tower itself. If she calls it in, the tower collapses. & Salt Syndicate Counting House, the \textit{Black Ledger}. \\
\hline
\end{longtable}

\vspace{2mm}
\textbf{The Ninth Heir's Path:}
\begin{itemize}
\item To find her, the party must:
  \begin{enumerate}
    \item Discover that a page was torn from the Ledger (any archive research, DV 4)
    \item Learn that the page named a child born in secret (canal mothers' gossip, DV 3)
    \item Locate the hidden vault beneath the Granary Tower (requires a key from the Widow's Scribe)
  \end{enumerate}
\item The Ninth Heir is a child of nine years. She has never seen the sun. She knows the Ledger by heart.
\item If the party helps her claim the throne, she will rule as regent for nine years — then abdicate to the party's chosen ally. Her boon is a \textbf{Blank Page}: one debt erased from the Ledger, permanently.
\end{itemize}

\vspace{2mm}
\textbf{Possible Outcomes:}

\begin{longtable}{|p{0.25\textwidth}|p{0.55\textwidth}|}
\hline
\textbf{Outcome} & \textbf{Consequences} \\
\hline
Kwame wins & Traditionalist crackdown. Tribes lose autonomy. +2 \textbf{Heat} in the provinces. \\
\hline
Adwoa wins & Vermilion Corps takes full control. Audits become weekly. Every PC gains 1 \textbf{Obligation} to the Corps. \\
\hline
Sefu wins & War with Kahfagia within the year. The party is conscripted as deniable assets. \\
\hline
Nia wins & The Salt Syndicate privatizes the grain trade. Prices double. The party gains a \textbf{Syndicate Contact} but loses \textbf{Popular Favour}. \\
\hline
Ninth Heir wins & A regency council rules. The Ledger is rewritten. The party gains a \textbf{Blank Page} and the enmity of all four heirs. \\
\hline
No heir (civil war) & The canals run red. Mbaro burns. The party escapes with one artifact from the Granary Tower — and a blood debt to whoever saved them. \\
\hline
\end{longtable}

\end{tcolorbox}

\begin{tcolorbox}[colback=blue!5!white,colframe=blue!60!black,title={Adventure Seed: The Ninth Page}]

\textbf{The Hook:} The Widow's Scribe summons the party to the Granary Tower at midnight. He is dying. His shadow has already left. ``The page was not destroyed,'' he whispers. ``It was hidden. In a place the Vermilion Corps cannot touch. Find it before they find you.''

\vspace{2mm}
\textbf{The Clues:}
\begin{enumerate}
\item The page is wrapped in oilcloth and sealed with the Widow's own breath-bond.
\item It is hidden in a canal boat that sank nine years ago — the boat of a Canal Mother who was executed for treason.
\item The boat's wreck lies in the \textbf{Mothers' Stepwell}, the ninth step down, where the water is black and the dead do not float.
\end{enumerate}

\vspace{2mm}
\textbf{The Twist:} The page does not name an heir. It names a \textit{condition}: \textit{``The throne passes to whichever of my children can count the grains of sand in the hourglass of empire.''} The hourglass is in the Widow's bedchamber. The sand is not sand. It is ground bone. The bone is from the first emperor.

\vspace{2mm}
\textbf{The Choice:} The party can:
\begin{itemize}
\item Bring the page to the Widow (she will read it aloud at the succession council — and die immediately after, as the breath-bond expires).
\item Destroy the page and let the heirs fight (civil war, but the party can broker a truce).
\item Rewrite the page (requires a \textit{Forger's Kit} and a drop of blood from each heir — and a truth the Ledger does not yet know).
\end{itemize}

\vspace{2mm}
\textbf{Reward:} If the party successfully resolves the succession without bloodshed (or with very \textit{careful} bloodshed), each PC gains a \textbf{Claim Token}: a single vote in the next succession council. The votes can be traded, sold, or used to blackmail future rulers.

\end{tcolorbox}

% Index entries
\index{Oshiira|see{Ledger Empire}}
\index{Theme generator}
\index{Ledger Empire generator}
\index{Places!Oshiira}
\index{People!Oshiira}
\index{Complications!Oshiira}
\index{Rewards!Oshiira}
\index{Widow of Oshiira}
\index{Vermilion Corps}
\index{Canal Mother}
\index{Narrow Captain}
\index{Salt Syndicate}
\index{Grain Ledger}
\index{Shoreless Bay}
\index{Granary Tower}
\index{Mudflat weigh-station}
\index{Kahfagian cold war}
\index{Stepwell of Mothers}
\index{Vermilion Hand}
\index{Ninth Heir}
\index{Succession clocks}
\index{Local Patrons!Oshiira}

\newpage

% -------------------- ASHAAN --------------------
\chapter{Ashaan}
\emph{The Veiled Throne.}

\noindent
\textbf{Content Warning:} Ashaan depicts systemic cruelty, slavery, state violence, and child sacrifice (secret cults). Read the \textit{Safety \& Sensibility} section before playing.

\noindent
Kofi mutters: \textit{“Once, the Ashaani Empire stretched from the Red Desert to the eastern isles. Now it’s a rump state propped up by spirit‑binding and fear. The Veiled Sisters call it order. I call it a slow suicide.”}

\noindent
Ashaan is a nation that died long ago. The Veiled Sisters keep its corpse walking through spirit‑binding, terror, and the endless consumption of their own people. The capital, Galanina, is a city of black marble and brass cages. The Red Desert grows larger every year. The Breath‑less resistance whispers in root cellars, learning to read ledgers without leaving marks. Do not expect easy answers here – only hard choices and heavier prices.

\vspace{0.5em}
\noindent\textit{“Never name the dead. The Sisters bind spirits to names. Speak a dead name aloud, and you may find yourself bound alongside them.”} \\
— Tamar, Breath‑less cell leader

\vspace{1em}
\begin{almanac}
\textbf{What do people eat for breakfast?} Black bread and thin soup – the bread from grain grown on the edge of the Red Desert, the soup flavoured with salt and memory. The memory is of green fields. The fields are gone.

\textbf{What does a child learn before they can walk?} To veil their face. Cover the mouth when speaking, the eyes when listening. A child who cannot veil will be taken by the Sisters.

\textbf{What is the first thing you notice about a stranger?} Their veil. Red for mourning, white for marriage, black for the Sisters. Grey means Breath‑less. Torn means fugitive. No veil means a fool – and fools are collected by the Black Hand.

\textbf{What is the most common crime?} Spirit‑trading – the real economy. The Sisters bind spirits to brass cages and sell them. Foreign buyers never know what they are buying.

\textbf{What does no one talk about but everyone knows?} The Red Desert was a garden before the Sisters came. They bound the spirits of the soil. The soil died. The sand still whispers – the garden’s ghosts.

\textbf{Where do people go when they die?} To brass cages. The dead are bound into spirit‑lanterns and hung from the walls of Galanina. They cannot speak, cannot scream. They can only watch.

\textbf{How do you apologize for a serious wrong?} You offer a name – a dead name, a forgotten name. The Sisters bind it to a new cage. The cage is placed in the Hall of Silence. The Hall is always full.

\textbf{What is the most important festival?} The Night of the Unbound (winter solstice). The spirits strain against their cages. The Sisters renew their bindings. The city is silent for three days. The silence is fear.

\textbf{What is the secret that could destroy a powerful person?} The High Sister is not a sister. She is a spirit – the first one bound, nine hundred years ago. She bound herself to the throne. She cannot leave, cannot die, cannot stop binding.
\end{almanac}

\newpage
\section{Ashaan −−− ``The Veiled Throne''}
łabel{chap:ashaan}

\textbf{Key Demographic: } Human
\textbf{Captial: } Ashaan (The City)

\begin{tcolorbox}[colback=red!5!white,colframe=red!60!black,title={Content Warning}]
Ashaan is a setting of systemic cruelty: chattel slavery, state−sanctioned violence, child sacrifice (secret cults), and the erasure of identity. The Veiled Sisters rule through terror, the Black Hand erases dissent, and the Breath−less resistance struggles in shadows. These themes are presented as horror and injustice, never glamourised. Before playing, use Lines & Veils to ensure everyone's comfort. If you need to adjust or omit elements, the GM has final authority. Your safety matters more than canon.
\end{tcolorbox}

\subsection*{Quick Reference}
\textbf{Tagline:} \textit{The Veil is not a choice. It is survival. Speak in thirds — two options and a hidden third. The ninth word is always a knife.}

\begin{quote}
  \textbf{Veiled Sister (Elite):}  
  ``The Veil is not a choice. It is survival. You will show your eyes and nothing more. Speak in thirds — offer two options and a hidden third. The ninth word is always a knife.''
\end{quote}

\begin{quote}
  \textbf{Breath−less Agent (Commoner):}  
  ``They say Ikasha sleeps. I say she is patient. The Maeri have ruled for generations; they believe the Veil is eternal. But the Veil is silk. Silk frays. And I am learning to bite through — one whisper at a time.''
\end{quote}

\noindent\textbf{Genre / Mood:} Grimdark spirit−horror, bureaucratic terror, whispered resistance, creeping divine punishment.

\noindent\textbf{Starting Location:} A root cellar beneath a spice merchant's shop, lit by a smuggled candle. A Breath−less agent presses a folded paper into your hand: ``My brother argued with an Esoiti about a misweighed sack of cardamom. His name was erased from the rolls by morning. Help me find him — or help me remember his name.''

\vspace{2mm}
\begin{tcolorbox}[colback=black!3,colframe=blue!60!black,title={Theme & Atmosphere}]
Once, the Ashaani Empire stretched across the eastern Amaranthine rim — a river−fed dominion of sorcerer−kings, the 999 Immortals, and the great twin−river metropolis of Khasurah. Three centuries ago, an invasion of Oshiira shattered the empire. The Immortals were annihilated; plague, blockade, and northern mercenaries finished the ruin. Today, what remains is a rump city−state called Ashaan (or As'Ashaan), a fortified river−island north of the Khesai delta. The old lands — the \textbf{Green Ribbon} floodplains, the \textbf{Great Desert} where glass hums with ancient curses — are a patchwork of loyal towns, restless hill peoples, and the ruins of swallowed cities like Khasurah, which still breathes under dunes and salt crusts.

But the fall of Ashaan was not merely conquest. It was \emph{punishment}. The priests whisper that the 999 Immortals committed a sin so profound that the gods turned their faces away — and the sand remembers what the living have forgotten. The Veil, the magical membrane that separates the living from the dead, is not a barrier. It is a \emph{wound} that never healed. And the wound is spreading.

\textbf{Neighbors and Rivals:}
\begin{itemize}
  ℹtem \textit{West – Oshiira:} The grain empire that broke Ashaan's back. The two nations maintain a cold peace, punctuated by Oshiiran grain barges that pay ruinous tariffs to pass the Brass Gate. Ashaani spies in Mbaro watch for any sign that the Widow might seek to finish what her ancestors started.
  ℹtem \textit{South – Taharka:} The highland kingdom of the Monsoon Crown. Taharkan passes are a graveyard for Ashaani armies; the Veiled Sisters have never broken those mountains. Now Ashaani agents slip through disguised as spice merchants, seeking to weaken Taharka from within.
  ℹtem \textit{East – Galanina (Sidhi ancestral homeland) and the Pereshi Highlands:} Galanina, the old eastern capital, lies under Ashaani control — but its Catacombs of the Unburied hold the bound spirits of Sidhi slaves, and the Breath−less dream of freeing them. Beyond, the Pereshi fire temples guard the secret of the Forgotten Prince, an Ashaani conqueror whose true name was severed. The Maeri want that name; the Pereshi will never give it.
\end{itemize}

The state is ruled by three masked beings known as the \textbf{Maeri} — the Veiled Sister (coins and winter), the Helmed Sister (steel and sickness), and the Masked Sister (red and black). No one knows what they are: demigods, archfey, lich−queens, or the last masks of the First Fey. They speak through the \textbf{Esoti}, a collegium of wizards who maintain the wards, govern the bureaucracy, and interpret the Maeri's will. In secret, certain Esoti circles conduct child sacrifice to Khemesh, a sea−demon cult imported in the empire's maritime days.

The \textbf{Black Hand} is the Maeri's instrument: a fanatical network of bards and warlocks divided into three Aspects — Veiled (accountants, heralds, ruin by paper), Helmed (assassins and enforcers, ruin by blade), and Masked (courtesans, impresarios, ruin by desire). The Hand removes over−mighty nobles, disappears rival magi, and unpersons threats. Within it, ambition is a short road to an unmarked grave. Unlike the Sisters, the Hand is \emph{human} — broken, desperate, or fanatical, but still capable of fear, doubt, and the occasional flicker of mercy.

Slavery is legal, widely practiced, and increasingly focused on wandered Sidhi lineages broken by the empire's fall. Enslaved Sidhi bear blue marks on the forearm; illegal to remove without a manumission writ, but earth−wise midwives know ways to fade it. In the catacombs beneath Galanina — the old eastern capital — the dead are mapped like a city, and the spirits of the unburied are bound into brass cages.

Yet the Veil is cracking. In Galanina, devotion to the Angel Mykkiel (a solar of unbending Law and Justice) spreads. His zealots are incorruptible — an acid slowly eating the old order. And beneath the capital, the goddess Ikasha (She Who Sleeps) stirs in her sealed vault. Her followers, the \textbf{Breath−less}, meet in root cellars, learn to read records without leaving marks, and whisper that a name restored can weaken the Veil.

\vspace{2mm}
\textbf{What you notice first:} The smell of incense and rust. The way every doorway has a veil of brass chains. The silence when a Veiled Sister passes — even the spirits hold their breath. The brass tags on the dead — one for each Sister. And the green glass that glints in every shadow, fragments of a desert that was once a garden.

\textbf{The one rule that kills newcomers:} Never name the dead. The Sisters bind spirits to names. Speak a dead name aloud, and you may find yourself bound alongside them. Also, never address a coin−veiled speaker at court unless you are ready to accept their price.
\end{tcolorbox}

\subsection*{Regional Mood: The Weight of Damnation}

Power here is measured in spirits bound, names erased, and the silence of witnesses. The Maeri rule through terror; the Esoti hold the records; the Black Hand hunts without uniform. The Breath−less whisper in cellars, learning to read records without leaving marks. Ashaan is not a place of open rebellion. It is a place of whispered warnings, vanished neighbors, and scrolls that rewrite themselves.

But beneath all of this runs the \emph{creeping}. The Red Desert expands by inches each year, its green glass humming with a forgotten curse. The spirits bound in brass cages grow restless; their whispers are prayers to gods who no longer answer. The Veil, once a clean wound, has begun to fester. No one remembers what sin caused it. Perhaps no one ever knew. But everyone feels the weight — a pressure behind the eyes, a dryness in the throat, a certainty that something is \emph{coming}.

\subsection*{Prominent NPCs of Ashaan}
łabel{sec:ashaan−npcs}

\begin{longtable}{|p{0.2\textwidth}|p{0.25\textwidth}|p{0.3\textwidth}|p{0.15\textwidth}|}
ℎline
\textbf{Name} & \textbf{Role/Title} & \textbf{Motivation} & \textbf{Rumoured Quirk} \\
ℎline
The Veiled Sister & Maeri of coins and winter & Maintain the Veil and the tribute system & She has never removed her veil of gold coins. The coins are from dead empires. She does not blink. Her voice is the same as the wind over a mass grave. \\
ℎline
The Helmed Sister & Maeri of steel and sickness & Enforce order through terror & Her bronze helm never opens. Those who see inside never speak again. Her shadow has teeth. She does not eat. \\
ℎline
The Masked Sister & Maeri of red and black & Weave desire into control & She knows every secret whispered in Ashaan's brothels and art houses. Her smile is a contract. She has no reflection. \\
ℎline
Magistrate Eso & Esoiti high magistrate & Control every record, every court, every gate & He rewrites his own memory each morning. His left hand is a stamp. His right hand is a key that fits no lock. \\
ℎline
The Black Hand (Aspect Commander) & Faceless hunter & Erase anyone who whispers Ikasha's name & No one knows their faces. They could be your neighbor, your clerk, your spouse. They are never seen eating. \\
ℎline
Tamar & Breath−less cell leader & Tear down the Veil from within & She has three names; each is for a different jurisdiction. Her shadow moves only when she is still. \\
ℎline
Jordan & Spirit−broken slave & A child taken for ``reevaluation,'' returned hollow−eyed & He no longer speaks. He only weeps when he hears a brass bell. His shadow is a brass cage. \\
ℎline
The Whisperer & Forbidden mage in hiding & Learn to read records without leaving marks & He writes on water. The words vanish, but the Sisters cannot read them. His reflection is missing. \\
ℎline
\end{longtable}

\subsection*{Names of Ashaan}
łabel{sec:ashaan−names}

Inspired by the eastern river−core and the long shadow of old empires: sharp consonants, vowel−rich endings, with echoes of ancient Egyptian and Seljuk naming conventions. Surnames often denote a lost lineage, a trade, or a place swallowed by the sand.

\begin{longtable}{|p{0.25\textwidth}|p{0.25\textwidth}|p{0.2\textwidth}|p{0.2\textwidth}|}
ℎline
\textbf{First Names (Male)} & \textbf{First Names (Female)} & \textbf{Surnames} & \textbf{Epithets} \\
ℎline
Jordan & Morai & al−Galanina & \textit{the Veiled} \\
ℎline
Tarek & Iman & Esoiti−broken & \textit{the Breath−less} \\
ℎline
Jasin & Nadira & Black−Hand−marked & \textit{the Unnamed} \\
ℎline
Zahir & Tamar & Spirit−bound & \textit{the Ninth Witness} \\
ℎline
Rashid & Layla & Ledger−erased & \textit{the Forgotten} \\
ℎline
Dawud & Amina & Cage−born & \textit{the Hollow} \\
ℎline
Seti & Neferet & Khasurah−lost & \textit{the Glass−Voiced} \\
ℎline
\end{longtable}

\paragraph*{(Gate/Sanctum/Pan)} Salt pan with a broken spirit−cage; black marble temple with steps stained red; the Brass Gate — a fortified strait with spirit−buoys.

\section*{\spadecard{A} Places}
łabel{sec:ashaan−places}
\begin{enumerate}
\setcounter{enumi}{1}
ℹtem \textbf{Salt Pan with Broken Spirit−Cage} −−− A white flat where the air hums. \texttt{[DEATH]} \texttt{[SPIRIT]}
  ℎfill \textit{The cage is empty, but it still rattles at night. Something got out. The salt tastes of iron.}
ℹtem \textbf{Black Marble Temple (Maeri)} −−− Steps stained red; the blood never washes away. \texttt{[DEATH]} \texttt{[HORROR]}
  ℎfill \textit{The temple has no windows. The Maeri do not need light. The steps are warm, as if a great beast sleeps beneath.}
ℹtem \textbf{The Brass Gate} −−− A fortified strait with spirit−buoys that scream when crossed. \texttt{[TRAVEL]} \texttt{[SPIRIT]}
  ℎfill \textit{Each buoy contains a bound spirit. They sing the names of those who pass. The ninth buoy is silent.}
ℹtem \textbf{Veiled Sister's Sanctum} −−− Entrance hidden by a mirage that shows your greatest fear. \texttt{[MAGIC]} \texttt{[FEAR]}
  ℎfill \textit{The mirage is not an illusion. It is a memory the Sister has stolen. The memory is always from a child.}
ℹtem \textbf{Galanina Slave Market} −−− A circular pit with cages of brass and bone. \texttt{[DEATH]} \texttt{[WEIGHT]}
  ℎfill \textit{The auctioneer wears a mask. The mask is made from the face of a debt−seller. The face still moves.}
ℹtem \textbf{Esoiti Record Hall} −−− A long chamber with shelves of silent scrolls. \texttt{[LAW]} \texttt{[WEIGHT]}
  ℎfill \textit{The scrolls write themselves. The ink is the blood of erased names. The ninth shelf is empty — reserved for the name that must never be recorded.}
ℹtem \textbf{Black Hand Garrison} −−− A nondescript building with no windows and a single door. \texttt{[DANGER]}
  ℎfill \textit{The door has no handle. It opens inward only when a name is whispered. The whisper is never forgotten.}
ℹtem \textbf{Breath−less Cellar} −−− A root cellar with a single candle and a hole in the floor. \texttt{[SECRET]}
  ℎfill \textit{The hole leads to a tunnel that leads to a tomb. The tomb is a safe house. The candle never burns lower.}
ℹtem \textbf{Red Desert's Edge} −−− Where the fertile land turns to dust and glass. \texttt{[WEATHER]} \texttt{[CURSE]}
  ℎfill \textit{The glass is green. It hums with the memory of the prince's curse. The sound is the same as a brass cage rattling.}
ℹtem[J] \textbf{The Ninth Name Shrine} −−− A ruin where a forbidden name is carved into stone. \texttt{[DEATH]} \texttt{[SECRET]}
  ℎfill \textit{The stone weeps water. The water tastes of iron and forgotten vows. The name is worn smooth — no one can read it, but everyone fears it.}
ℹtem[Q] \textbf{Spirit−Binding Foundry} −−− A forge where brass cages are made. \texttt{[CRAFT]} \texttt{[SPIRIT]}
  ℎfill \textit{The forge is heated by a bound fire−spirit. It screams when the bellows pump. The screams are prayers in a dead language.}
ℹtem[K] \textbf{The Hollow Vessel Nursery} −−− An orphanage where children are prepared for possession. \texttt{[HORROR]} \texttt{[DEATH]}
  ℎfill \textit{The children have no shadows. Their shadows have been sold. The ninth cot is always empty — for the child who will be the Hollow's vessel.}
ℹtem[A] \textbf{Ikasha's Sleeping Chamber} −−− A sealed vault beneath Galanina. \texttt{[MYSTERY]} \texttt{[PATRON]}
  ℎfill \textit{Inside, a stone sarcophagus. The stone is warm. She Who Sleeps dreams of freedom. Her dreams leak through the cracks — green glass grows where they touch stone.}
\end{enumerate}

\paragraph*{(Hunter/Agent/Magician)} Slave−hunter (Black Hand) with a branding iron; Breath−less agent (freed slave, now spy); an Ashaani magician who offers dangerous knowledge.

\section*{ℎeartcard{A} People & Factions}
łabel{sec:ashaan−people}
\begin{enumerate}
\setcounter{enumi}{1}
ℹtem \textbf{Terrified Baker} −−− Kneads dough with shaking hands. He knows which Black Hand guards take bribes and which take only silence. He will sell you bread, but he will not meet your eyes. If the Hand asks, he did not see you. He will weep afterward. \texttt{[SURVIVAL]} \texttt{[FEAR]}
  ℎfill \textit{His son was taken for ``reeducation'' last year. He does not know if the boy is alive. He is afraid to ask.}
ℹtem \textbf{Ledger Clerk (Esoiti)} −−− Ink−stained fingers, hollow eyes. Copies erasure orders all day. He knows whose names are on next week's list. He will trade that knowledge for a single assurance that his own name isn't on it — but he cannot be sure you can deliver. He is not brave. He is merely desperate. \texttt{[LAW]} \texttt{[WEIGHT]}
  ℎfill \textit{He has a wife and a daughter. The daughter's name is written on a scrap of paper hidden in his shoe. That is his only insurance.}
ℹtem \textbf{Black Hand Recruit (Fresh)} −−− Still has a family. Still flinches when the brand touches his arm. He joined because the mines are a death sentence. He does not know what he will be asked to do next week. He is terrified. He will let you pass if you promise to tell his mother he didn't suffer. \texttt{[COMBAT]} \texttt{[TRAGEDY]}
  ℎfill \textit{He has never killed anyone. He does not know if he can.}
ℹtem \textbf{Mother at the Well} −−− Gave her youngest child's name to the Sisters. Not because she is evil. Because the alternative was her whole family being taken. She draws water at dawn when acolytes aren't watching. She will not speak of it. But if you ask about the \textit{ninth child}, she will weep. \texttt{[DEATH]} \texttt{[SORROW]}
  ℎfill \textit{The child's name was Laila. She hums her daughter's lullaby to the rope. The well remembers.}
ℹtem \textbf{Black Hand Veteran} −−− Has served twenty years. Her hands do not shake. She has seen the Sisters' work up close. She knows there is no escape, only delay. She will not help you escape — but she might sell you a day's head start. The price is a story. The story must be true. \texttt{[DANGER]} \texttt{[WORLDWEARY]}
  ℎfill \textit{She has forgotten the taste of honey. She would like to remember.}
ℹtem \textbf{Esoti Junior Magistrate} −−− Inherited her post from an uncle who disappeared. She knows the laws are lies. She enforces them anyway because the alternative is the Binding Foundry. She drinks too much and weeps in the archive. She has a key to a vault. She will lend it — for a story about what freedom tastes like. \texttt{[LAW]} \texttt{[SECRET]}
  ℎfill \textit{The key is a brass finger bone. It fits a lock in the Ninth Ledger's vault. She does not know what it opens.}
ℹtem \textbf{Slave (Household)} −−− Not a rebel. Not a hero. Just a woman who has learned to be invisible. She knows the secret passages behind the kitchen. She will trade directions for a single piece of news from outside — any news, even bad. She wants to know if her village still exists. \texttt{[STEALTH]} \texttt{[SURVIVAL]}
  ℎfill \textit{She wears a brass tag with a number. She has memorised the number of every slave in the household. She is waiting for the day the numbers no longer matter.}
ℹtem \textbf{Black Hand Officer (Compromised)} −−− He takes bribes. He looks the other way. He has a daughter in the Sisters' care as a ``ward.'' He knows she is a hostage. He will let you pass if you carry a letter to her — a letter that says nothing, but is written on paper that smells of home. He is not your ally. He is a man drowning. \texttt{[AUTHORITY]} \texttt{[TRAP]}
  ℎfill \textit{The paper is his daughter's old schoolwork. The ink is his own blood. He does not know if she can still read.}
ℹtem \textbf{Veiled Sister (Acolyte, not a full Sister)} −−− She was given to the temple at nine. She does not remember her mother. She believes the Veil protects the empire from something worse. She has never seen a slave beaten — they hide that from novices. She is not evil. She is \emph{cultivated}. A crack in her certainty may take years to open. You do not have years. \texttt{[MYSTERY]} \texttt{[DANGER]}
  ℎfill \textit{She wears a veil of hammered silver. Underneath, her face is smooth — no scars, no wrinkles. She has not aged since she took her vows.}
ℹtem[J] \textbf{The Whisperer (Breath−less cell leader)} −−− She has three names, each for a different jurisdiction. She moves through the catacombs like smoke. She will trust you only after you have broken a small law in her presence — not to prove your courage, but to prove you have something to lose. \texttt{[SECRET]} \texttt{[RESISTANCE]}
  ℎfill \textit{She writes messages on water. The words vanish, but the Sisters cannot read them. Her reflection is missing. She sold it for a name.}
ℹtem[Q] \textbf{The Black Hand's Shadow} −−− A Hand agent who has stopped sleeping. He sees the Sisters' work in his dreams. He knows he cannot leave. He will not betray his oath — but he might fail to see you, if the price is a single night of dreamless rest. The price is a memory. He will take the happiest one. \texttt{[DEATH]} \texttt{[HORROR]}
  ℎfill \textit{He carries a brass bell that rings once every hour. It is the only thing that keeps him awake. He has not heard it ring in three days.}
ℹtem[K] \textbf{The Esoti Censor (Desperate)} −−− A high−ranking magistrate who has been ordered to falsify a lineage. If she refuses, her family is erased. If she complies, she knows the truth will outlive her. She will give you the original document — for a promise that her grandson will be told, one day, why his grandmother's name was struck from the rolls. \texttt{[LAW]} \texttt{[TRAGEDY]}
  ℎfill \textit{The document is written on human skin. The skin was her own, cut from her thigh when she took office. She has no feeling there anymore.}
ℹtem[A] \textbf{The Veiled Sister (Full)} −−− Malignant divine justice made flesh. She does not threaten. She does not bargain. She states what will happen. She is always correct. She has never been seen eating, sleeping, or blinking. She is not a woman. She is a function of the Veil. There is no complexity here. There is only the weight of a god's attention. \textit{(No debt. No price. Only survival.)} \texttt{[PATRON]} \texttt{[APOCALYPSE]}
  ℎfill \textit{Her veil is made of gold coins from a dozen dead empires. They do not clink. They weep when she speaks.}
\end{enumerate}

\paragraph*{(Dust/Gate/Spirit)} Spirit−dust storm (Navigation DV +1, mark Fatigue); Brass Gate closed for ``inspection'' — delay 2 segments; a bound spirit breaks free and attacks.

\section*{\clubcard{A} Complications/Threats}
łabel{sec:ashaan−complications}
\begin{enumerate}
\setcounter{enumi}{1}
ℹtem \textbf{Spirit−Dust Storm} −−− Red sand from the desert carries whispering spirits. \texttt{[WEATHER]} \texttt{[DEATH]}
  ℎfill \textit{Navigation DV +1. Mark Fatigue. The voices ask for your name. Do not answer. They already know it.}
ℹtem \textbf{Brass Gate Closed} −−− The strait is sealed for ``inspection.'' Delay 2 segments. \texttt{[LAW]} \texttt{[DELAY]}
  ℎfill \textit{The ``inspection'' is a hunt for a Breath−less agent. The agent is in your party. The gate knows. The gate is patient.}
ℹtem \textbf{Bound Spirit Breaks Free} −−− A brass cage cracks. The spirit attacks. \texttt{[DEATH]} \texttt{[COMBAT]}
  ℎfill \textit{The spirit was once a child. It weeps as it fights. Killing it feels like murder. Freeing it feels like surrender.}
ℹtem \textbf{The Veil Fails} −−− For a moment, the Maeri's magic falters. Spirit−visions haunt all. \texttt{[MAGIC]} \texttt{[FEAR]}
  ℎfill \textit{You see the face of someone you wronged. They are asking for an apology. You cannot remember what you did. That is the punishment.}
ℹtem \textbf{Black Hand Sweep} −−− A sudden search of the district. Everyone's papers are checked. \texttt{[LAW]} \texttt{[DANGER]}
  ℎfill \textit{Your papers are in order. The Black Hand does not care. They take you anyway. The order was the trap.}
ℹtem \textbf{Esoiti Audit} −−− A magistrate demands to see your records. Your records are lies. \texttt{[LAW]} \texttt{[WEIGHT]}
  ℎfill \textit{The magistrate smiles. ``Everyone lies. I just need to know how much.'' His quill is already moving.}
ℹtem \textbf{Spirit−Hound Scent} −−− A pack of bound hounds is on your trail. \texttt{[DEATH]} \texttt{[TRACKING]}
  ℎfill \textit{The hounds are invisible. You can hear their breathing, close and wet. Their breath smells of brass and old prayers.}
ℹtem \textbf{Slave Raid} −−− A Veiled Sister's caravan is taking villagers. You are in the way. \texttt{[COMBAT]} \texttt{[WEIGHT]}
  ℎfill \textit{The villagers are armed. They have been expecting you. You are the bait. The Sister is not here. She is watching.}
ℹtem \textbf{Ikasha's Whisper} −−− In a quiet moment, you hear a voice: ``Free me.'' \texttt{[OMEN]} \texttt{[PATRON]}
  ℎfill \textit{The whisper comes from a crack in the floor. A single green eye looks up. The eye is glass. It blinks.}
ℹtem[J] \textbf{The Ninth Name Spoken} −−− Someone in the party speaks a dead name aloud. \texttt{[DEATH]} \texttt{[WEIGHT]}
  ℎfill \textit{The air grows cold. A spirit appears. It wants to know why you called. It does not accept silence as an answer.}
ℹtem[Q] \textbf{Sister's Test} −−− A Veiled Sister offers a deal: betray a Breath−less contact for safety. \texttt{[SOCIAL]} \texttt{[WEIGHT]}
  ℎfill \textit{The offer is genuine. The price is your soul. The Sister watches you decide. She is not in a hurry. She has all of eternity.}
ℹtem[K] \textbf{Mass Haunting} −−− The bound spirits rise as one. The Veil is tearing. \texttt{[DEATH]} \texttt{[HORROR]}
  ℎfill \textit{The spirits do not attack. They simply ask: ``Why did you let this happen?'' They do not mean the empire. They mean the children.}
ℹtem[A] \textbf{The Demon's Weight Called} −−− Varash the Unmaker (or a fragment) demands payment. \texttt{[DEATH]} \texttt{[WEIGHT]}
  ℎfill \textit{The amulet Kalrantha glows. A voice whispers: ``Speak your desire.'' The whisper is your own. You do not remember giving it permission.}
\end{enumerate}

\paragraph*{(Name/Lamp/Seal)} A slave's true name (free them or bind them further); spirit−lamp (one question to a bound spirit); brass seal (guest status in a Sister's sanctum).

\section*{⋄card{A} Rewards/Leverage}
łabel{sec:ashaan−rewards}
\begin{enumerate}
\setcounter{enumi}{1}
ℹtem \textbf{Slave's True Name} −−− Use it to free a bound person — or bind them further. \texttt{[WEIGHT]} \texttt{[POWER]}
  ℎfill \textit{The name is written on a scrap of skin. The skin was the slave's own. The letters are still warm.}
ℹtem \textbf{Spirit−Lamp} −−− Ask one question of a bound spirit. It cannot lie. \texttt{[MAGIC]} \texttt{[TRUTH]}
  ℎfill \textit{The lamp burns green. The spirit's face appears in the flame. It is weeping. The tears evaporate into glass.}
ℹtem \textbf{Brass Seal (Guest Status)} −−− Safe passage in a Sister's sanctum for one visit. \texttt{[SOCIAL]} \texttt{[PROTECTION]}
  ℎfill \textit{The seal is warm. It glows when a Sister is lying. It also glows when you are lying. The Sister does not ask which.}
ℹtem \textbf{Veil−Clip} −−− Speak without being silenced by Maeri magic. \texttt{[MAGIC]} \texttt{[PROTECTION]}
  ℎfill \textit{The clip is a small brass spider. It crawls onto your collar and whispers. It whispers the names of those who cannot speak.}
ℹtem \textbf{Erased Name Page} −−− A page from the Esoiti record with a name struck through. \texttt{[SECRET]} \texttt{[WEIGHT]}
  ℎfill \textit{The ink is still wet. The name was erased ten minutes ago. The page is warm. It has a pulse.}
ℹtem \textbf{Breath−less Cipher} −−− A codebook for reading secret messages on water. \texttt{[SECRET]} \texttt{[STEALTH]}
  ℎfill \textit{The cipher is written on a leaf. The leaf dissolves when read. The water remembers. The water does not forgive.}
ℹtem \textbf{Spirit−Hound Whistle} −−− Call off a pack of bound hounds once. \texttt{[PROTECTION]} \texttt{[DEATH]}
  ℎfill \textit{The whistle is made from a child's bone. The child was the hounds' last victim. The hounds remember her scent.}
ℹtem \textbf{Ikasha's Dream} −−− One prophetic vision of the Maeri's next move. \texttt{[MAGIC]} \texttt{[KNOWLEDGE]}
  ℎfill \textit{The dream arrives as a headache. The vision is clear. The cost is one memory. The memory is always of a door closing.}
ℹtem \textbf{Black Hand Mask} −−− Wear it to pass as a Hunter for one scene. \texttt{[STEALTH]} \texttt{[DECEPTION]}
  ℎfill \textit{The mask has no eye holes. The mask sees for you. It also remembers. It will report what it saw to the Hand.}
ℹtem[J] \textbf{Cage−Maker's Flaw} −−− A key that opens any spirit−cage. One use. \texttt{[MAGIC]} \texttt{[PROTECTION]}
  ℎfill \textit{The key is a bent nail. It was the flaw the cage−maker left. The flaw is also a prayer. The cage−maker does not know to whom.}
ℹtem[Q] \textbf{Drowned Sister's Secret} −−− The location of a hidden Breath−less archive. \texttt{[SECRET]} \texttt{[KNOWLEDGE]}
  ℎfill \textit{The archive is in a well. The well is guarded by a ghost. The ghost's name is Mercy. She does not remember why.}
ℹtem[K] \textbf{Ninth Sister's Veil} −−− Wear it to appear as a Veiled Sister for one hour. \texttt{[DECEPTION]} \texttt{[DANGER]}
  ℎfill \textit{The veil is also a noose. Do not tie it too tight. The Sister who gave it to you is already dead. She did not untie it.}
ℹtem[A] \textbf{The Unnamed Prince's Map} −−− A map to the original Kalrantha temple. \texttt{[LEGEND]} \texttt{[QUEST]}
  ℎfill \textit{The map is drawn on human skin. The skin was the prince's own. The map glows under moonlight. It shows the way — and the cost.}
\end{enumerate}

\section*{Quick use notes}
łabel{sec:ashaan−quick−use}
\begin{itemize}
ℹtem Draw until you have all four suits: Spade = place, Heart = actor, Club = pressure, Diamond = leverage. Highest rank sets the main Timer (2−−5 \(→\) 4, 6−−10 \(→\) 6, J/Q/K \(→\) 8, A \(→\) 10).
ℹtem Diamonds are codified outcomes (names, lamps, seals) that change position rather than call for a roll.
ℹtem If any A appears, echo \textbf{veil−omens}—brass chains that rattle without wind, scrolls that open to the wrong page, a whisper that sounds like your own voice, green glass growing where no glass should be.
\end{itemize}

\subsection*{Additional Features (Cultural Mechanics)}
łabel{sec:ashaan−mechanics}

\begin{itemize}
ℹtem \textbf{The Veil:} The Maeri's control is maintained through spirit−binding and terror. The Veil is a magical membrane that separates the living from the dead. When the Veil fails, spirits walk. Mechanically, the GM may spend 3 SB to trigger a \emph{Veil Failure} scene: all spirits become visible, bound spirits may break free, and social rolls against the Maeri's agents gain +1 Position. The Veil fails more often now. No one knows why.

ℹtem \textbf{The Veil's Attention (Divine Threat Timer ):} Unlike mundane factions, the Veiled Sisters do not have a favour or debt track. They have a timer: \textbf{Veil's Attention [6]}. It ticks when the party draws the Sisters' gaze — by speaking a forbidden name, by freeing a slave, by resisting a Black Hand sweep, or simply by being in the wrong place at the wrong time. When the timer ticks, the GM may impose an omen: a mirror fogs, a bell rings in still air, a shadow falls where no light is cast. When the timer reaches 6, a Veiled Sister appears. She does not offer a bargain. She pronounces a \emph{judgment}: a name to be erased, a memory to be taken, a life to be claimed. The party cannot fight her. They can only flee, or hide, or offer a substitution — a name they are willing to lose. The Sisters are not enemies. They are a natural disaster with a ledger.

ℹtem \textbf{The Breath−less (Ikasha's Cult as Active Resistance):} She Who Sleeps is not a patron of open rebellion. Her followers meet in root cellars, learn to read records without leaving marks, and whisper secrets under the cover of a cough. A character who whispers Ikasha's name before sleeping gains +1 die to resist fear for the next day, but marks 1 \textbf{Exposure} (the Black Hand may notice). Their symbol is a broken brass chain. Aiding the Breath−less may earn a \textbf{Breath−less Cipher} (see Diamonds) but advances a \textbf{Black Hand Interest} timer.

ℹtem \textbf{The 999 Immortals:} The failed invasion of Oshiira is commemorated in the Temple of the Golden Sepulcher, where the gem−eyed skulls of the 999 are interred. Each New Flood, magistrates process with gold−veiled reliquaries, praying not to repeat their hubris. The guardless helmed blades of the Immortals became trophies; many now hang at the hips of Oshiira's greatest warriors. Some say the skulls whisper at night. They whisper a name that no one will speak.

ℹtem \textbf{Veil Etiquette:} At court, a face fully veiled in coin signals ``I speak as the treasury.'' To address such a speaker is to accept their price. The ninth cup at elite tables is lifted without drinking — a silent remembrance of the 999 and a warning against hubris. Masktime (the hour before sunrise) is when contracts are sealed, lovers dismissed, and bribes counted; unmasked business at that hour is considered offensive. Those who break the ninth cup rule find that their next drink tastes of glass.

ℹtem \textbf{Local Patrons (Syncretic Names):}
  \begin{itemize}
    ℹtem \textbf{The Unbound Chain} −− Ikasha. \textit{She Who Sleeps; her dream is freedom. Her dream is leaking.}
    ℹtem \textbf{The Gilded Scale} −− Mykkiel. \textit{The Maeri's twisted version of divine law; the zealots of Galanina. Their scales weigh not coin, but names.}
    ℹtem \textbf{The Drowned Prince} −− Khemesh. \textit{Pressure from the deep; the Red Desert was once his seabed. His cult practices child sacrifice in secret. The sacrifices are always the ninth child.}
    ℹtem \textbf{The Silvered Tongue} −− Malachai. \textit{Curses sold as blessings, whispered in the slave markets. The blessing always includes a ninth clause. No one reads it.}
  \end{itemize}
\end{itemize}

\subsection*{Lore Echoes (Cross−Expansion Hooks)}
\begin{itemize}
  ℹtem \textbf{The Kalrantha Amulet (Dhahara Connection):} The Red Desert was once fertile land, destroyed by an Ashaani prince who stole the Kalrantha amulet. The Maeri want the amulet to bind the demon again — or to use it for themselves. The Breath−less want to destroy it. The prince's true name, long severed by the Pereshi, is said to be written on a scroll sealed beneath Mount Khvarena. The Maeri would pay anything to learn it. The sand remembers his crime. The sand is still hungry.
  
  ℹtem \textbf{The Sidhi of the Brass Coast (Way of Silk):} The Sidhi are coastal cousins of the Ashaani who fled slavery generations ago. They maintain lighthouse courts and freedman's seals. They remember every raid. A Sidhi captain may be a Breath−less contact — and some whisper that the green flame of their lighthouses is a fragment of Ikasha's dream, stolen from the Veiled Sisters. The flame flickers when the Veil tears.
  
  ℹtem \textbf{The Veiled Sisters and the Serpent's Turn (Isoka):} The oldest Veiled Sister legends speak of a time when name−erasure was a sacred judgment, witnessed by the Serpent (Isoka). The severed name was buried in a well of tears, and the Serpent's gaze kept the Hollow at bay. But when the Sisters began erasing names for silver and power, the Serpent turned away. Now the Serpent's gaze falls only on the Breath−less — who whisper that Ikasha's dream is to restore what the Serpent abandoned. The Sisters do not speak of this. The Hollow does not care. The Serpent's shed skin is said to be buried beneath the Brass Gate.
  
  ℹtem \textbf{The Hollow and the Ninth Name (Eastern Realms):} In Ashaan, the Hollow is the space left by an erased name. The Maeri feed the Hollow with forgotten weights. The Breath−less know that to restore a name is to weaken the Veil — and that the Hollow's hunger is the price of the Serpent's silence. A name restored can weaken a Sister's binding for a season. The Hollow is patient. The Hollow is not hungry. The Hollow is \emph{waiting}.
  
  ℹtem \textbf{Oshiira's Cold Peace:} Grain barges from Oshiira pay ruinous tariffs at the Brass Gate. In return, Ashaani spies watch every weigh−house in Mbaro, waiting for the Widow to falter. A party that smuggles grain past the Gate could earn the gratitude of starving villages — and the ire of the Veiled Sisters. The Breath−less sometimes bribe the gatekeepers with stories, not silver; stories are harder to trace. The gatekeepers have a ninth story they never tell.
  
  ℹtem \textbf{Taharka's Unconquered Passes:} Ashaani armies have shattered against the Monsoon Crown's highlands. Now the Maeri send agents disguised as spice merchants to spread poison and dissent. A Breath−less cell in the Silver Pass may hold the only copy of a treaty that could unite Taharka and Ashaan against a common foe — if anyone is brave enough to read it. The treaty is written on a scroll sealed with a knotted cord; the knot is Tulkani work, and it cannot be cut without a story. The ninth page of the treaty is blank. It is reserved for the name of the traitor.
  
  ℹtem \textbf{The Pereshi Highlands and the Forgotten Prince:} The fire temples of Pereshi guard the severed name of the Ashaani prince who tried to conquer them. The Maeri would pay anything to learn that name — for to speak it would be to unmake the Veiled Throne's own claim to authority. The Pereshi will never give it. A party caught between the two might steal the name — or destroy it forever. The Breath−less believe that destroying the name would free the Hollow's attention from Ashaan; the Maeri believe it would let them bind the Hollow itself. The prince's name is carved on a ninth coal that has never been lit. It is warm.
\end{itemize}

\subsection*{Ashaan Superstitions (Burn for +1 die once)}
\begin{itemize}
  ℹtem \textbf{Bread to the Veil:} Crumble a crust at a brass cage. Ignore the first \emph{Spirit−Dust Storm} penalty. \texttt{[SUPERSTITION]}
  ℹtem \textbf{Knot the Veil−Clip:} Tie a single knot in your veil. Cancel \emph{Black Hand Sweep} or take −−1 die on your next stealth roll (your choice). \texttt{[SUPERSTITION]}
  ℹtem \textbf{Name to the Well:} Whisper a dead Breath−less agent's name into a well. Gain +1 Effect when evading the Maeri this scene, but mark \emph{Obligation} to Ikasha. \texttt{[SUPERSTITION]}
  ℹtem \textbf{Glass to the Sand:} Drop a shard of green glass into the Red Desert at noon. The next time the Veil fails, you will see it coming one minute in advance. The glass will hum the warning. \texttt{[SUPERSTITION]}
\end{itemize}

\subsection*{Thematic SB Spend Table}
łabel{sec:ashaan−sb}

\paragraph{Minor Complications (1 SB)}
\begin{itemize}
ℹtem \textbf{Exposure:} Your actions draw unwanted attention from \textbf{Black Hand agents or Esoiti clerks}.
ℹtem \textbf{Noise:} Sounds of your actions alert nearby \textbf{spirit−hounds or slave−catchers}.
ℹtem \textbf{Trace:} Evidence of your passage marks your route for \textbf{Veiled Sisters or bound spirits}.
ℹtem \textbf{Delay:} A brief but meaningful setback costs you \textbf{time or safe passage}.
ℹtem \textbf{Supply Strain:} Mark +1 segment on a relevant \textbf{resource timer}.
\end{itemize}

\paragraph{Moderate Setbacks (2 SB)}
\begin{itemize}
ℹtem \textbf{Alarm Raised:} \textbf{Local Black Hand or Esoiti magistrate} becomes aware and begins responding.
ℹtem \textbf{Position Lost:} You lose advantageous ground/cover/stealth due to \textbf{veil failure or spirit manifestation}.
ℹtem \textbf{Foe Appears:} A \textbf{bound spirit or slave−hunter} arrives on scene.
ℹtem \textbf{Gear Trouble:} A piece of equipment becomes \textbf{Compromised/Neglected}.
ℹtem \textbf{Lock/Barrier:} A simple obstacle now requires a test to overcome.
\end{itemize}

\paragraph{Serious Trouble (3 SB)}
\begin{itemize}
ℹtem \textbf{Reinforcements:} Additional \textbf{Black Hand, Esoiti guards, or bound spirits} arrive.
ℹtem \textbf{Key Gear Breaks:} A crucial tool/weapon becomes temporarily unusable.
ℹtem \textbf{Major Twist:} The situation fundamentally changes — \textbf{Veil fails / spirit breaks free / name is spoken}.
ℹtem \textbf{Rail Tick:} Advance a relevant campaign/front timer by 1 segment.
ℹtem \textbf{Condition Applied:} Mark \textbf{Fatigue 1/Harm 1/Condition} appropriate to fiction.
\end{itemize}

\paragraph{Major Turns (4+ SB)}
\begin{itemize}
ℹtem \textbf{Trap Springs:} A prepared danger activates with full effect.
ℹtem \textbf{Authority Arrival:} \textbf{A Maeri (Veiled, Helmed, or Masked), Esoiti High Magistrate, or Ikasha's manifestation} intervenes.
ℹtem \textbf{Scene Shift:} The environment changes dramatically — \textbf{Veil tears / sandstorm rises / Brass Gate seals}.
ℹtem \textbf{Patron Omen:} Divine/arcane forces take notice — \textbf{omen appears / blessing lost / curse manifests}.
ℹtem \textbf{Narrative Pivot:} The story takes an unexpected turn that reframes objectives.
\end{itemize}

\paragraph{Region−Specific SB Options}
\begin{itemize}
ℹtem \textbf{Ashaan (The Veil):} Brass chains rattle, scrolls open to the wrong page, a whisper sounds like your own voice, a green glass shard grows from a crack in the wall.
ℹtem \textbf{Ashaan (Red Desert):} Glass shards hum, sand shifts to reveal bones, a mirage shows a memory you had forgotten, the sky turns the colour of a healing bruise.
ℹtem \textbf{Ashaan (Breath−less):} A cough from a shadow, a knock in a cellar, a name written on water that does not dissolve, a candle that burns without a wick.
\end{itemize}

\subsection*{Pursuit Ladder (Near / Mid / Far)}
\begin{itemize}
ℹtem \textbf{Advance/Withdraw (Wits + Stealth):} On a hit, shift one band. On a strong hit, also impose \emph{False Trail}.
ℹtem \textbf{Cut the Veil (Presence + Sway):} On a hit, freeze range for one exchange; if a Sister's seal is in your possession, +1 die.
ℹtem \textbf{Weather the Haunting (Resolve + Spirit):} On a hit in \emph{Mass Haunting}, you pick who is targeted; on a miss, the GM does.
\end{itemize}

\subsection*{Boss Seeds (Spirit & Veil)}
łabel{sec:ashaan−bosses}

\begin{itemize}
ℹtem \textbf{The Record−Eater (Nine Erasures)} — \emph{Phases}: Audit Thresholds \(→\) Orchestrated Erasure \(→\) Veil Coup. \emph{Cracks}: cut a lineage page; reveal swapped seal; survive a staged haunting. It appears as a scribe whose hands are brass.
ℹtem \textbf{The Unbound Sister (Brine−Veil)} — \emph{Phases}: Drownings Called Mercy \(→\) Procession at Brass Gate \(→\) Spirit−Baptism of a Magistrate. \emph{Cracks}: relic spirit−lamp, survivor testimony, the Veil turning against them. Her veil is made of chain, not silk. It clinks when she moves.
ℹtem \textbf{The Forgotten Prince's Echo (Pereshi Connection)} — \emph{Phases}: Whisper of the True Name \(→\) Fire Temple Cracks \(→\) The Prince Walks. \emph{Cracks}: burn a page of the Pereshi epic, speak a different name into the Ninth Coal, or convince the Maeri that the prince was never real. His shadow has no feet. He walks on glass.
\end{itemize}

\begin{tcolorbox}[colback=black!3,colframe=black!40!white,title={Patronage & Power}]
Among the Ashaani, power flows through fear, spirit−binding, and the erasure of names. The Maeri hold ultimate authority, but the Esoti control the records, the Black Hand enforces silence, and the Breath−less whisper in cellars. A name can be a weapon or a shield. The Mykkielites of Galanina are an acid slowly eating the old order, and Ikasha stirs in her sleep.

But the true power is the \emph{creep}. No one knows when the Veil will tear next. No one knows which name will be the ninth. No one is safe. Not the Maeri. Not the Breath−less. Not the slaves, not the masters, not the glass that grows in the dark. The punishment is spreading. The sin is forgotten. The sand remembers. The sand is patient.

\textbf{For the GM:}  
Power in Ashaani society is hidden and treacherous. Patronage often takes the form of true names, spirit−lamps, or erased pages — each carrying the weight of damnation. To emphasize this:
\begin{itemize}
ℹtem Tie rewards to secrets (names, whispers, visions) that can be betrayed, stolen, or used against the giver.
ℹtem Let rival factions (Maeri vs. Esoti vs. Breath−less vs. Mykkielites) offer conflicting paths, forcing players to choose whose weight to incur.
ℹtem Use the Brass Gate, the Record Hall, and the root cellars as arenas for social contests, where truth is a weapon and silence is a shield.
ℹtem Track \textbf{Black Hand Interest} as a campaign timer. Each time the party aids the Breath−less or escapes a Sister, tick the timer. When it fills, a Black Hand squad arrives with a warrant for their erasure — but the warrant is signed by a Sister who has been dead for a year.
ℹtem Let the Veil fail unpredictably. The GM may spend SB to trigger a tear at the worst moment. The tear is always in the shape of a ninth door.
\end{itemize}
In Ashaani society, your name is your bond — and your bond can be erased. But the erasure leaves a scar. The scar grows glass.
\end{tcolorbox}

% Index entries
∈dex{Ashaan|see{Veiled Throne}}
∈dex{Theme generator}
∈dex{Veiled Throne generator}
∈dex{Places!Ashaan}
∈dex{People!Ashaan}
∈dex{Complications!Ashaan}
∈dex{Rewards!Ashaan}
∈dex{Maeri (Three Sisters)}
∈dex{Veiled Sister}
∈dex{Helmed Sister}
∈dex{Masked Sister}
∈dex{Esoiti}
∈dex{Black Hand}
∈dex{Breath−less}
∈dex{Ikasha}
∈dex{Mykkiel}
∈dex{Red Desert}
∈dex{Brass Gate}
∈dex{Galanina}
∈dex{Khasurah (swallowed city)}
∈dex{Spirit−binding}
∈dex{999 Immortals}
∈dex{Kalrantha amulet}
∈dex{Green glass}
∈dex{Local Patrons!Ashaan}

\newpage

% -------------------- NGOMEBE --------------------
\chapter{Ngomebe}
\emph{The Stone Men.}

\noindent
Kofi says: \textit{“They call their cities ‘walking’. I call them stubborn. The stone remembers every mile, every roller, every oath. You want to trade with Ngomebe? Learn the stone‑song. And never, ever carve your name into living rock.”}

\noindent
Ngomebe are a people of deep, iron‑dark skin and older ties to the southern ranges. Their cities walk – colossal stone structures on rollers and sledges, dragged by teams of oxen, magic, and Aeler‑designed capstans. A walking city is a living argument against stillness. The stone remembers every mile, every roller, every oath sworn in the shadow of the keystone. The ancestors speak through cracks; the tokoloshe fouls the bearings at night. Do not mistake their patience for weakness.

\vspace{0.5em}
\noindent\textit{“Never carve a name into living rock. The stone remembers. It will whisper that name to every spirit of the earth, and the weight will follow your bloodline for seven generations.”} \\
— Kaelen, Stone‑Cutter

\vspace{1em}
\begin{almanac}
\textbf{What do people eat for breakfast?} Stone‑ground meal and fermented milk – the meal ground between the city’s own millstones, the milk from goats that walk alongside. The goats remember where the city has been.

\textbf{What does a child learn before they can walk?} To sing the city’s song. Eight beats to a mile, eight miles to a day, eight days to a season. The ninth beat is the city’s secret – never sung aloud.

\textbf{What is the first thing you notice about a stranger?} Their dust. White dust from the mountains, red dust from the desert. No dust means never walked with a city. The cities do not trust the dustless.

\textbf{What is the most common crime?} Stone‑theft – not a crime but a wound. A stolen stone is a forgotten year. The forgetting spreads. The city’s song stutters.

\textbf{What does no one talk about but everyone knows?} The city is not moving toward anything. It is moving away. The Ngomebe are fleeing something that lies to the south. Something that does not walk. It waits.

\textbf{Where do people go when they die?} Into the keystone. Ashes mixed into mortar and pressed into the foundation. The dead become the city’s weight. It walks heavier each year.

\textbf{How do you apologize for a serious wrong?} You carve a new roller – a single stone with no cracks, no flaws. If it rolls true, you are forgiven. If it crumbles, you are added to the mortar.

\textbf{What is the most important festival?} The Turning of the Rollers (spring equinox). The city stops for one day. The rollers are inspected, greased, blessed. The song has eight verses. The ninth is silence – for the dead.

\textbf{What is the secret that could destroy a powerful person?} The city is not a city. It is a tomb – a monument to a king who died before the first stone was laid. His name is carved into the keystone, at the city’s inaccessible heart. He is waiting to be remembered.
\end{almanac}

\subsection*{The Aelinnel of the Southern Coast}
\emph{Copper, bone‑rolls, and the dawn pact.}

\noindent
Along the southeastern coast, where the savannah meets the mangrove, a handful of fae‑touched gnome enclaves maintain the ancient pact that keeps the walking cities from sinking into soft ground. Every dawn, their bone‑spinners roll carved femurs to determine the day’s safe path. The Ngomebe tolerate them – the gnomes hold a debt older than memory, and their engineers can find a bearing’s true pitch by ear alone.

\vspace{0.5em}
\noindent\textit{“The bone‑roll never lies. If it says the city will sink, you shore the rollers. If it says the path is clear, you don’t ask why. You just walk.”} \\
— Vent‑Prior Thrain, Aeler engineer (on the gnomish bone‑roll)

\vspace{0.5em}
For full generator tables for the Aelinnel of the Inland Sea, see the \textit{Fate’s Edge Core Amaranthine Worldbook}. In the Southern Realms, they appear only as occasional allies or mysterious toll‑keepers at the edge of Ngomebe’s migration routes.

\newpage
\section{Ngomebe --- ``The Stone Men''}
\label{chap:ngomebe}

\subsection*{Quick Reference}
\textbf{Tagline:} \textit{The stone remembers. The ancestors walk with the rollers. The city that stops dies.}

\begin{quote}
  \textbf{Mason-Kin Elder (Elite):}  
  ``A city that stops moving is a city that has chosen to die. The stone remembers every mile, every roller, every oath sworn in the shadow of the keystone. We do not build cities. We raise them from the earth and teach them to walk.''
\end{quote}

\begin{quote}
  \textbf{Stone-Cutter (Commoner):}  
  ``My grandfather carved the keystone for this city. My father greased the rollers. I inherited a sledge and a song. The song is the city's heartbeat. If I stop singing, the city stops walking — and so do I.''
\end{quote}

\begin{quote}
  \textbf{Ancestral Healer (Izinyanga):}  
  ``The ancestors do not speak in words. They speak in cracks. A fissure in the keystone is not a flaw — it is a message. Learn to read the stone's wounds, or the city will read your bones.''
\end{quote}

\vspace{2mm}
\begin{tcolorbox}[colback=black!3,colframe=blue!60!black,title={Theme \& Atmosphere}]
The Ngomebe are a people of deep, iron-dark skin and older ties to the southern ranges. Their cities walk — colossal stone structures on rollers and sledges, dragged by teams of oxen, magic, and Aeler-designed capstans. This partnership is one of the great cross-cultural achievements of the age: Aeler engineers maintain the bearings; Ngomebe masons carve the keystones. A walking city is not a curiosity; it is a living argument against stillness, a monument to survival in a land where staying in one place invites famine, war, or the slow forgetting of ancestral stone.

Spirits are not distant — they live in the rock, the rollers, the cracks. The ancestors speak through keystone fissures. The \textbf{tokoloshe} fouls the bearings at night. The \textbf{impundulu} (lightning bird) shatters entire sledges with a single strike. The Ngomebe do not fight these spirits. They bargain with them: a grindstone of fresh meat, a pot of home-brewed sorghum beer, a promise whispered into a roller-pit.

\vspace{2mm}
\textbf{What you notice first:} The groan of stone on stone. The smell of warm granite, ox sweat, and sorghum beer offerings. The way the ground trembles slightly, even a mile away, as the city inches across the plateau.

\textbf{The one rule that kills newcomers:} Never carve a name into living rock. The stone remembers. It will whisper that name to every spirit of the earth, and the debt will follow your bloodline for seven generations.
\end{tcolorbox}

\subsection*{Regional Mood: The Weight of Movement}

Power here is measured in keystones, migration routes, and the loyalty of the Mason-Kin. A walking city is a polity, a temple, and a weapon all at once. The Mason-Kin Houses control the rights to carve, to move, and to settle. \textbf{Root-binders} — heretics who plant acorns in the rollers — believe the cities should stop and let the forest reclaim the stone. The \textbf{Aeler engineers} are neither masters nor servants; they are partners in a pact older than memory. But beneath them all, the \textbf{Ancestral Healers (Izinyanga)} speak for the dead, and the dead speak through the keystone's cracks. When a city stops, it is because the ancestors have been offended. When a city moves true, it is because the \textbf{Rain Queen} has blessed the path.

\textbf{Starting Location:} The edge of a moving city's path, where the stone is still warm and the dust has not yet settled. A Mason-Kin elder approaches, wiping granite dust from his face: ``A keystone has cracked. The city will stop before dawn. The ancestors are restless. Find the saboteur — a Root-binder, a tokoloshe, or a traitor among us — or push the rollers yourselves. The plateau does not forgive stillness.''

\subsection*{Prominent NPCs of Ngomebe}
\label{sec:ngomebe-npcs}

\begin{longtable}{|p{0.2\textwidth}|p{0.25\textwidth}|p{0.3\textwidth}|p{0.15\textwidth}|}
\hline
\textbf{Name} & \textbf{Role/Title} & \textbf{Motivation} & \textbf{Rumoured Quirk} \\
\hline
Elder Morvath & Mason-Kin elder of the eastern caravan & Keep the city moving south, away from the Red Desert's breath & He has not slept in forty days; the city's rhythm keeps him awake. \\
\hline
Kaelen & Stone-cutter whose city left him behind & Catch up to his family's city before it crosses the Salt Scar & He chisels a new keystone every night. None have been accepted. \\
\hline
Vent-Prior Thrain & Aeler engineer, grease-stained and level-carrying & Maintain the bearings that keep the city alive & She speaks to the rollers in a language of taps. They answer in groans. \\
\hline
The Root-binder & A heretic planting acorns in the roller-pits & Let the forest reclaim the stone; let the cities sleep & He has a single green leaf growing from his left ear. \\
\hline
Ox-Master Senn & Handler of the eight-ox teams that drag the city & Keep the beasts healthy, the yokes oiled, and the pace steady & He names each ox after a dead mason. The oxen remember. \\
\hline
The Still Voice & A spirit of abandoned keystones & Witness the cities that chose to stop & It appears only when a city's migration clock reaches zero. It asks one question: ``Why did you stay?'' \\
\hline
Nomvula (Rain Queen) & Mythical rainmaker, keeper of the drought & Ensure the plateau never dries while the cities walk & She has not been seen in a century. Her staff is a petrified lightning bolt. \\
\hline
Tokoloshe Catcher & Old woman who traps water spirits & Keep the bearings free of malevolent tricksters & She catches tokoloshe in clay pots; they sing when the pot is shaken. \\
\hline
\end{longtable}

\subsection*{Names of Ngomebe}
\label{sec:ngomebe-names}

Inspired by the deep southern plateau: short, earth-toned syllables with clicks and tonal shifts. Surnames often denote a stone type, a city of origin, or a Mason-Kin House.

\begin{longtable}{|p{0.25\textwidth}|p{0.25\textwidth}|p{0.2\textwidth}|p{0.2\textwidth}|}
\hline
\textbf{First Names (Male)} & \textbf{First Names (Female)} & \textbf{Surnames} & \textbf{Epithets} \\
\hline
Morvath & Thandi & Keystone-carver & \textit{the Unstill} \\
\hline
Kaelen & Nia & Roller-born & \textit{the Grey-Faced} \\
\hline
Senn & Zola & Mason-Kin & \textit{the Slow Walker} \\
\hline
Duma & Amara & Stone-song & \textit{the Ninth Roller} \\
\hline
Temba & Lindi & Sledge-heir & \textit{the Root-Breaker} \\
\hline
Jengo & Penda & Ox-kin & \textit{the Warm Stone} \\
\hline
Nomvula & Lethabo & Rain-touched & \textit{the Drought's Daughter} \\
\hline
\end{longtable}

\paragraph*{(Quarry/Hall/Edge)} Abandoned quarry where stone-song still echoes; Mason-Kin hall with walls carved in lineage; the edge of a city's path — the stone is warm to the touch.

\section*{Spades --- Places}
\label{sec:ngomebe-places}
\begin{enumerate}
\setcounter{enumi}{1}
\item \textbf{City on the Move} --- Rollers groan, oxen low, the ground trembles.
  \hfill \textit{The city's shadow moves faster than the city itself. Stay in the shadow to keep pace.}
\item \textbf{Abandoned Quarry} --- Stone-song echoes from the walls.
  \hfill \textit{The song is the memory of the last stone cut here. It sings of a mason who fell.}
\item \textbf{Mason-Kin Hall} --- Walls carved with lineage, floor inlaid with sledge-ruts.
  \hfill \textit{The carvings weep water when a city is in danger. The water tastes of granite.}
\item \textbf{Edge of the Path} --- The stone is warm. The warmth fades when the city leaves.
  \hfill \textit{If you place your palm on the stone before it cools, you may ask the city one question.}
\item \textbf{Roller-Pit} --- A trench where the great rollers turn, greased with rendered ox-fat.
  \hfill \textit{The grease is mixed with ancestor ash. The rollers remember every mason who ever touched them.}
\item \textbf{Keystone Chamber} --- The heart of the city, a single stone that bears all weight.
  \hfill \textit{The keystone hums. The pitch changes when a mason swears a false oath.}
\item \textbf{Sledge-Track Scar} --- A groove worn into the bedrock by centuries of passage.
  \hfill \textit{The groove is deep enough to hide a corpse. Some say it already does.}
\item \textbf{Oxen Corral (Mobile)} --- A pen on sledges, dragged behind the city.
  \hfill \textit{The oxen never sleep. They walk in their sleep. Their dreams move the city.}
\item \textbf{Root-binder's Grove} --- A hidden stand of trees planted in an abandoned roller-pit.
  \hfill \textit{The trees have stone roots. The roots grip the rollers and will not let go.}
\item[J] \textbf{Abandoned City (Dead Keystone)} --- A stone shell on a hill, silent and cold.
  \hfill \textit{The keystone is cracked. A single green shoot grows from the crack. The Root-binder's work.}
\item[Q] \textbf{Bearing Foundry (Aeler)} --- A mobile forge where new rollers are cast.
  \hfill \textit{The forge is pulled by a separate team of oxen. It is always last in the procession.}
\item[K] \textbf{The Thing-Holm of Stone} --- A flat rock where Mason-Kin elders meet.
  \hfill \textit{The rock is a single piece of fallen sky-metal. Oaths sworn here cannot be broken.}
\item[A] \textbf{The First City (Myth)} --- A walking city that never stopped. It is said to circle the plateau, just beyond the horizon.
  \hfill \textit{Those who see it cannot speak of it for a year and a day. The stone seals their lips.}
\end{enumerate}

\paragraph*{(Cutter/Elder/Engineer)} Stone-cutter whose city left them behind; Mason-Kin elder, grey with granite dust; Aeler engineer: grease-stained, carrying a level.

\section*{Hearts --- People \& Factions}
\label{sec:ngomebe-people}
\begin{enumerate}
\setcounter{enumi}{1}
\item \textbf{Stone-Cutter (Left Behind)} --- Her city moved south. She stayed to carve a monument.
  \hfill \textit{She carves her own face into the rock. ``So the city will remember me,'' she says.}
\item \textbf{Mason-Kin Elder} --- Grey with granite dust, eyes the colour of slate.
  \hfill \textit{He carries a hammer that has never struck a false blow. The hammer knows.}
\item \textbf{Root-binder Heretic} --- Planting acorns in the roller-pits, singing to the earth.
  \hfill \textit{His hands are bark. His breath smells of green wood. The oxen fear him.}
\item \textbf{Aeler Engineer} --- Tuning fork in one hand, level in the other.
  \hfill \textit{She can hear a crack in a bearing from a mile away. She never walks under a keystone.}
\item \textbf{Ox-Master} --- Whip of braided leather, voice like a landslide.
  \hfill \textit{He names his oxen after dead masons. The oxen stop when he speaks their names.}
\item \textbf{Ancestral Healer (Izinyanga)} --- Reads the cracks in the keystone, speaks for the dead.
  \hfill \textit{She wears a necklace of ground quartz. The ancestors whisper through the stones.}
\item \textbf{Tokoloshe Catcher} --- Old woman, blind in one eye, carries a clay pot.
  \hfill \textit{She traps tokoloshe in her pots. They sing when shaken. She uses their songs to find sabotage.}
\item \textbf{Impundulu Tamer} --- A legendary bird-whisperer, with scars from lightning.
  \hfill \textit{He can call the lightning bird with a whistle. It demands payment: the blood of a mason.}
\item \textbf{Impi Captain} --- Warrior of the city's flank guard, carrying a shield of braided leather.
  \hfill \textit{She lost two fingers to a Root-binder ambush. She still holds her spear with the remaining three.}
\item[J] \textbf{The First Mason's Ghost} --- A translucent figure in a mason's apron.
  \hfill \textit{He appears when a keystone is first laid. He offers advice. The advice is always wrong.}
\item[Q] \textbf{The Ninth Ox} --- A mythic beast said to pull the First City.
  \hfill \textit{Its hooves never touch the ground. It walks on the shadow of clouds.}
\item[K] \textbf{Bearing-Caster (Aeler)} --- A dwarf who pours molten metal into roller-molds.
  \hfill \textit{He signs each bearing with his own thumbprint. The thumbprint is a prayer.}
\item[A] \textbf{The Still Voice} --- A spirit made of silence and cold stone.
  \hfill \textit{It asks: ``Why did you stop?'' Answer truthfully, or it will never let you leave.}
\end{enumerate}

\paragraph*{(Crack/Sabotage/Tremor)} A keystone cracks — the city stops; Root-binder sabotage — a wheel jams; an earth tremor (all actions Desperate for a scene).

\section*{Clubs --- Complications/Threats}
\label{sec:ngomebe-complications}
\begin{enumerate}
\setcounter{enumi}{1}
\item \textbf{Cracked Keystone} --- The heart of the city fractures. The city stops.
  \hfill \textit{The crack is singing. The song is the name of the mason who carved it. He is long dead.}
\item \textbf{Root-binder Sabotage} --- Acorns sprout in the roller-pits. A wheel jams.
  \hfill \textit{The roots are intelligent. They twist toward the sound of oaths.}
\item \textbf{Earth Tremor} --- The ground shakes. All actions are Desperate for a scene.
  \hfill \textit{The tremor was not natural. The Still Voice is waking.}
\item \textbf{Mason-Kin Feud} --- Two Houses argue over the next migration route.
  \hfill \textit{The argument is carved into the keystone. The stone is splitting.}
\item \textbf{Oxen Plague} --- A sickness spreads through the teams. Speed halves.
  \hfill \textit{The sickness came from a Root-binder's grove. The leaves are in the fodder.}
\item \textbf{Bearing Failure (Aeler)} --- A roller seizes. The engineer blames the stone-cutter.
  \hfill \textit{The bearing was sabotaged with sand from the Red Desert. The sand whispers.}
\item \textbf{Sledge-Track Collapse} --- The ground gives way. A section of the city sinks.
  \hfill \textit{The collapse reveals an older city buried beneath. The older city is still walking.}
\item \textbf{Migration Blockade} --- Another city has stopped on the planned route.
  \hfill \textit{The stopped city's keystone is dead. Its people have become Root-binders.}
\item \textbf{Ancestral Debt Called} --- A Mason-Kin House demands a sledge-right from a century ago.
  \hfill \textit{The debt is carved on a tablet of slate. The slate is warm. The ancestors are watching.}
\item[J] \textbf{Tokoloshe in the Bearings} --- A water spirit fouls the roller-pits. The city groans.
  \hfill \textit{The tokoloshe is small, fast, and invisible to the unaware. Only a catcher can see it.}
\item[Q] \textbf{Impundulu Strikes} --- A lightning bird shatters a keystone with a thunderous blow.
  \hfill \textit{The bird is angry. Someone broke a vow to the Rain Queen. The city bleeds stone.}
\item[K] \textbf{Root-binder Revolt} --- A grove of walking trees attacks the roller-pits.
  \hfill \textit{The trees have stone bark. The stones are carved with Mason-Kin faces.}
\item[A] \textbf{The First City Approaches} --- A city older than memory appears on the horizon.
  \hfill \textit{It is made of obsidian. Its rollers are giant femurs. It wants to merge.}
\end{enumerate}

\paragraph*{(Charter/Wax/Song)} Keystone charter (right to dwell in a city for a season); sledge-wax (greases negotiation with the Mason-Kin); stone-song fragment (one reroll on a survival test).

\section*{Diamonds --- Rewards/Leverage}
\label{sec:ngomebe-rewards}
\begin{enumerate}
\setcounter{enumi}{1}
\item \textbf{Keystone Charter} --- Right to dwell in a city for a season, tax-free.
  \hfill \textit{The charter is a chip of the keystone itself. Carry it to be welcomed.}
\item \textbf{Sledge-Wax} --- A lump of grey grease. Use it to gain +1 Position in Mason-Kin negotiations.
  \hfill \textit{The wax is rendered from the fat of an ox that walked a thousand miles.}
\item \textbf{Stone-Song Fragment} --- A recorded hum. Reroll a failed Survival test once.
  \hfill \textit{The song is the city's heartbeat. Hum it, and the stone will guide you.}
\item \textbf{Roller-Bearing (Aeler-crafted)} --- A spare bearing. Gift it to an engineer for a favour.
  \hfill \textit{The bearing is stamped with a rune that means ``no debt.'' The rune is a lie.}
\item \textbf{Mason-Kin Hammer} --- A tool that never misses a keystone's true face.
  \hfill \textit{The hammer's head is a piece of the First City's keystone. It whispers warnings.}
\item \textbf{Ox-Master's Whip} --- Crack it to make oxen move faster for one leg.
  \hfill \textit{The whip's braid includes a hair from every ox that ever pulled a city.}
\item \textbf{Root-binder's Acorn} --- Plant it to stop a city's roller for one day.
  \hfill \textit{The acorn grows into a sapling within an hour. The sapling has stone roots.}
\item \textbf{Abandoned City Map} --- A chart of dead cities and their keystone graves.
  \hfill \textit{The map is drawn on stone. The stone is a page. The pages are bound in granite.}
\item \textbf{Elder's Verdict} --- A ruling that settles a Mason-Kin dispute in your favour.
  \hfill \textit{The verdict is carved on a slate tablet. Carry it as a shield.}
\item[J] \textbf{Bearing-Foundry Pass} --- Access to the Aeler mobile forge for one repair.
  \hfill \textit{The pass is a single steel coin. The coin is still hot.}
\item[Q] \textbf{Still Voice's Silence} --- One question answered truthfully by the spirit of abandoned stone.
  \hfill \textit{The answer comes as a cold breath. The breath smells of geode crystals.}
\item[K] \textbf{First City's Echo} --- A piece of obsidian that hums with ancient rhythm.
  \hfill \textit{Once per campaign, touch it to make a city move one mile in a single night.}
\item[A] \textbf{The Living Keystone} --- A fragment from a city that still walks.
  \hfill \textit{Plant it in dead stone, and a new city will grow. The growth takes a century.}
\end{enumerate}

\section*{Quick use notes}
\label{sec:ngomebe-quick-use}
\begin{itemize}
\item Draw until you have all four suits: Spade = place, Heart = actor, Club = pressure, Diamond = leverage. Highest rank sets the main Clock (2–5 → 4, 6–10 → 6, J/Q/K → 8, A → 10).
\item Diamonds are codified outcomes (charters, wax, songs) that change position rather than call for a roll.
\item If any A appears, echo \textbf{stone-omens}—rollers that groan in harmony, keystones that weep water, a city that leaves no track.
\end{itemize}

\subsection*{Additional Features (Cultural Mechanics)}
\label{sec:ngomebe-mechanics}

\begin{itemize}
\item \textbf{The Keystone Bond (Relationship Mechanic):} Each walking city has a Keystone — a single stone that bears the weight of the roof. The Keystone is also the physical anchor of the city's ancestral spirit. A character who bonds with the Keystone (by spending a downtime scene in meditation at its base) gains:
  \begin{itemize}
    \item \textbf{Keystone Sense:} Once per scene, the character can ask the Keystone one question about the city's current state (``Is the eastern roller cracking?'' ``Is there a saboteur in the foundry?''). The Keystone answers truthfully, but only about the present.
    \item \textbf{Keystone Debt:} The character gains a Debt Clock (size 4) to the ancestor. When it fills, the ancestor demands a service — usually a pilgrimage to a forgotten quarry or the retrieval of a lost roller.
  \end{itemize}
\item \textbf{The Roller-Pits (Hazard Mechanic):} The city moves on rollers — massive wooden cylinders that must be greased and rotated. The roller-pits are dangerous, loud, and prone to sabotage.
  \begin{itemize}
    \item Entering a roller-pit requires a \textbf{Body + Athletics test (DV 4)} to avoid being crushed. On a failure, the character takes 1 Harm.
    \item Sabotaging a roller-pit (to stop the city or redirect it) requires a \textbf{Wits + Craft test (DV 5)}. On a success, the city moves in the desired direction; on a failure, the roller jams and the city stops — and the ancestors are angry.
    \item Repairing a jammed roller takes 3 actions. Each action requires a \textbf{Body + Craft test (DV 4)}. On a failure, the roller shifts and the character takes 1 Fatigue.
  \end{itemize}
\item \textbf{The Migration Clock (Campaign Clock, size 10):} Every walking city has a Migration Clock. It advances each day the city moves, and retreats when the city stops.
  \begin{itemize}
    \item At 2: The city enters fertile land. The party gains +1 die to foraging and hunting.
    \item At 4: The city passes a rival settlement. The party may trade or raid.
    \item At 6: The city enters a narrow pass. All navigation rolls are Desperate.
    \item At 8: The city approaches a sacred site. The ancestors grow restless — all rolls involving the Keystone gain +1 die, but failure ticks the Keystone Debt clock by 2.
    \item At 10: The city reaches its destination — the end of the migration. The ancestors are satisfied. The party may claim a \textbf{Keystone Blessing} (a permanent +1 die to all rolls involving stonework or negotiation with Ngomebe).
  \end{itemize}
  The party can influence the Migration Clock by:
  \begin{itemize}
    \item \textbf{Scouting ahead} (Wits + Survival) – success advances the clock by 1, failure ticks it back by 1.
    \item \textbf{Repairing rollers} (Craft test) – success advances the clock by 2, failure advances it by 0 but ticks the Keystone Debt.
    \item \textbf{Paying tribute} (a Honey Jar, a story, or a Boon) – advances the clock by 1 without risk.
  \end{itemize}
\item \textbf{The Root-binder Heresy (Faction Mechanic):} Root-binders are heretics who plant acorns in the roller-pits, hoping to stop the cities and let the forest reclaim the stone. They have a \textbf{Sabotage Clock (size 6)}. Each time the party fails a Craft or Survival roll in a roller-pit, the Sabotage Clock ticks.
  \begin{itemize}
    \item At 2: A roller cracks. The city loses 1 Speed (travel takes longer).
    \item At 4: A roller-pit floods. The city must stop for repairs (lose 1 day).
    \item At 6: A major collapse. The city cannot move until the party clears the pit (requires a Strength + Craft test, DV 5, and costs 1 day).
  \end{itemize}
  Players can work with the Root-binders (gain information, sabotage rival cities) or against them (hunt them, protect the rollers). The Root-binders offer \textbf{Green Blessings} – a permanent +1 die to survival in the forest, but at the cost of a Debt Clock to ``the forest itself'' (the Hollow's attention).
\item \textbf{The Mason-Kin Houses:} The walking cities are not nations but extended families. Each House controls a keystone, a migration route, and the rights to carve. Houses feud over sledge-tracks, water rights, and the honour of the First City.
\item \textbf{The Aeler Pact:} The Aeler engineers of the southern ranges maintain the bearings and capstans that make the walking cities possible. The Ngomebe provide keystone hearts and stone-song. It is a partnership of equals, sealed in slate and iron.
\item \textbf{Ancestral Cracks:} Once per session, a character may press their ear to a keystone and roll Spirit + Lore (DV 3). On success, the ancestors whisper a warning or a clue (GM reveals a hidden complication or a shortcut). On a miss, the stone is silent — but the ancestors are listening.
\item \textbf{Tokoloshe Pots:} A character who carries a clay pot with a tokoloshe inside gains +1 die to detect sabotage or traps in the city. The tokoloshe sings when danger is near. However, if the pot breaks, the tokoloshe escapes and will haunt the carrier for a full season (all actions Desperate at night).
\item \textbf{Local Patrons (Syncretic Names):}
  \begin{itemize}
    \item \textbf{The First Stone} — The Sealed Gate. *The keystone that never cracks; the oath that never breaks.*
    \item \textbf{The Rain Queen} — Raéyn. *The bringer of storms and the drought; the mistress of the path.*
    \item \textbf{The Lightning Bird} — The Storm King. *The shatterer of false keystones; the judge of masons.*
    \item \textbf{The Ancestor's Voice} — The Witness. *The whisper in the crack; the memory in the stone.*
  \end{itemize}
\end{itemize}

\subsection*{Lore Echoes (Cross-Expansion Hooks)}
\begin{itemize}
  \item \textbf{The Red Desert's Breath (Ashaan):} The Red Desert is spreading. Some Ngomebe cities are moving south to escape the dust. The Keystone Charter of a southern city may be the only thing standing between a tribe and annihilation.
  \item \textbf{The Sledge-Wax Trade (Way of Silk):} Sledge-wax is a valuable commodity on the caravan routes. Tulkani merchants trade it for story-knots; Fhara water-clerks use it to seal well-contracts. A missing shipment of sledge-wax could stop a city.
  \item \textbf{The Iron Mines of Vajrapur (Dhahara):} The Aeler engineers of Ngomebe source their bearing-metal from the Dhaharan iron mines. A dispute between a Mason-Kin House and a Dhaharan raja could choke the supply of rollers.
  \item \textbf{The Hollow and the Still City (Eastern Realms):} The Hollow is attracted to still places. A city that stops moving becomes a beacon. The Aelaerem whisper that the First City never stopped because it is running from something.
  \item \textbf{Tokoloshe and the Bitter Root:} Daughters of the Cord sometimes capture tokoloshe and weave them into Hollowed children. The tokoloshe inside the child sings when a debt is due. The Ngomebe avoid Daughter-run hospices for this reason.
\end{itemize}

\subsection*{Ngomebe Superstitions (burn for +1 die once)}
\begin{itemize}
  \item \textbf{Bread to the Roller:} Crumble a crust into the grease-pit. Ignore the first \emph{Earth Tremor} penalty. `[SUPERSTITION]`
  \item \textbf{Knot the Sledge-Rope:} Tie a single knot in an ox's halter. Cancel \emph{Oxen Plague} or take –1 die on your next handling roll (your choice). `[SUPERSTITION]`
  \item \textbf{Name to the Keystone:} Whisper a dead mason's name into the keystone's crack. Gain +1 Effect when negotiating with the Mason-Kin this scene, but mark \emph{Obligation} to the First City. `[SUPERSTITION]`
\end{itemize}

\subsection*{Thematic SB Spend Table}
\label{sec:ngomebe-sb}

\paragraph{Minor Complications (1 SB)}
\begin{itemize}
\item \textbf{Exposure:} Your actions draw unwanted attention from \textbf{Mason-Kin guards or Aeler engineers}.
\item \textbf{Noise:} Sounds of your actions alert nearby \textbf{ox-drivers or stone-cutters}.
\item \textbf{Trace:} Evidence of your passage marks your route for \textbf{Root-binders or keystone guardians}.
\item \textbf{Delay:} A brief but meaningful setback costs you \textbf{time or favorable migration pace}.
\item \textbf{Supply Strain:} Mark +1 segment on a relevant \textbf{resource clock}.
\end{itemize}

\paragraph{Moderate Setbacks (2 SB)}
\begin{itemize}
\item \textbf{Alarm Raised:} \textbf{Local Mason-Kin elder or Aeler vent-prior} becomes aware and begins responding.
\item \textbf{Position Lost:} You lose advantageous ground/cover/stealth due to \textbf{earth tremor or roller jam}.
\item \textbf{Foe Appears:} A \textbf{Root-binder cell or rival House guard} arrives on scene.
\item \textbf{Gear Trouble:} A piece of equipment becomes \textbf{Compromised/Neglected}.
\item \textbf{Lock/Barrier:} A simple obstacle now requires a test to overcome.
\end{itemize}

\paragraph{Serious Trouble (3 SB)}
\begin{itemize}
\item \textbf{Reinforcements:} Additional \textbf{Mason-Kin, Root-binders, or Aeler engineers} arrive.
\item \textbf{Key Gear Breaks:} A crucial tool/weapon becomes temporarily unusable.
\item \textbf{Major Twist:} The situation fundamentally changes—\textbf{keystone cracks/city stops/First City sighted}.
\item \textbf{Rail Tick:} Advance a relevant campaign/front clock by 1 segment.
\item \textbf{Condition Applied:} Mark \textbf{Fatigue 1/Harm 1/Condition} appropriate to fiction.
\end{itemize}

\paragraph{Major Turns (4+ SB)}
\begin{itemize}
\item \textbf{Trap Springs:} A prepared danger activates with full effect.
\item \textbf{Authority Arrival:} \textbf{Mason-Kin High Council, Still Voice, or First City's herald} intervenes.
\item \textbf{Scene Shift:} The environment changes dramatically—\textbf{tremor opens a chasm/city splits in two/keystone begins to sing}.
\item \textbf{Patron Omen:} Divine/arcane forces take notice—\textbf{omen appears/blessing lost/curse manifests}.
\item \textbf{Narrative Pivot:} The story takes an unexpected turn that reframes objectives.
\end{itemize}

\paragraph{Region-Specific SB Options}
\begin{itemize}
\item \textbf{Ngomebe (Stone Law):} Keystone hums a warning, roller-grooves glow faintly, the ground trembles without quake.
\item \textbf{Ngomebe (Root-binder Heresy):} Acorns sprout from cracks, vines write warnings in the dust, trees move against the wind.
\item \textbf{Ngomebe (Aeler Pact):} Bearings click in counterpoint, tuning forks ring without being struck, the smell of hot metal hangs in still air.
\item \textbf{Ngomebe (Ancestral Spirits):} A cold breath whispers a name, the keystone weeps water, a tokoloshe laugh echoes from a dry roller-pit.
\end{itemize}

\subsection*{Pursuit Ladder (Near / Mid / Far)}
\begin{itemize}
\item \textbf{Advance/Withdraw (Wits + Survival):} On a hit, shift one band. On a strong hit, also impose \emph{False Trail}.
\item \textbf{Cut the Sledge (Presence + Sway):} On a hit, freeze range for one exchange; if a keystone charter is in your possession, +1 die.
\item \textbf{Weather the Tremor (Body + Athletics):} On a hit in \emph{Earth Tremor}, you pick who stays standing; on a miss, the GM does.
\end{itemize}

\subsection*{Boss Seeds (Stone \& Root)}
\begin{itemize}
\item \textbf{The Root-King (Nine Groves)} — \emph{Phases}: Acorn Sabotage $\rightarrow$ Orchestrated Stillness $\rightarrow$ Forest Coup. \emph{Cracks}: burn a grove, expose a false acorn, survive a landslide.
\item \textbf{The Still Voice (Silence of Stone)} — \emph{Phases}: Abandoned Keystones $\rightarrow$ Procession of Dead Cities $\rightarrow$ Stillness Baptism of an Elder. \emph{Cracks}: relic stone-song, survivor testimony, the city moving against the Voice.
\item \textbf{The Tokoloshe King (Master of Tricksters)} — \emph{Phases}: Bearings Foul $\rightarrow$ Roller Pit Flooded $\rightarrow$ Oxen Stampede. \emph{Cracks}: trap him in a clay pot, negotiate with a pot of sorghum beer, feed him a false name.
\end{itemize}

\begin{tcolorbox}[colback=black!3,colframe=black!40!white,title={Patronage \& Power}]
Among the Ngomebe, power flows through keystones, migration rights, and the loyalty of the Mason-Kin Houses. An elder's authority rests on the age of his House's keystone and the number of cities that follow his route. The Aeler engineers hold technical authority; the Root-binders offer a heretical alternative; and the Ancestral Healers speak for the dead, who still vote through the cracks in the stone.

\textbf{For the GM:}  
Power in Ngomebe society is measured in movement and stone. Patronage often takes the form of keystone charters, sledge-wax, or stone-song fragments—each tied to a specific city and House. To emphasize this:
\begin{itemize}
\item Tie rewards to physical tokens (chips of keystone, lumps of wax, recorded songs) that can be challenged, stolen, or voided.
\item Let rival Houses issue conflicting claims, forcing players to choose whose city shelters them.
\item Use the roller-pits, keystone chambers, and Mason-Kin halls as arenas for social contests, where migration routes are debated and saboteurs are judged.
\item Introduce \textbf{Ancestral Cracks} as a way for players to gain clues from the dead — but at the cost of the ancestors' attention.
\end{itemize}
In Ngomebe society, the stone remembers everything — including your debts. The ancestors do not forget. The tokoloshe does not forgive. And the city that stops dies.
\end{tcolorbox}

% Index entries
\index{Ngomebe|see{Stone Men}}
\index{Theme generator}
\index{Stone Men generator}
\index{Places!Ngomebe}
\index{People!Ngomebe}
\index{Complications!Ngomebe}
\index{Rewards!Ngomebe}
\index{Mason-Kin}
\index{Root-binder}
\index{Aeler engineers}
\index{Keystone}
\index{Walking city}
\index{First City}
\index{Still Voice}
\index{Ox-Master}
\index{Stone-song}
\index{Sledge-wax}
\index{Roller-bearing}
\index{Ancestral Healers}
\index{Izinyanga}
\index{Tokoloshe}
\index{Impundulu}
\index{Rain Queen}
\index{Local Patrons!Ngomebe}
\newpage

% -------------------- TAHARKA --------------------
\chapter{Taharka}
\emph{The Monsoon Crown.}

\noindent
Kofi admits: \textit{“The passes are a graveyard for armies. The horses speak in riddles. The coffee ritual has three cups – and a fourth, silent one, for the spirits. I have never understood it. I just pay my toll and keep my mouth shut.”}

\noindent
Taharka grows rich on monsoon trade – spices, coffee, incense – and has never been conquered. Its high passes are a graveyard for Ashaani armies, a barrier that Oshiiran grain barges cannot cross. The kingdom is a fortress of rock and thin air, where every valley speaks a different dialect and every ridge has a watchtower. The horse is not a mount; it is a lineage, a treaty, a weapon. And the old spirits – \textit{mizimu}, \textit{kichocheo}, \textit{pepo} – are not legends. They are neighbours. Be polite. Count to eight. Never pour the fourth cup.

\vspace{0.5em}
\noindent\textit{“Never refuse the second cup of coffee. The first is welcome; the second is a test of honour; the third is a contract. Refuse the second, and you have insulted the house – and the house’s ancestors, who listen through the floorboards.”} \\
— Ridge‑Watcher Captain Anwar

\vspace{1em}
\begin{almanac}
\textbf{What do people eat for breakfast?} Barley flatbread and spiced tea – the bread baked on iron griddles, the tea brewed with cardamom and mountain herbs that only the Ridge‑Watchers know.

\textbf{What does a child learn before they can walk?} To read the wind. The eight winds have names. The ninth wind has no name – it is the mountain’s breath. The mountain’s breath is death.

\textbf{What is the first thing you notice about a stranger?} Their breath. The air is thin. Fast breathers have not acclimated. Slow breathers are hiding something. Those who do not breathe are spirits.

\textbf{What is the most common crime?} Horse theft – a declaration of war. A stolen horse is a stolen ancestor. The thief’s clan pays in blood or trade.

\textbf{What does no one talk about but everyone knows?} The Crown Council has not met in ten years. The passes are closed by snow, by war, by the simple fact that the councillors cannot agree on a meeting place. The Ridge‑Watchers rule by the mountain’s voice.

\textbf{Where do people go when they die?} To the high passes. Bodies laid on stone platforms, wrapped in woollen shrouds. The wind carries their names to the peaks. The peaks remember.

\textbf{How do you apologize for a serious wrong?} You offer a horse – the best in your stable. If the wronged party accepts, the debt is paid. If they kill the horse, the debt is blood. Blood is never repaid.

\textbf{What is the most important festival?} The First Rain (monsoon’s arrival). The passes open, the caravans descend, the lowland ports receive the first spice ships. The celebration lasts nine days. The ninth day is silent – for the dead of the passes.

\textbf{What is the secret that could destroy a powerful person?} The Crown Council’s missing member is not missing. He is buried under a landslide that was not an accident. The Ridge‑Watchers caused it – because he learned the passes’ secret: they are not natural. Something older carved them, and the Ridge‑Watchers are its servants.
\end{almanac}

\subsection*{The Aelaerem of the Southern Hills}
The Aelaerem – pygmy halfling coffee‑tenders of the highest passes – are described in full with their own generator tables at the end of this chapter (see the \textbf{Aelaerem of the Southern Hills} subsection). Their red threads, ninth‑cup rites, and mist‑shrouded terraces are a constant reminder that the smallest folk often hold the oldest keys.

\newpage
\section{Taharka --- ``The Monsoon Crown''}
\label{chap:taharka}

\begin{quote}
  \textbf{Ridge-Watcher Captain (Elite):}  
  ``The passes remember every army that tried to climb them. The mountains do not forgive haste, and they do not forget a foreign banner. We do not guard the highlands. We listen to them. They tell us who to let through.''
\end{quote}

\begin{quote}
  \textbf{Horse Trader (Commoner):}  
  ``Before you praise the rider, praise the horse. Before you praise the horse, praise the mountain that grew its lungs. A Taharkan foal is born breathing thin air. That is why our cavalry has never been broken.''
\end{quote}

\begin{tcolorbox}[colback=black!3,colframe=blue!60!black,title={Theme \& Atmosphere}]
Taharka grows rich on monsoon trade — spices, coffee, and incense that travel east to the distant ocean and west to the Gulf. The kingdom has never been conquered; its passes are a graveyard for Ashaani armies and a barrier that Oshiiran grain barges cannot cross. The highlands are a fortress of rock and thin air, where every valley speaks a different dialect and every ridge has a watchtower. On the coast, the monsoon winds bring ships from Dhahara, Ayokha, and the eastern isles, and a trade creole has grown from a dozen tongues — the language of the wharves, where a merchant can bargain in verse and a captain can curse in seven accents.

\vspace{2mm}
\textbf{What you notice first:} The smell of roasting coffee and frankincense. The way the air thins as you climb, until each breath is a measured gift. The constant lowing of pack mules on terraced paths.

\textbf{The one rule that kills newcomers:} Never refuse a second cup of tea. The first is hospitality; the second is a test of honour. The third is a promise. Refuse the second, and you have insulted the house.
\end{tcolorbox}

\subsection*{Regional Mood: The Weight of the Pass}

Power here is measured in passes, horse lines, and the favour of the monsoon. The Crown Council rules from the high plateau, but each valley has its own lord, each pass its own toll-keeper. The Ridge-Watchers are the eyes of the kingdom, stationed in rock-chapels that double as signal towers. The horse is not a mount; it is a lineage, a treaty, a weapon. To control the highland breeds is to control the passes — and to control the passes is to remain free.

\textbf{Starting Location:} A rock-chapel at the head of the Silver Pass, where the air is cold and the wind smells of snow. A Ridge-Watcher captain blows a horn of ibex horn: ``The Crown Council has declared a muster. The passes are closing to all foreign traffic. You have until dusk to declare your business — or turn back.''

\subsection*{Prominent NPCs of Taharka}
\label{sec:taharka-npcs}

\begin{longtable}{|p{0.2\textwidth}|p{0.25\textwidth}|p{0.3\textwidth}|p{0.15\textwidth}|}
\hline
\textbf{Name} & \textbf{Role/Title} & \textbf{Motivation} & \textbf{Rumoured Quirk} \\
\hline
Ridge-Watcher Captain Anwar & Commander of the Silver Pass garrison & Keep the passes closed to Ashaani spies & He has not descended below the snow line in twenty years. \\
\hline
Merchant-Queen Safiya & Spice House matriarch of the western coast & Control the monsoon trade before the Crown taxes it & She speaks the trade creole with a coastal accent that no one can place. \\
\hline
Horse Lord Tamrat & Breeder of the royal stables & Sell a foal to the Dhaharan court for a pass-right & He names his horses after dead poets. The horses recite verses in their gait. \\
\hline
Sister Hana & Covenant judge of the highland circuit & Resolve disputes between lords without drawing Crown attention & She carries a water-clock that runs on melted snow. \\
\hline
The Astrologer of the Monsoon & Crown advisor who reads the winds & Predict the first rains to set the trade season & He has not seen the sun in thirty years. He works by candlelight and omen. \\
\hline
The Pass Ghost & A spirit of a soldier who froze on watch & Warn travellers of avalanches and ambushes & It appears as a column of steam in the cold air. It points with a skeletal finger. \\
\hline
\end{longtable}

\subsection*{Names of Taharka}
\label{sec:taharka-names}

Inspired by the highland plateau and the coastal creole: melodic, with doubled consonants and honorific prefixes. Surnames often denote a clan, a pass, or a horse-line.

\begin{longtable}{|p{0.25\textwidth}|p{0.25\textwidth}|p{0.2\textwidth}|p{0.2\textwidth}|}
\hline
\textbf{First Names (Male)} & \textbf{First Names (Female)} & \textbf{Surnames} & \textbf{Epithets} \\
\hline
Tamrat & Safiya & Berhanu & \textit{the Ridge-Born} \\
\hline
Anwar & Hana & Tesfaye & \textit{the Horse-Tongue} \\
\hline
Dawit & Amara & Gebre & \textit{the Monsoon-Watcher} \\
\hline
Kassim & Lij & Mekonnen & \textit{the Ninth Pass} \\
\hline
Solomon & Tsion & Desta & \textit{the Coffee-Blade} \\
\hline
Yonas & Selam & Asfaw & \textit{the Thin-Air} \\
\hline
\end{longtable}

\paragraph*{(Farm/Chapel/Pass)} Terraced farm with coffee shrubs on every step; rock-chapel with walls that drip condensation; the high pass where snow lingers even in summer.

\section*{Spades --- Places}
\label{sec:taharka-places}
\begin{enumerate}
\setcounter{enumi}{1}
\item \textbf{Terraced Farm (Coffee)} --- Shrubs planted on stone-step shelves.
  \hfill \textit{The steps are a thousand years old. Each step is a different family's debt.}
\item \textbf{Rock-Chapel (Dripping)} --- A cave-temple where the walls weep condensation.
  \hfill \textit{The water is holy. Drink it, and you will remember a past life. The water tastes of iron.}
\item \textbf{High Pass (Summer Snow)} --- A saddle between peaks where the wind never stops.
  \hfill \textit{The snow is blue. It crunches like glass underfoot. Do not eat it.}
\item \textbf{King's Stable} --- A long building of fitted stone, horses in brass halters.
  \hfill \textit{Each stall has a nameplate. The nameplate is a defeated enemy's shield.}
\item \textbf{Spice Warehouse (Coast)} --- A wood-and-coral building that smells of cardamom.
  \hfill \textit{The floor is layered with spices that fell from broken sacks. The layers are a century thick.}
\item \textbf{Monsoon Wharf} --- A stone pier that is underwater half the year.
  \hfill \textit{The pier is carved with tide marks. Each mark is a year when a ship never returned.}
\item \textbf{Crown Council Hall} --- A roundhouse of stone and horn, built into a hillside.
  \hfill \textit{The hall has nine doors. Each door leads to a different valley. The ninth door is sealed.}
\item \textbf{Horse Cemetery} --- A field of cairns, each topped with a bronze bell.
  \hfill \textit{The bells ring when a foal is born to the same bloodline. The dead horses speak through the bells.}
\item \textbf{Ridge-Watcher Tower} --- A stone needle on a peak, with a signal fire that never goes out.
  \hfill \textit{The fire burns blue when the passes are closed. It burns green when a foreign army is sighted.}
\item[J] \textbf{Trade Creole School} --- A thatched hut where children learn the coastal tongue.
  \hfill \textit{The teacher is a one-eyed sailor. He signs his lessons with knots, not words.}
\item[Q] \textbf{The Astrologer's Observatory} --- A roofless tower with a bronze armillary sphere.
  \hfill \textit{The sphere has nine rings. The ninth ring is broken. The astrologer refuses to fix it.}
\item[K] \textbf{The Pass of the Ninth Army} --- A defile where an Ashaani host froze to the rock.
  \hfill \textit{On cold nights, you can still see their shapes in the ice. They are still standing.}
\item[A] \textbf{The King's Tomb (Hidden)} --- A cave sealed with a single keystone.
  \hfill \textit{The keystone has a single word carved in the trade creole: ``Wait.''}
\end{enumerate}

\paragraph*{(Trader/Watcher/Envoy)} Horse trader with a copper bell; Ridge-Watcher captain with a bow over his shoulder; Spice House merchant in robes smelling of cardamom.

\section*{Hearts --- People \& Factions}
\label{sec:taharka-people}
\begin{enumerate}
\setcounter{enumi}{1}
\item \textbf{Horse Trader} --- Copper bell on his belt, a ledger of bloodlines in his head.
  \hfill \textit{He names his price in coffee, not coin. The coffee is grown on his own terrace.}
\item \textbf{Ridge-Watcher Captain} --- Bow of horn and sinew, arrows fletched with eagle feathers.
  \hfill \textit{He can draw a bow at full gallop. He has never missed.}
\item \textbf{Spice House Merchant} --- Robes stained with cinnamon, fingers yellow with turmeric.
  \hfill \textit{She can identify a spice's origin by taste. She has been wrong twice. Both times, the ship sank.}
\item \textbf{Crown Council Envoy} --- Careful, neutral, carrying a staff of office.
  \hfill \textit{He speaks in questions. ``Would the honoured traveller perhaps consent to...''}
\item \textbf{Coffee Farmer} --- Hands stained black from hulling beans, eyes crinkled from sun.
  \hfill \textit{He offers a cup to every passerby. The cup is clay. The clay is from a sacred spring.}
\item \textbf{Monsoon Captain} --- A sailor with a brass ring in one ear and a scar across his throat.
  \hfill \textit{He speaks the trade creole with a Dhaharan accent. His ship is named ``Ninth Wind.''}
\item \textbf{Pass Guide} --- A hooded figure with a staff tipped with a horse skull.
  \hfill \textit{She never speaks above a whisper. The wind carries her voice exactly as far as needed.}
\item \textbf{Sister Hana (Covenant Judge)} --- Grey robe, water-clock, a scale that never balances.
  \hfill \textit{She has ruled on nine border disputes. The losing party always apologizes.}
\item \textbf{Horse Lord Tamrat} --- Grey beard braided with horsehair, missing two fingers.
  \hfill \textit{The missing fingers were bitten off by a foal. He kept the foal.}
\item[J] \textbf{The Astrologer} --- Blind, pale, wrapped in wool. He smells of wax and ozone.
  \hfill \textit{He predicts the monsoon by the pulse of a spider in a jar. The spider has never been wrong.}
\item[Q] \textbf{The Pass Ghost} --- A column of steam in the cold air, shaped like a soldier.
  \hfill \textit{It points at travellers. The direction indicates the safest route. It never points the same way twice.}
\item[K] \textbf{The King (Offstage)} --- Never seen, rarely heard, but his name opens every door.
  \hfill \textit{Rumour says he is a horse. Rumour says he is a spirit. Rumour says he is a child of nine.}
\item[A] \textbf{The Ninth Heir (Exiled)} --- A prince who lost a succession war, now living as a horse-trader.
  \hfill \textit{He carries a seal made of bone. The bone was his father's thumb.}
\end{enumerate}

\paragraph*{(Landslide/Sickness/Assassination)} A landslide blocks the pass (delay 2 segments); horse sickness spreads — all mounts lose 1 Speed; the king is assassinated — passes close to all foreign traffic.

\section*{Clubs --- Complications/Threats}
\label{sec:taharka-complications}
\begin{enumerate}
\setcounter{enumi}{1}
\item \textbf{Landslide Block} --- A rockslide covers the pass. Delay 2 segments.
  \hfill \textit{The rocks are warm. The landslide was caused by a signal fire on the peak.}
\item \textbf{Horse Sickness} --- A cough spreads through the stables. All mounts lose 1 Speed.
  \hfill \textit{The sickness came from a Dhaharan ship. The captain is a spy.}
\item \textbf{King Assassinated} --- The Crown Council declares a succession crisis. All passes close.
  \hfill \textit{The assassin was a horse. The horse was poisoned. The poison was in the coffee.}
\item \textbf{Ashaani Raiders (Disguised)} --- A caravan of ``merchants'' is actually a Veiled Sister's scout team.
  \hfill \textit{They wear Taharkan dress but do not remove their shoes indoors. The innkeeper noticed.}
\item \textbf{Coffee Blight} --- A fungus attacks the terraces. Prices triple. Starvation looms.
  \hfill \textit{The blight was spread by a trader from the coast. The trader's ship flew no flag.}
\item \textbf{Monsoon Delay} --- The winds shift. The first ships are three weeks late.
  \hfill \textit{The delay was predicted by the astrologer. He warned the Crown. The Crown did not listen.}
\item \textbf{Spice Warehouse Fire} --- A stack of cinnamon ignites. The whole wharf is at risk.
  \hfill \textit{The fire was set by a rival House. The House has a monopoly on pepper.}
\item \textbf{Trade Creole Dispute} --- A contract is voided because of a mistranslated word.
  \hfill \textit{The word means ``shipment'' in one dialect and ``hostage'' in another. The merchant knew.}
\item \textbf{Ridge-Watcher Betrayal} --- A captain sells the pass schedule to an Ashaani agent.
  \hfill \textit{The captain's daughter is a spirit-broken slave in Galanina. He had no choice.}
\item[J] \textbf{The Pass Ghost's Warning} --- The ghost points at a traveller. The traveller will die before sunset.
  \hfill \textit{The traveller is a party member. The ghost is pointing at a shadow, not a person.}
\item[Q] \textbf{Crown Council Schism} --- The nine lords cannot agree on a successor. Civil war looms.
  \hfill \textit{The ninth lord is a child. The child speaks only the trade creole. No one understands.}
\item[K] \textbf{Horse Line Extinction} --- A plague kills the last mare of the royal bloodline.
  \hfill \textit{The mare's foal was stolen by a rival kingdom. The foal is in a Dhaharan stable.}
\item[A] \textbf{The Ninth Pass Opens} --- A sealed defile cracks open. An ancient army of ice-walkers emerges.
  \hfill \textit{The army is made of frozen Ashaani corpses. They are still looking for a way home.}
\end{enumerate}

\paragraph*{(Bond/Stone/Pass)} Horse-bond — a foal promised (future String); pass-stone — crossing a high pass at half toll; coffee measure — trade goods at +1 Position for one scene.

\section*{Diamonds --- Rewards/Leverage}
\label{sec:taharka-rewards}
\begin{enumerate}
\setcounter{enumi}{1}
\item \textbf{Horse-Bond (Promised Foal)} --- A future String. The foal has not yet been born.
  \hfill \textit{The bond is a braid of horsehair. Wear it, and the mare will remember you.}
\item \textbf{Pass-Stone} --- A carved token. Cross any high pass at half toll.
  \hfill \textit{The stone is a piece of the Pass Ghost's staff. It is cold to the touch.}
\item \textbf{Coffee Measure (Royal Seal)} --- A clay cup. Trade goods at +1 Position for one scene.
  \hfill \textit{The cup was used by the king. He drank from it once. The crack is from his tooth.}
\item \textbf{Ridge-Watcher's Horn} --- Blow it to summon a guide from the nearest tower.
  \hfill \textit{The horn is made from an ibex that the captain killed with a single arrow.}
\item \textbf{Spice House Letter of Credit} --- 5 Measures of trade value, redeemable at any coast warehouse.
  \hfill \textit{The letter is sealed with cinnamon oil. The seal smells of betrayal.}
\item \textbf{Crown Council Verdict} --- A ruling in your favour on any border dispute.
  \hfill \textit{The verdict is carved on a slate tile. The tile is also a signal flag.}
\item \textbf{Monsoon Schedule (Astrologer's Copy)} --- Predicts safe sailing days for the next season.
  \hfill \textit{The schedule is written on a dried spiderweb. The web was the astrologer's pet.}
\item \textbf{Trade Creole Phrasebook} --- Gain +1 die to all social rolls on the coast for one journey.
  \hfill \textit{The phrasebook is a string of knots. Each knot is a curse.}
\item \textbf{Horse Lord's Favour} --- One foal from the royal stables, delivered in a year.
  \hfill \textit{The favour is a single white hair from the stallion's tail. The stallion is dead.}
\item[J] \textbf{Pass Ghost's Pointing} --- One safe route through any high pass, guaranteed.
  \hfill \textit{The ghost writes the route in frost on a rock. The frost melts at noon.}
\item[Q] \textbf{The Astrologer's Spider} --- A jarred spider that pulses with the monsoon's heartbeat.
  \hfill \textit{Once per campaign, consult the spider to predict a storm with perfect accuracy.}
\item[K] \textbf{The Ninth Heir's Seal} --- A bone thumbprint. Commands loyalty from the exiled prince's network.
  \hfill \textit{The seal is also a key. It opens a lock in the Crown Council Hall. The lock is the ninth door.}
\item[A] \textbf{The King's Blessing} --- A royal decree written in trade creole. One national favour.
  \hfill \textit{The decree is written on a horse's hoof. The hoof was the king's favourite mount.}
\end{enumerate}

\section*{Quick use notes}
\label{sec:taharka-quick-use}
\begin{itemize}
\item Draw until you have all four suits: Spade = place, Heart = actor, Club = pressure, Diamond = leverage. Highest rank sets the main Clock (2–5 → 4, 6–10 → 6, J/Q/K → 8, A → 10).
\item Diamonds are codified outcomes (bonds, stones, measures) that change position rather than call for a roll.
\item If any A appears, echo \textbf{monsoon-omens}—horses that speak in riddles, coffee that tastes of snow, a pass that closes behind you.
\end{itemize}

\subsection*{Additional Features}
\begin{itemize}
\item \textbf{The Ridge-Watchers:} A network of signal towers on every peak. A fire, a mirror, or a horn can send a message across the entire kingdom in a day. The Ridge-Watchers are neutral in Crown politics; they serve only the passes.
\item \textbf{The Trade Creole:} A language born on the monsoon wharves, blending words from a dozen coastal tongues. It has no written form; contracts are sealed with knots or witnessed by a Lethai-ar. The Creole is the tongue of negotiation, insult, and poetry — and a single mistranslation can start a war.
\item \textbf{The Horse Lineages:} Taharkan horses are bred for thin air and steep trails. Each bloodline has a name, a song, and a debt. To own a Taharkan horse is to inherit its lineage's obligations — and its enemies.
\end{itemize}

\subsection*{Lore Echoes (Cross-Expansion Hooks)}
\begin{itemize}
  \item \textbf{The Monsoon Trade (Dhahara/Ayokha):} Taharka's spice and coffee travel east on the monsoon winds, reaching the Dhaharan coast and the Ayokhan river markets. A delay in Taharka means famine in a dozen ports.
  \item \textbf{The Pass Ghost and the Hollow (Eastern Realms):} The Pass Ghost is a spirit of the threshold — a border between living and dead. The Aelaerem say it is a cousin to the Hollow, bound by an oath to the first king.
  \item \textbf{Ashaani Spies and the Veiled Sisters:} Ashaani agents often pose as spice merchants. The Ridge-Watchers have learned to test them by offering tea — a foreigner who does not know the ritual of three cups reveals themselves.
  \item \textbf{The Ninth Heir and the Kalrantha (Dhahara):} The exiled prince is rumoured to have bought his horse from a Dhaharan smuggler — a horse that carried a certain green glass amulet. The Crown Council wants to know where the amulet is now.
\end{itemize}

\subsection*{Taharka Superstitions (burn for +1 die once)}
\begin{itemize}
  \item \textbf{Bread to the Pass:} Crumble a crust at the highest point of the pass. Ignore the first \emph{Landslide Block} penalty.
  \item \textbf{Knot the Horse's Tail:} Tie a single knot in a horse's tail hair. Cancel \emph{Horse Sickness} or take –1 die on your next riding roll (your choice).
  \item \textbf{Name to the Coffee:} Whisper a dead Ridge-Watcher's name into a coffee cup. Gain +1 Effect when negotiating passage this scene, but mark \emph{Obligation} to the Crown.
\end{itemize}

\begin{tcolorbox}[colback=black!3,colframe=black!40!white,title={Patronage \& Power}]
Among the Taharka, power flows through passes, horse lines, and the favour of the monsoon. The Crown Council rules by consensus, but each valley lord commands his own Ridge-Watchers. The spice merchants of the coast operate independently, paying tariffs to the Crown but answering to no single lord.

\textbf{For the GM:}  
Power in Taharkan society is measured in altitude and access. Patronage often takes the form of pass-stones, horse-bonds, or coffee measures—each tied to a specific valley, bloodline, or trade route. To emphasize this:
\begin{itemize}
\item Tie rewards to physical tokens (carved stone, braided hair, clay cups) that can be challenged, stolen, or voided.
\item Let rival lords issue conflicting pass-rights, forcing players to choose whose path to follow.
\item Use the rock-chapels, horse cemeteries, and monsoon wharves as arenas for social contests, where hospitality is law and a cup of tea is a contract.
\end{itemize}
In Taharkan society, the mountain remembers every step—and the horse remembers every rider.
\end{tcolorbox}

% Index entries
\index{Taharka|see{Monsoon Crown}}
\index{Theme generator}
\index{Monsoon Crown generator}
\index{Places!Taharka}
\index{People!Taharka}
\index{Complications!Taharka}
\index{Rewards!Taharka}
\index{Ridge-Watchers}
\index{Crown Council}
\index{Horse lineages}
\index{Trade creole}
\index{Monsoon trade}
\index{Pass Ghost}
\index{Coffee terraces}
\index{Rock-chapels}
\index{Horse Lord}
\index{Astrologer of the Monsoon}
\index{Ninth Pass}
\index{Saint Adar!Taharka}
\index{Saint Hervor!Taharka}
\index{Tide-Washed Covenant!Taharka}

\section*{Thematic SB Spend Table}
\label{sec:taharka-sb}

\subsection*{Minor Complications (1 SB)}
\begin{itemize}
\item \textbf{Exposure:} Your actions draw unwanted attention from \textbf{Ridge-Watchers or Crown Council envoys}.
\item \textbf{Noise:} Sounds of your actions alert nearby \textbf{horse guards or spice warehouse clerks}.
\item \textbf{Trace:} Evidence of your passage marks your route for \textbf{pass guides or rival horse traders}.
\item \textbf{Delay:} A brief but meaningful setback costs you \textbf{time or favourable monsoon window}.
\item \textbf{Supply Strain:} Mark +1 segment on a relevant \textbf{resource clock}.
\end{itemize}

\subsection*{Moderate Setbacks (2 SB)}
\begin{itemize}
\item \textbf{Alarm Raised:} \textbf{Local Ridge-Watcher captain or Spice House factor} becomes aware and begins responding.
\item \textbf{Position Lost:} You lose advantageous ground/cover/stealth due to \textbf{landslide or fog}.
\item \textbf{Foe Appears:} A \textbf{rival horse lord or Ashaani agent} arrives on scene.
\item \textbf{Gear Trouble:} A piece of equipment becomes \textbf{Compromised/Neglected}.
\item \textbf{Lock/Barrier:} A simple obstacle now requires a test to overcome.
\end{itemize}

\subsection*{Serious Trouble (3 SB)}
\begin{itemize}
\item \textbf{Reinforcements:} Additional \textbf{Ridge-Watchers, Crown guards, or spice house enforcers} arrive.
\item \textbf{Key Gear Breaks:} A crucial tool/weapon becomes temporarily unusable.
\item \textbf{Major Twist:} The situation fundamentally changes—\textbf{king assassinated/monsoon fails/pass closes forever}.
\item \textbf{Rail Tick:} Advance a relevant campaign/front clock by 1 segment.
\item \textbf{Condition Applied:} Mark \textbf{Fatigue 1/Harm 1/Condition} appropriate to fiction.
\end{itemize}

\subsection*{Major Turns (4+ SB)}
\begin{itemize}
\item \textbf{Trap Springs:} A prepared danger activates with full effect.
\item \textbf{Authority Arrival:} \textbf{Crown Council, the Astrologer, or the Pass Ghost} intervenes.
\item \textbf{Scene Shift:} The environment changes dramatically—\textbf{avalanche buries the pass/monsoon surf floods the wharf/coffee terraces catch fire}.
\item \textbf{Patron Omen:} Divine/arcane forces take notice—\textbf{omen appears/blessing lost/curse manifests}.
\item \textbf{Narrative Pivot:} The story takes an unexpected turn that reframes objectives.
\end{itemize}

\subsection*{Region-Specific SB Options}
\begin{itemize}
\item \textbf{Taharka (Highlands):} Pass stones glow faintly, horse bells ring without wind, snow falls upward.
\item \textbf{Taharka (Coast):} Monsoon winds carry whispers, spice dust forms faces, the trade creole sprouts new words unbidden.
\item \textbf{Taharka (Crown Politics):} Council doors slam, verdicts are spoken in rhyme, a ninth cup appears on every table.
\end{itemize}

\subsection*{Pursuit Ladder (Near / Mid / Far)}
\begin{itemize}
\item \textbf{Advance/Withdraw (Wits + Survival):} On a hit, shift one band. On a strong hit, also impose \emph{Thin Air}.
\item \textbf{Cut the Pass (Presence + Sway):} On a hit, freeze range for one exchange; if a pass-stone is in your possession, +1 die.
\item \textbf{Weather the Monsoon (Resolve + Sail):} On a hit in \emph{Monsoon Delay}, you pick which ships dock; on a miss, the GM does.
\end{itemize}

\subsection*{Boss Seeds (Horse \& Monsoon)}
\begin{itemize}
\item \textbf{The Horse-Line Thief (Nine Bloodlines)} — \emph{Phases}: Foal Thefts $\rightarrow$ Orchestrated Plague $\rightarrow$ Stable Coup. \emph{Cracks}: recover a stolen foal, expose a false lineage, survive a stampede.
\item \textbf{The Monsoon-Breaker (Storm That Stalls)} — \emph{Phases}: Delayed Ships $\rightarrow$ Procession of Empty Wharves $\rightarrow$ Famine Baptism of a Merchant. \emph{Cracks}: relic monsoon schedule, survivor testimony, the wind turning against them.
\end{itemize}

\newpage

% -------------------- SEKOGO --------------------
\chapter{Sekogo}
\emph{The Green Council.}

\noindent
Kofi shrugs: \textit{“The jungle is not a place. It is a council. The Ukwe – the Speaking Tree – decides whose debts are collected and whose are forgiven. I have seen a man lie before the Ukwe. He did not speak again for a year.”}

\noindent
Sekogo lies between the Sahel grasslands and the southern plateau – a dry savannah that grows wetter as you travel east, until the baobabs give way to deep jungle. The Green Council rules through consensus; honey‑orchards are sacred; old spirits still walk. Slavery is forbidden – the jungle itself would rise against any slaver who crossed its border. The Ukwe, the Speaking Tree, remembers every promise. The Leopard Society hunts those who break the deepest taboos. Honey is a bond. A mask is a law.

\vspace{0.5em}
\noindent\textit{“Never cut a marked tree. The bark grows over the names of debtors, but the wood remembers. Cut it, and the Ukwe will add your name to its ledger.”} \\
— Hunter Okoro, Leopard Society mask‑wearer

\vspace{1em}
\begin{almanac}
\textbf{What do people eat for breakfast?} Mashed yams and groundnut stew – the yams from clearings that shift with the seasons, the groundnuts harvested by monkey‑tribes who speak the old tongue. The monkeys are not paid. They are remembered.

\textbf{What does a child learn before they can walk?} To name the spirits. River, Tree, Stone, Rain, Fire, Wind, Honey, and the ninth – the spirit that has no name. The ninth spirit is the Ukwe’s shadow. It never forgives.

\textbf{What is the first thing you notice about a stranger?} Their mask. Carved wood, painted with lineage symbols. A cracked mask means a broken taboo. A blank mask means no lineage. No mask means not a person. The Leopard Society hunts the maskless.

\textbf{What is the most common crime?} Debt‑forgetting – a sickness. The Ukwe keeps the Calendar of Debt. A forgotten debt grows roots and strangles the debtor’s house.

\textbf{What does no one talk about but everyone knows?} The Green Council does not rule. The Ukwe rules. The Council is its voice – eight speakers chosen by the tree’s roots. The ninth seat is empty. The ninth speaker is the tree itself. It does not speak. It waits.

\textbf{Where do people go when they die?} To the Ukwe’s roots. Bodies laid at the base of the Speaking Tree. The roots absorb them. A new leaf grows for each soul. The leaves whisper the names of the dead.

\textbf{How do you apologize for a serious wrong?} You hammer a nail into a \textit{nkisi} power figure. The nail is your wrong. The spirit inside accepts it. When the nail rusts away, your debt is paid. The nail never rusts.

\textbf{What is the most important festival?} The Honey Harvest (autumn). The bees are sung to sleep. The first bowl is offered to the Ukwe. The tree drinks. The tree does not thank.

\textbf{What is the secret that could destroy a powerful person?} The Ukwe is not a tree. It is a prison – the binding of a spirit that tried to consume the jungle nine centuries ago. The spirit is still alive. The tree is its cage. The Green Council has forgotten what it guards.
\end{almanac}

\newpage
\section{Sekogo −−− ``The Green Council''}
łabel{chap:sekogo}

\textbf{Key Demographics: } Humans
\textbf{Capital City: } Ukwe (The City)

\section*{Region Diagnostic Template}

\subsection*{\textbf{\textsf{Sekogo}}}

\paragraph{Current Hook (one sentence)}
\framebox[łinewidth][l]{%
\begin{minipage}{łinewidth}
\vspace{2mm}
The Ukwe (Speaking Tree) demands a name: a caravan master broke his oath to the grove, and the roots are hunting him down.
\vspace{2mm}
\end{minipage}}

\paragraph{Core Conflict (one sentence)}
\framebox[łinewidth][l]{%
\begin{minipage}{łinewidth}
\vspace{2mm}
Players are caught between the Conservators (who demand stasis) and the Wakers (who demand growth through fever and blood), while the Leopard Society hunts anyone who breaks the balance.
\vspace{2mm}
\end{minipage}}

\paragraph{The One Thing Only This Region Does}
\framebox[łinewidth][l]{%
\begin{minipage}{łinewidth}
\vspace{2mm}
Promises are physical objects (\textit{mukanda} bark−cords) that tighten when strained and rot when released; breaking them attracts the \textit{Mwitu} (Hollow).
\vspace{2mm}
\end{minipage}}

\paragraph{Interactive Mechanics (player−facing)}
\begin{itemize}
ℹtem[\card{♠}{1}] \textbf{The Ukwe's Witness:} Speak a secret to a marked tree to gain \textbf{1 Boon}, but the tree remembers it forever (GM gains \textbf{1 SB} when that secret is relevant).
ℹtem[\card{♥}{2}] \textbf{Honey−Tithe:} Offer honey to gain safe passage; refusing marks \textbf{1 Obligation} to the Grove Spirit.
ℹtem[\card{♣}{3}] \textbf{Mask Law:} Wearing a ceremonial mask grants authority but binds you to its taboos; breaking them causes the mask to crack (\textbf{−−1 die} social).
ℹtem[\card{♦}{4}] \textbf{Ananse's Bargain:} Renegotiate a promise once per session (Spirit + Lore, DV 4); success shifts terms, failure doubles the burden.
\end{itemize}

\paragraph{Campaign Timers (GM−facing)}
\begin{itemize}
ℹtem \textbf{Waker Fever [6]} – \textit{Advances when:} Players ignore spirit warnings or use Waker spores. \textit{At full:} The jungle becomes hostile; navigation rolls at \textbf{−−1 die}.
ℹtem \textbf{Ukwe's Memory [8]} – \textit{Advances when:} Players break promises or cut marked trees. \textit{At full:} The \textit{Mwitu} manifests; the party is lost until a debt is paid.
\end{itemize}

\paragraph{Faction Triad}
\begin{tabular}{|p{0.2łinewidth}|p{0.25łinewidth}|p{0.25łinewidth}|p{0.2łinewidth}|}
ℎline
\textbf{Faction} & \textbf{Goal} & \textbf{Method} & \textbf{Why players care} \\
ℎline
Leopard Society & Restore Balance & Hunt oath−breakers and slavers & They offer protection but demand justice. \\
ℎline
Earth Society & Preserve Memory & Keep the Calendar of Witness & They hold the records of past debts. \\
ℎline
Wakers & Awaken the Jungle & Use fever−mists and blood rites & They offer power but risk corruption. \\
ℎline
\end{tabular}

\paragraph{The Price}
\framebox[łinewidth][l]{%
\begin{minipage}{łinewidth}
\vspace{2mm}
Power requires memory loss, Obligation to spirits, or physical marks (scars, thorns); honey is currency, but the bees remember who took it.
\vspace{2mm}
\end{minipage}}

\paragraph{Geography as Constraint}
\framebox[łinewidth][l]{%
\begin{minipage}{łinewidth}
\vspace{2mm}
The jungle is a witness; marked trees cannot be cut, and paths shift if promises are broken. The Crimson Valley bridges only hold if web−law is honored.
\vspace{2mm}
\end{minipage}}

\paragraph{The Ninth Taboo manifestation}
\framebox[łinewidth][l]{%
\begin{minipage}{łinewidth}
\vspace{2mm}
The Ninth Root searches for erased names; the ninth knot on a \textit{mukanda} is left loose for the Hollow; the ninth bell of Ananse's Child rings when a promise is broken.
\vspace{2mm}
\end{minipage}}

\paragraph{Current Weaknesses (honest self−assessment)}
\begin{itemize}
ℹtem Too many spirit types (\textit{Mzuka}, \textit{Kizimba}, \textit{Ngozi}, etc.) may confuse players.
ℹtem Faction goals overlap too much (all care about promises).
ℹtem Risk of "magic everywhere" diluting the cost of spellcasting.
\end{itemize}

\vspace{1cm}
ℎrule
\vspace{0.5cm}
\noindent\emph{Design Fixes Applied: Consolidated spirits into three categories (Ancestral, Nature, Vengeant). Clarified faction distinctions (Stasis vs. Growth vs. Enforcement). Integrated core \textbf{Fate's Edge} mechanics (Boons, SB, Obligation) into all rewards.}

\subsection*{Quick Reference}
\textbf{Tagline:} \textit{The jungle remembers every promise. The mask is the law. The honey is a bond.}

\begin{quote}
  \textbf{Leopard Society Hunter (Elite):}  
  ``The jungle has a memory. When a fever rises, we cut it out. When an oath is broken, we restore the balance. We are not cruel. We are necessary. The tree remembers every deed, and we are its teeth.''
\end{quote}

\begin{quote}
  \textbf{Honey−Gatherer (Commoner):}  
  ``The bees know my name. They do not sting me because I sing to them the song of the first blossom. My grandmother taught me that the honey is not a gift; it is a trust. One day, the grove will need something from you, and you will remember the taste.''
\end{quote}

\textbf{Genre / Mood:} Jungle folk horror, spirit diplomacy, mask−law.

\textbf{Starting Location:} A honey−orchard at the jungle's edge, where the trees hum and the air is thick with bees. A masked figure in a leopard pelt steps from the shadows: ``The Ukwe has asked for a name. A caravan master broke his oath to the grove. Find him before the roots do −−− or become his witness.''

\vspace{2mm}
\begin{tcolorbox}[colback=black!3,colframe=blue!60!black,title={Theme & Atmosphere}]
Sekogo lies between the Sahel grasslands and the southern plateau −−− a dry savannah that grows wetter as you travel east, until the baobabs give way to deep jungle and the air turns thick with the smell of damp earth and overripe fruit. Here, the Green Council rules through consensus, honey−orchards are sacred, and the old spirits still walk. Slavery is forbidden; the jungle itself would rise against any slaver who crossed its border. But the Council is divided: Conservators want to protect the old pacts, while Wakers argue that the jungle has grown sluggish and needs to be roused with fever−mists, thorn−spirits, and older, bloodier rites. At the eastern edge lies the Crimson Valley, home to the Lethai−ar −−− chestnut−brown elves whose silk bridges span red canyons and whose web−law keeps both empires from using the valley as a shortcut to war.

\vspace{2mm}
\textbf{Neighbors and the Unclaimed Ground:}
\begin{itemize}
  ℹtem \textit{North −− Ameria and the Neutral Kingdoms:} Beyond Sekogo's northern savannah lies a patchwork of small, independent kingdoms −−− buffer states that pay tribute to neither Sekogo nor Kahfagia. Their courts are weak, their laws a creole of jungle custom and coastal tariff. Smugglers and exiles thrive there. The Leopard Society occasionally hunts across the border, but always with a local guide.
  ℹtem \textit{East −− The Crimson Valley (Lethai−ar):} The Ukwe's roots extend into the valley's red stone. The Lethai−ar silk−wardens respect the Speaking Tree, and the Green Council's witnessed promises are sometimes recorded on a Lethai−ar cord. A bridge toll paid in Sekogo honey is honored on both sides of the gorge.
  ℹtem \textit{South and West −− The Jungle's Own:} Sekogo's borders are not lines but gradients. The deeper you go, the more the Ukwe's roots thicken and the Leopard Society's law becomes the only law.
\end{itemize}

\vspace{2mm}
\textbf{What you notice first:} The sound of bees −−− a constant, low drone that never stops, even at night. The way the trees lean toward you when you speak a secret. The smell of honey and rotting fruit, sweet and sharp at once.

\textbf{The one rule that kills newcomers:} Never cut a marked tree. Each such tree is a living witness. The bark grows over the names of those who have broken harmony, but the wood remembers. Cut it, and the Ukwe will add your name to its memory.
\end{tcolorbox}

\subsection*{Regional Mood: The Weight of the Root}

Power here is measured in witnessed promises, honey tithes, and the balance between jungle and council. The Leopard Society hunts those who break the deepest taboos −−− oath−breakers, slavers, outsiders who take more than they give. The Earth Society keeps the \textbf{Calendar of Witness}, a sacred list written on bark and buried at the Ukwe's roots, preserving memories that must not be forgotten.

The spirits of Sekogo are consolidated into three kinds for play: \textit{Ancestral} (bound to groves, speak in dead voices), \textit{Nature} (guardians like crocodiles or baobabs, show rather than speak), and \textit{Vengeant} (born from blood−promises, hunt until answered). The Hollow manifests here as the \textit{Mwitu}, the forest that forgets where the village ends −−− and that forgetfulness is a hunger.

\subsection*{Prominent NPCs of Sekogo}
łabel{sec:sekogo−npcs}

\begin{longtable}{|p{0.2\textwidth}|p{0.25\textwidth}|p{0.3\textwidth}|p{0.15\textwidth}|}
ℎline
\textbf{Name} & \textbf{Role/Title} & \textbf{Motivation} & \textbf{Rumored Quirk} \\
ℎline
Hunter Okoro & Leopard Society mask−wearer & Hunt down a Waker who poisoned a foreign caravan & He never removes his mask. The mask is his real face now. \\
ℎline
Elder Imani & Earth Society Ogboni & Keep the Calendar of Witness from being forgotten & She speaks only while holding a piece of the Ukwe's bark. She cannot lie. \\
ℎline
Nganga Malika & Spirit−doctor of the \textit{minkisi} & Heal a wound caused by a broken \textit{nkisi} & She carries a carved figure with nine nails. The ninth nail is loose. \\
ℎline
Waker Sefu & Radical who wants to wake the jungle & Poison the wells of a Conservator village & A green shoot grows from his shoulder. It bleeds when he lies. \\
ℎline
Lethai−ar Silk−Warden Anara & Guard of the Crimson Valley bridges & Keep the web−law intact against Oshiiran pressure & Her skin−tattoos record a witnessed promise her grandmother carried. She will not say to whom. \\
ℎline
The Ukwe's Voice & A spirit that speaks through a cracked root & Remind the Council of a promise nine years overdue & It speaks in the voice of a dead child. The child was a witness. \\
ℎline
Ananse's Child (Trickster) & A spider−spirit in human form & Unravel oaths for the joy of chaos & It wears nine tiny bells on its ankles; each bell is a broken promise. \\
ℎline
Mami Wata & Water spirit of the Crimson Falls & Collect a memory from every soul that crosses her pool & Her hair is black water; her eyes are river stones. She never blinks. \\
ℎline
\end{longtable}

\subsection*{Names of Sekogo}
łabel{sec:sekogo−names}

Inspired by the savannah and jungle regions: melodic, with nasal consonants and tonal shifts. Surnames often denote a totem animal, a grove, or a witnessed deed.

\begin{longtable}{|p{0.25\textwidth}|p{0.25\textwidth}|p{0.2\textwidth}|p{0.2\textwidth}|}
ℎline
\textbf{First Names (Male)} & \textbf{First Names (Female)} & \textbf{Surnames} & \textbf{Epithets} \\
ℎline
Okoro & Imani & Ngoma & \textit{the Masked} \\
ℎline
Sefu & Malika & Ukwe−born & \textit{the Honey−Tongue} \\
ℎline
Chima & Anara & Leopard−kin & \textit{the Oath−Keeper} \\
ℎline
Kofi & Zola & Root−walker & \textit{the Ninth Nail} \\
ℎline
Tunde & Amara & Grove−singer & \textit{the Unmarked} \\
ℎline
Biko & Nkiru & Minkisi−hand & \textit{the Still Witness} \\
ℎline
Adofo & Adaeze & Ananse−fo & \textit{the Trickster's Child} \\
ℎline
\end{longtable}

\paragraph*{(Orchard/Circle/Bridge)} Honey−orchard where the trees hum; druid circle with moss−covered stones and a leopard skull at center; Lethai−ar silk bridge spanning a crimson gorge with the Ukwe's roots clutching the water below.

\section*{Spades −−− Places}
łabel{sec:sekogo−places}
\begin{enumerate}
\setcounter{enumi}{1}
ℹtem \textbf{Honey−Orchard} −−− The trees hum. The air is thick with bees.
  ℎfill \textit{The honey is dark red. It tastes of iron and forgotten promises. The ninth hive is always empty −−− the bees say the Hollow has taken the queen.} \texttt{[SPIRIT]}
ℹtem \textbf{Druid Circle (Leopard Skull)} −−− A ring of moss−covered stones. A skull sits at the center.
  ℎfill \textit{The skull is warm. When the Leopard Society hunts, the eye sockets glow −−− and the \textit{mzuka} of the skull whispers the quarry's true name.} \texttt{[RITUAL]}
ℹtem \textbf{Lethai−ar Silk Bridge} −−− A span of woven silk across a crimson gorge.
  ℎfill \textit{The silk is made from spider−sap and oath−hair. It will not break unless an oath is called. The ninth strand is invisible −−− it is the Hollow's witness.} \texttt{[OATH]}
ℹtem \textbf{The Ukwe (Speaking Tree)} −−− A baobab so old its bark is carved with a thousand names.
  ℎfill \textit{The tree speaks in the voices of the dead. It never gives answers −−− only new questions and old memories. The ninth root never stops growing; it searches for the name that was erased.} \texttt{[WITNESS]}
ℹtem \textbf{Nkisi Grove} −−− A clearing where broken power−objects are buried.
  ℎfill \textit{The ground hums. The nails that sealed the \textit{minkisi} have grown into iron thorns. The ninth thorn is always warm; it marks the grave of the nganga who drove the first nail.} \texttt{[MAGIC]}
ℹtem \textbf{Calendar of Witness (Earth Society Vault)} −−− A hollow log sealed with wax and bone.
  ℎfill \textit{Inside are strips of bark. Each strip records a broken harmony. The bark is still growing −−− and the ninth strip has no writing. It waits for a promise that has not yet been broken.} \texttt{[MEMORY]}
ℹtem \textbf{Waker's Fever−Mist Pool} −−− A stagnant pond that steams even in winter.
  ℎfill \textit{The mist smells of rotting flowers. Breathe it, and you will see the jungle as the Wakers see it −−− a place that is starving, waiting to be woken. The ninth breath of the mist shows the \textit{Mwitu}.} \texttt{[WEATHER]}
ℹtem \textbf{Leopard Society Lodge} −−− A longhouse hidden in a thicket of thorn−vines.
  ℎfill \textit{The walls are made of human jawbones. The jaws are from oath−breakers. They still click when a new name is spoken −−− and the ninth click is always the name of the next quarry.} \texttt{[SECRET]}
ℹtem \textbf{Crimson Valley Gorge} −−− A canyon of red stone, crossed by silk bridges at three points.
  ℎfill \textit{The stone is warm. It bleeds rust when the web−law is broken. The ninth bridge is a myth −−− it appears only when a promise is witnessed by a ghost.} \texttt{[THRESHOLD]}
ℹtem[J] \textbf{Conservator Archive} −−− A cave filled with sealed clay pots. Each pot contains a memory.
  ℎfill \textit{The memories are for sharing. The price is a memory of equal weight. The ninth pot is cracked −−− the memory inside is leaking, and no one remembers what it was.} \texttt{[MEMORY]}
ℹtem[Q] \textbf{The Grove of Sleeping Spirits} −−− A forest where the \textit{minkisi} go to die.
  ℎfill \textit{The spirits do not wake. They dream of the first nail. The dream is contagious −−− and the ninth dream is the Hollow's own.} \texttt{[DEATH]}
ℹtem[K] \textbf{The Ukwe's Root−Chamber} −−− A cavern beneath the tree where the roots form a throne.
  ℎfill \textit{The throne is made of braided roots. Sit on it, and the tree will ask one question. Answer poorly, and you never leave. Answer with a lie, and the ninth root will bind your tongue.} \texttt{[JUDGMENT]}
ℹtem[A] \textbf{The Crimson Falls} −−− A waterfall that runs red after the first rain of spring.
  ℎfill \textit{The water is not blood. It is rust from the canyon walls. The Lethai−ar say it is the earth's memory of a war −−− and the ninth drop is always cold, because it fell on a broken oath.} \texttt{[SPIRIT]}
ℹtem[A+] \textbf{The Pool of Mami Wata} −−− A still, black pool behind the Crimson Falls, reachable only by diving underwater.
  ℎfill \textit{The water is warm. It shows not your reflection, but the face of the person whose memory the spirit holds. The ninth ripple is always your own face −−− but older, and weeping.} \texttt{[WATER]}
\end{enumerate}

\paragraph*{(Gatherer/Druid/Warden)} Honey−gatherer with a fruit offering and a leopard claw necklace; Green Council druid with a thorn−staff and milky eyes that see the weight of promises; Lethai−ar silk−warden, masked and quiet, his mask painted with the Ukwe's roots.

\section*{Hearts −−− People & Factions}
łabel{sec:sekogo−people}
\begin{enumerate}
\setcounter{enumi}{1}
ℹtem \textbf{Honey−Gatherer} −−− She wears a leopard claw around her neck. The claw is from a cat she killed with her bare hands.
  ℎfill \textit{She offers a fruit to the trees before each harvest. The fruit is always a little rotten −−− the ninth bite is for the \textit{Mwitu}.} \texttt{[COMMUNITY]}
ℹtem \textbf{Green Council Druid (Conservator)} −−− Milk−white eyes, a staff of living wood.
  ℎfill \textit{She cannot see the living. She sees only the threads of promises −−− glowing connections between every person and the Ukwe. The ninth thread is always frayed.} \texttt{[JUSTICE]}
ℹtem \textbf{Waker Radical} −−− Carries a bag of fever−mist spores. His eyes are yellow.
  ℎfill \textit{He plants thorns in the paths of Conservator caravans. The thorns grow into walls overnight −−− and the ninth thorn is always a \textit{nkisi} nail, driven into a forgotten promise.} \texttt{[DANGER]}
ℹtem \textbf{Lethai−ar Silk−Warden} −−− Mask carved from the Ukwe's root, painted with oath−names.
  ℎfill \textit{He never speaks above a whisper. The bridge hears him anyway −−− and the ninth whisper is always the name of someone who will break a promise before dawn.} \texttt{[OATH]}
ℹtem \textbf{Leopard Society Hunter} −−− Wears a pelt that shifts and moves as if alive.
  ℎfill \textit{He can transform into a leopard at night. Or so he claims. No one has survived a hunt to confirm it −−− but the ninth transformation is said to be permanent.} \texttt{[FEAR]}
ℹtem \textbf{Earth Society Elder} −−− Holds a piece of the Ukwe's bark. The bark is warm.
  ℎfill \textit{She asks every visitor: ``What have you forgotten to honor?'' The answer is carved into her memory −−− and the ninth answer is always the same: a name that cannot be spoken.} \texttt{[TRUTH]}
ℹtem \textbf{Nganga (Spirit−Doctor)} −−− Carries a \textit{nkisi} carved with nine nails. The ninth nail is loose.
  ℎfill \textit{She hammers a nail into the figure to call a spirit. Each nail is a cost. She has nine costs left −−− and the ninth nail, if driven, will call the Hollow.} \texttt{[MAGIC]}
ℹtem \textbf{Oath−Keeper (Hollow Manifestation)} −−− A figure made of bark and silence. It appears when the Calendar of Witness grows heavy.
  ℎfill \textit{It does not speak. It points at the one who broke harmony. That person's name appears on the nearest tree −−− and the ninth letter of the name glows green.} \texttt{[DEATH]}
ℹtem \textbf{Lethai−ar Mark−Cutter} −−− A tattooist who records promises on skin.
  ℎfill \textit{She uses a needle of silk−fiber and ink made from crushed crimson stone. The ink glows when a promise is broken −−− and the ninth mark she ever cut was her own.} \texttt{[MEMORY]}
ℹtem[J] \textbf{The Ukwe's Voice (Spirit)} −−− A crack in the tree's bark that speaks.
  ℎfill \textit{The voice changes. Sometimes it is a child. Sometimes it is a jarl. Sometimes it is the listener's own. The ninth word it speaks is always the name of a promise you forgot to make.} \texttt{[SPIRIT]}
ℹtem[Q] \textbf{The First Nail (Legend)} −−− A mythic nganga who sealed a spirit into the Ukwe itself.
  ℎfill \textit{She is said to still walk the jungle, carrying a hammer made of root and bone. The ninth nail she drove was never found −−− it is still searching for the promise it was meant to seal.} \texttt{[LEGEND]}
ℹtem[K] \textbf{The Leopard King} −−− A spirit that leads the Society's midnight hunts.
  ℎfill \textit{It has no mask. Its face is a leopard's. Its eyes are human. The ninth hunter who follows it never returns −−− they become the Leopard King's shadow.} \texttt{[FEAR]}
ℹtem[A] \textbf{The Crimson Valley's First Weaver} −−− A ghost who teaches the web−law to new wardens.
  ℎfill \textit{She appears at dawn on the central bridge. She asks one question: ``Who watches the watchers?'' The ninth warden who answers wrongly is woven into the bridge.} \texttt{[JUSTICE]}
ℹtem[A+] \textbf{Ananse's Child (Trickster)} −−− A small figure with too−long fingers, wearing nine tiny bells.
  ℎfill \textit{It offers to ``renegotiate'' any promise −−− but its renegotiation always adds a new, hidden clause. The clause is a joke. The joke is always on you −−− and the ninth bell rings only when the joke is already done.} \texttt{[TRICKERY]}
\end{enumerate}

\paragraph*{(Spirit/Offer/Name)} A thorn−spirit blocks the path −−− must offer honey or fight; the Ukwe demands a name before letting you pass; a Lethai−ar silk−walker challenges you to recite the web−law.

\section*{Clubs −−− Complications/Threats}
łabel{sec:sekogo−complications}
\begin{enumerate}
\setcounter{enumi}{1}
ℹtem \textbf{Thorn−Spirit Block} −−− A wall of living thorns grows across the path.
  ℎfill \textit{The spirit −−− a \textit{kizimba} of the hedge −−− demands a name. Give a false one, and the thorns grow inward. Give a true one, and the thorns will part −−− but the spirit will remember your face.} \texttt{[SPIRIT]} \texttt{[OBLIGATION]}
ℹtem \textbf{Ukwe Demands a Name} −−− The tree speaks: ``I have forgotten a promise. Remind me.''
  ℎfill \textit{If you cannot name a broken promise, the tree will take your name instead −−− and the ninth root will grow a new leaf with your face on it.} \texttt{[WITNESS]} \texttt{[TRUTH]}
ℹtem \textbf{Lethai−ar Web−Law Challenge} −−− A silk−warden blocks the bridge: ``Recite the ninth clause.''
  ℎfill \textit{The ninth clause is unwritten. It changes with each warden. You must guess −−− but the ninth guess is always wrong, because the clause is the warden's own secret.} \texttt{[OATH]}
ℹtem \textbf{Leopard Society Hunt} −−− A masked hunter appears at dusk. ``You are called to witness.''
  ℎfill \textit{Refusal earns a Leopard's Mark. The jungle becomes hostile to you −−− and the ninth beast you meet will be wearing a mask of your face.} \texttt{[FEAR]}
ℹtem \textbf{Nkisi Hunger} −−− A spirit−object you carry begins whispering. It wants a memory.
  ℎfill \textit{Refuse, and the \textit{nkisi} shatters. The spirit inside is not friendly −−− and the ninth shard will embed itself in your shadow, following you until you feed it.} \texttt{[MAGIC]}
ℹtem \textbf{Waker Poisoning} −−− A well is found with a green film on the water.
  ℎfill \textit{The poison is not lethal. It makes you dream of the jungle's hunger −−− and the ninth dream is always the \textit{Mwitu} asking, ``Why are you still here?''} \texttt{[WEATHER]}
ℹtem \textbf{Earth Society Reckoning} −−− An elder demands to see what you carry in your memory.
  ℎfill \textit{The record is written in scars on your skin. You have no scars. That is suspicious −−− so the elder will give you nine scars, one for each forgotten promise you cannot name.} \texttt{[TRUTH]}
ℹtem \textbf{Honey−Tithe Demanded} −−− A grove spirit appears: ``Pay the ninth jar.''
  ℎfill \textit{You have no honey. The spirit offers a bargain. The price is a story −−− and the ninth story you tell will become a new root in the Ukwe's memory.} \texttt{[TRADE]} \texttt{[OBLIGATION]}
ℹtem \textbf{Crimson Valley Border Dispute} −−− Oshiiran surveyors cross the web−line. The Lethai−ar arm their bows.
  ℎfill \textit{The party is caught between. Whoever you side with, the other will remember −−− and the ninth bridge will close to you forever.} \texttt{[POLITICS]}
ℹtem[J] \textbf{The Ukwe's Root Breaches} −−− A root bursts through a village floor. It is looking for a specific name.
  ℎfill \textit{The name is carved on the root. The root is weeping. The one who broke the promise is still alive −−− and the ninth night they sleep, the root will find them.} \texttt{[DEATH]}
ℹtem[Q] \textbf{Leopard Society Coup} −−− The Wakers have infiltrated the Lodge. The masks are fighting.
  ℎfill \textit{You cannot tell who is who. Both sides wear leopard pelts. Both sides smell of blood −−− and the ninth mask you see will be the one that speaks with the Leopard King's voice.} \texttt{[CONFLICT]}
ℹtem[K] \textbf{Calendar Fire} −−− A spark catches the Earth Society vault. The memories are burning.
  ℎfill \textit{The smoke forms faces. The faces are those who broke harmony. They are screaming −−− and the ninth scream is your own, from a promise you made and then forgot.} \texttt{[DISASTER]}
ℹtem[A] \textbf{The Crimson Falls Run Red} −−− The waterfall turns to blood. The Lethai−ar say the web−law has been broken.
  ℎfill \textit{The break was centuries ago. The falls are just catching up −−− and the ninth drop of blood carries the name of the first oath−breaker.} \texttt{[OMEN]}
ℹtem[A+] \textbf{Mami Wata's Gaze} −−− A figure rises from the black pool behind the falls. It asks: ``What promise have you never been able to keep?''
  ℎfill \textit{If you lie, the water will take your voice. If you tell the truth, it will ask for a memory as price −−− and the ninth memory you offer will be the one you need most.} \texttt{[SPIRIT]} \texttt{[OBLIGATION]}
\end{enumerate}

\paragraph*{(Favor/Trust/Shard)} Leopard's Favor (once, call on the Leopard Society for a hunt −−− you are bound to them for a future hunt); Ukwe Memory (reroll a failed Survival test once per session, but GM gains 1 SB for a jungle complication); Nkisi Shard (add $+1$ Effect to a magical working, but the shard shatters and costs 1 Corruption).

\section*{Diamonds −−− Rewards/Leverage}
łabel{sec:sekogo−rewards}
\begin{enumerate}
\setcounter{enumi}{1}
ℹtem \textbf{Leopard's Favor} −−− A leopard claw wrapped in red thread. One use to summon a Society hunt.
  ℎfill \textit{The favor is a bond. The Society will remember. The price is a future hunt −−− and the ninth hunt will be for your own name, if you break the bond.} \texttt{[BOND]}
ℹtem \textbf{Ukwe Memory (Leaf)} −−− A single leaf from the Speaking Tree. Reroll a failed Survival test once per session.
  ℎfill \textit{The leaf is green. It will not wilt. The tree remembers who took it −−− and the ninth time you use it, the leaf will turn black and whisper a name you have forgotten.} \texttt{[MEMORY]}
ℹtem \textbf{Nkisi Shard} −−− A piece of a broken power−object. Add $+1$ Effect to a magical working, then it shatters.
  ℎfill \textit{The shard is warm. It whispers secrets in a language you almost understand −−− and the ninth secret is always about a promise you made to yourself and broke.} \texttt{[MAGIC]}
ℹtem \textbf{Honey−Jar (Royal Tithe)} −−− A clay pot of dark red honey. Trade it for safe passage through any grove.
  ℎfill \textit{The honey is thick. It moves in the jar as if alive. Do not open it −−− the ninth time you open it, the bees will remember your name and ask for a tithe.} \texttt{[TRADE]}
ℹtem \textbf{Earth Society Tablet} −−− A strip of bark with a broken promise recorded. Present it to recall a forgotten harmony.
  ℎfill \textit{The bark is still growing. The memory is still fresh −−− and the ninth time you use it, the bark will grow over your hand and show you the promise you yourself broke.} \texttt{[MEMORY]}
ℹtem \textbf{Lethai−ar Silk Cord} −−− A braided cord that records a witnessed promise. Use it to enforce any oath.
  ℎfill \textit{The cord tightens when the oath is broken. It will not stop until blood is drawn −−− and the ninth drop of blood will be your own, if you witnessed a false promise.} \texttt{[OATH]}
ℹtem \textbf{Leopard Mask (Temporary)} −−− Wear it to pass as a Society hunter for one scene.
  ℎfill \textit{The mask has eye holes that see the weight of promises. It also sees your own −−− and the ninth time you wear it, the mask will not come off.} \texttt{[SECRET]}
ℹtem \textbf{Waker's Spore Pouch} −−− A bag of fever−mist spores. Release them to confuse pursuers.
  ℎfill \textit{The spores cause vivid hallucinations. The hallucinations are always of the Ukwe −−− and the ninth hallucination will be a question: ``Why did you lie?''} \texttt{[STEALTH]}
ℹtem \textbf{Nganga's Hammer} −−− A small iron hammer. Drive a nail into any object to seal a spirit inside.
  ℎfill \textit{The hammer has a cost. Each nail requires a memory −−− and the ninth nail you drive will seal a fragment of your own soul into the object.} \texttt{[MAGIC]}
ℹtem[J] \textbf{The Ukwe's Answer} −−− A single truth from the Speaking Tree, whispered on the wind.
  ℎfill \textit{The truth is always about a broken promise. The promise may not be yours −−− but the ninth time you listen, the tree will speak a promise you made and forgot.} \texttt{[TRUTH]}
ℹtem[Q] \textbf{Silk−Warden's Mask} −−− A carved mask that grants safe passage across any Lethai−ar bridge.
  ℎfill \textit{The mask is also a witness. It will remember every promise you make while wearing it −−− and the ninth promise will be witnessed by the Hollow.} \texttt{[OATH]}
ℹtem[K] \textbf{Oath−Keeper's Finger} −−− A dried root shaped like a digit. Use it to call an Oath−Keeper once.
  ℎfill \textit{The Oath−Keeper will settle one broken promise for you. Then it will ask for yours −−− and the ninth finger of its hand will point at the promise you most fear to keep.} \texttt{[DEATH]}
ℹtem[A] \textbf{The First Nail} −−− A legendary iron nail that sealed the first spirit into the Ukwe.
  ℎfill \textit{Drive it into any tree to create a new Speaking Tree. The tree will remember everything −−− and the ninth name it speaks will be your own, unless you offer a memory in its place.} \texttt{[LEGEND]}
ℹtem[A+] \textbf{Mami Wata's Tear} −−− A drop of water that never dries, said to be the spirit's grief for a promise broken long ago.
  ℎfill \textit{Drink it to forget one broken promise permanently −−− but you will dream of drowning once each season. The ninth dream will be the promise itself, remembered.} \texttt{[WATER]} \texttt{[MEMORY]}
\end{enumerate}

\section*{Quick use notes}
łabel{sec:sekogo−quick−use}
\begin{itemize}
ℹtem Draw until you have all four suits: Spade = place, Heart = actor, Club = pressure, Diamond = leverage. Highest rank sets the main Timer (2−−5 $→$ 4, 6−−10 $→$ 6, J/Q/K $→$ 8, A $→$ 10).
ℹtem Diamonds are codified outcomes (favors, memories, shards) that change position rather than call for a roll.
ℹtem If any A appears, echo \textbf{grove−omens} −−− bees that fly backward, roots that bleed sap, a mask that watches from the trees, the ninth bell of Ananse's Child ringing in still air. If an A+ appears (special high−rank result), a spirit of the deep jungle (Mami Wata, Ananse's Child, the Leopard King) manifests directly.
\end{itemize}

\subsection*{Additional Features (Cultural Mechanics)}
łabel{sec:sekogo−mechanics}

\begin{itemize}
ℹtem \textbf{The Three Societies:} Sekogo is governed by three overlapping authorities: the Leopard Society (hunters who restore balance), the Earth Society (elders who keep the Calendar of Witness), and the nganga (spirit−doctors who work with \textit{minkisi}). A single elder often belongs to two societies, and the tension between their roles is where drama lives.
ℹtem \textbf{The Ukwe (Speaking Tree):} A living witness whose roots extend across Sekogo. Every broken promise, every forgotten harmony stains its leaves. It answers questions −−− but never with kindness, and never for free. Accepting its gift creates a reciprocal bond that the tree may call upon at any moment. The ninth root is never still −−− it searches for the name that was erased.
ℹtem \textbf{The Crimson Valley and Web−Law:} The Lethai−ar guard the valley with fierce neutrality. Their skin−tattoos (Marks) record promises witnessed. When a Lethai−ar dies, their Marks are cut from the skin and buried at the Ukwe's roots. The web−law is unwritten; each warden carries a different clause in memory. The ninth clause is always a secret known only to the First Weaver's ghost.
ℹtem \textbf{Mask Law:} In Sekogo, a mask is not a disguise −−− it is a role. Wearing the mask of a Leopard Society hunter, a Waker, or a spirit entitles you to speak for that office, but also binds you to its taboos. Breaking a mask's taboo (e.g., a Leopard mask showing mercy to a slaver) causes the mask to crack and the wearer to lose $1$ die on all social rolls until they atone. The ninth mask is the Hollow's own −−− it has no face, and it never cracks because it was never whole.
ℹtem \textbf{Bark Promise−Cords (\textit{Mukanda}):} Promises are recorded on strips of bark called \textit{mukanda}. Each strip is knotted; the knot tightens when the promise is strained, loosens when it is kept. A \textit{mukanda} that rots means the bond is released; one that hardens into wood means the obligation has set. The ninth knot on a \textit{mukanda} is never tied −−− it is left loose for the Hollow to judge.
ℹtem \textbf{Ananse's Bargain (Trickster Mechanic):} Once per session, a character may call upon Ananse's Child to ``renegotiate'' any promise, oath, or contract. Roll Spirit + Lore (DV 4). On a success, the terms shift in your favor −−− but the GM gains 1 SB and may add a hidden clause that will trigger at the worst possible moment. On a miss, the burden doubles −−− and the ninth clause added will be the Hollow's own terms.
ℹtem \textbf{The \textit{Mwitu} (Hollow of the Forest):} When the Ninth manifests in Sekogo, it is as the \textit{Mwitu} −−− the forest that forgets where the village ends. It appears as a wall of gray mist between the trees, and it whispers: ``Why are you still here?'' Those who cannot answer lose their sense of direction (all navigation rolls at $−1$ die for one scene). The \textit{Mwitu} can be appeased by leaving a \textit{mukanda} at the edge of the mist −−− a promise to remember a forgotten name.
ℹtem \textbf{Local Patrons (Syncretic Names):}
  \begin{itemize}
    ℹtem \textbf{The Root Father} −−− Grimmir. \textit{The ancient pact between the jungle and the first honey−gatherer. His ninth root is the Hollow's own.}
    ℹtem \textbf{The Masked One} −−− Mykkiel. \textit{The law that wears a leopard's face; the judge who hunts. His ninth mask is a blank −−− for the case he cannot judge.}
    ℹtem \textbf{The Spider's Whisper} −−− Inaea. \textit{The trickster who spins webs in the spaces between promises. Her ninth knot is never tied.}
    ℹtem \textbf{The Water Mother} −−− Raéyn. \textit{Mami Wata, who holds the memories of the drowned. Her ninth tear is the one she never sheds.}
  \end{itemize}
\end{itemize}

\subsection*{Lore Echoes (Cross−Expansion Hooks)}
\begin{itemize}
  ℹtem \textbf{The Ukwe's Shadow and the Hollow (Eastern Realms):} The Hollow is said to be the shadow of the Ukwe −−− the space where broken promises go when the tree forgets them. Aelaerem hedge−keepers leave offerings at the Ukwe's roots to keep the Hollow from wandering. The ninth offering is always a \textit{mukanda} with a single knot −−− a promise to remember.
  ℹtem \textbf{Lethai−ar Silk and the Way of Silk:} The silk bridges of the Crimson Valley are made from spider−sap and oath−hair. Tulkani traders prize the silk for memory−cords; Fhara water−clerks use it to seal well−contracts. A bolt of Crimson silk can buy safe passage across half the continent −−− but the ninth bolt is always reserved for the Hollow's bridge.
  ℹtem \textbf{Leopard Society and Ashaani Slavers:} The Leopard Society has a standing grudge against Ashaani slavers. A masked hunter may offer the party a deal: help free a group of captives from a Black Hand caravan, and the Society will forget an old broken promise. The ninth captive freed will be marked with a Leopard's claw −−− a sign that they are bound to the Society for a hunt.
  ℹtem \textbf{Nganga and the Kalrantha Amulet (Dhahara):} The nganga of Sekogo believe the Kalrantha amulet is a rogue \textit{nkisi} −−− a spirit−object that was never properly sealed. They want to drive a nail into it to bind the demon. The ninth nail would seal the amulet forever −−− but the nganga who drives it will become the amulet's new keeper.
  ℹtem \textbf{Ananse's Child and the Grey Wanderer's Grimoire:} The trickster spirit is known to trade fragments of forgotten prophecies for stories. Some psions claim that Ananse's Child taught the first Echo−Readers the art of reading broken promises in the wind. The ninth prophecy it ever told was about the Hollow's return −−− but it told it as a joke, and no one laughed.
  ℹtem \textbf{The Neutral Kingdoms (North):} Smugglers from the unclaimed lands often carry Sekogo honey across the savannah. A Leopard Society hunt that crosses the border may find itself negotiating with a petty king who broke a promise to the Ukwe three generations ago. The king's ninth ancestor was a witness −−− and his ghost still sits on the throne, waiting for the promise to be kept.
\end{itemize}

\subsection*{Sekogo Superstitions (burn for $+1$ die once)}
\begin{itemize}
  ℹtem \textbf{Bread to the Root:} Crumble a crust at the Ukwe's base. Ignore the first \emph{Ukwe Demands a Name} penalty. \texttt{[SUPERSTITION]}
  ℹtem \textbf{Knot the Silk:} Tie a single knot in a Lethai−ar bridge cord. Cancel \emph{Web−Law Challenge} or take $−1$ die on your next negotiation (your choice). \texttt{[SUPERSTITION]}
  ℹtem \textbf{Name to the Nail:} Whisper a dead nganga's name into a \textit{nkisi}. Gain $+1$ Effect when using a spirit−object this scene, but the ancestors will expect an answer −−− and the ninth answer must be a promise. \texttt{[SUPERSTITION]}
  ℹtem \textbf{Spit to the Mask:} Before donning a ceremonial mask, spit on the inside. If the mask accepts, you gain $+1$ die to all rolls while wearing it. If it rejects, the mask will not come off for $1\text{d}4$ hours −−− and the ninth hour, the mask will speak. \texttt{[SUPERSTITION]}
\end{itemize}

\subsection*{Thematic SB Spend Table}
łabel{sec:sekogo−sb}

\paragraph{Minor Complications (1 SB)}
\begin{itemize}
ℹtem \textbf{Exposure:} Your actions draw unwanted attention from \textbf{Leopard Society hunters or Earth Society elders}.
ℹtem \textbf{Noise:} Sounds of your actions alert nearby \textbf{bees or Lethai−ar wardens}.
ℹtem \textbf{Trace:} Evidence of your passage marks your route for \textbf{Oath−Keepers or Waker spies}.
ℹtem \textbf{Delay:} A brief but meaningful setback costs you \textbf{time or safe passage through a grove}.
ℹtem \textbf{Supply Strain:} Mark $+1$ segment on a relevant \textbf{resource timer}.
\end{itemize}

\paragraph{Moderate Setbacks (2 SB)}
\begin{itemize}
ℹtem \textbf{Alarm Raised:} \textbf{Local Leopard Society or Green Council druid} becomes aware and begins responding.
ℹtem \textbf{Position Lost:} You lose advantageous ground/cover/stealth due to \textbf{thorn−spirit or fog from the Crimson Falls}.
ℹtem \textbf{Foe Appears:} A \textbf{Waker cell or Oath−Keeper} arrives on scene.
ℹtem \textbf{Gear Trouble:} A piece of equipment becomes \textbf{Compromised/Neglected}.
ℹtem \textbf{Lock/Barrier:} A simple obstacle now requires a test to overcome.
\end{itemize}

\paragraph{Serious Trouble (3 SB)}
\begin{itemize}
ℹtem \textbf{Reinforcements:} Additional \textbf{Leopard Society, Earth Society guardians, or Lethai−ar archers} arrive.
ℹtem \textbf{Key Gear Breaks:} A crucial tool/weapon becomes temporarily unusable.
ℹtem \textbf{Major Twist:} The situation fundamentally changes −−− \textbf{Ukwe speaks/Waker poison spreads/web−law broken}.
ℹtem \textbf{Rail Tick:} Advance a relevant campaign/front timer by 1 segment.
ℹtem \textbf{Condition Applied:} Mark \textbf{Fatigue 1/Harm 1/Condition} appropriate to fiction.
\end{itemize}

\paragraph{Major Turns (4+ SB)}
\begin{itemize}
ℹtem \textbf{Trap Springs:} A prepared danger activates with full effect.
ℹtem \textbf{Authority Arrival:} \textbf{Leopard King, Ukwe's Voice, or Crimson Valley First Weaver} intervenes.
ℹtem \textbf{Scene Shift:} The environment changes dramatically −−− \textbf{\textit{nkisi} shatter/honey turns to blood/Ukwe's root breaches the ground}.
ℹtem \textbf{Patron Omen:} Divine/arcane forces take notice −−− \textbf{omen appears/blessing lost/curse manifests}.
ℹtem \textbf{Narrative Pivot:} The story takes an unexpected turn that reframes objectives.
\end{itemize}

\paragraph{Region−Specific SB Options}
\begin{itemize}
ℹtem \textbf{Sekogo (Jungle Law):} Bees form words, tree bark weeps sap, the Ukwe's root trembles.
ℹtem \textbf{Sekogo (Crimson Valley):} Silk cords tighten, canyon walls bleed rust, the web−law recites itself.
ℹtem \textbf{Sekogo (Society Conflict):} Masks shift expressions, bark strips burn without flame, a ninth nail appears in every \textit{nkisi}.
\end{itemize}

\subsection*{Pursuit Ladder (Near / Mid / Far)}
\begin{itemize}
ℹtem \textbf{Advance/Withdraw (Wits + Survival):} On a hit, shift one band. On a strong hit, also impose \emph{Thorn Wall}.
ℹtem \textbf{Cut the Web (Presence + Sway):} On a hit, freeze range for one exchange; if a silk cord is in your possession, $+1$ die.
ℹtem \textbf{Weather the Hunt (Resolve + Spirit):} On a hit in \emph{Leopard Society Hunt}, you pick who is targeted; on a miss, the GM does.
\end{itemize}

\subsection*{Boss Seeds (Root & Mask)}
\begin{itemize}
ℹtem \textbf{The Root−King (Nine Broken Promises)} −−− \emph{Phases}: Promise Calls $→$ Orchestrated Famine $→$ Grove Coup. \emph{Cracks}: burn a bark strip, expose a false memory, survive a harvest without honey. \textit{The ninth root follows you home.}
ℹtem \textbf{The Mask that Walks (Leopard's Shadow)} −−− \emph{Phases}: Midnight Hunts $→$ Procession of Masks $→$ Skin−Baptism of a Warden. \emph{Cracks}: relic leopard claw, survivor testimony, the mask turning against them −−− and the ninth mask is always the Hollow's own. 
ℹtem \textbf{Mami Wata's Wrath (Water Spirit)} −−− \emph{Phases}: Pool Rises $→$ Crimson Falls Reversed $→$ Memory of the Drowned. \emph{Cracks}: offer a true memory of loss, speak the name of a slaver ship she sank, return a stolen mirror to her pool. \textit{The ninth memory she takes is your own.}
\end{itemize}

\begin{tcolorbox}[colback=black!3,colframe=black!40!white,title={Patronage & Power}]
Among Sekogo, power flows through witnessed promises, honey, and the consent of the jungle. The Green Council rules by consensus, but the Leopard Society restores balance, the Earth Society keeps the Calendar of Witness, and the nganga negotiate with spirits. No single faction holds all the power −−− and that is by design.

\textbf{For the GM:}  
Power in Sekogo society turns on broken promises and witnessed oaths. Patronage often takes the form of favors, memories, or shards −−− each carrying a future cost. To emphasize this:
\begin{itemize}
ℹtem Tie rewards to objects that create reciprocal bonds (leopard claws, Ukwe leaves, \textit{nkisi} shards, \textit{mukanda} strips) that can be called in at inconvenient moments.
ℹtem Let rival societies offer conflicting judgments, forcing players to choose whose law to follow −−− and whose ninth clause to honor.
ℹtem Use the honey−orchards, druid circles, and silk bridges as arenas for social contests, where a gift of honey is a treaty and a broken promise is a death sentence.
ℹtem Introduce trickster spirits (Ananse's Child) to complicate simple resolutions −−− they are excellent wild cards, and their ninth bell is always the Hollow's cue.
ℹtem Remember that Mami Wata is not evil; she is ancient. She collects memories, not lives. A bargain with her may be the only way to escape a blood−feud −−− but the ninth memory she asks for will be the one you cannot spare.
\end{itemize}
In Sekogo society, the jungle remembers everything −−− and the Ukwe forgets nothing. The \textit{Mwitu} waits at the edge of the ninth root, and the ninth bell of Ananse's Child rings only when a promise is already broken.
\end{tcolorbox}

% Index entries
∈dex{Sekogo|see{Green Council}}
∈dex{Theme generator}
∈dex{Green Council generator}
∈dex{Places!Sekogo}
∈dex{People!Sekogo}
∈dex{Complications!Sekogo}
∈dex{Rewards!Sekogo}
∈dex{Leopard Society}
∈dex{Earth Society}
∈dex{Nganga}
∈dex{Ukwe (Speaking Tree)}
∈dex{Crimson Valley}
∈dex{Lethai−ar}
∈dex{Web−law}
∈dex{Marks (tattoos)}
∈dex{Honey−orchards}
∈dex{Minkisi (power objects)}
∈dex{Wakers vs Conservators}
∈dex{Ananse (trickster)}
∈dex{Mami Wata}
∈dex{Ninki Nanka}
∈dex{Mwitu (Hollow of the forest)}
∈dex{Mukanda (promise−cords)}
∈dex{Local Patrons!Sekogo}
\newpage

% -------------------- LETHAI-AR OF THE CRIMSON VALLEY --------------------
\chapter{Lethai-ar of the Crimson Valley}
\emph{Silk Bridges and Web‑Law.}

\noindent
Kofi whispers: \textit{“The Crimson Valley is a wound of red stone, carved by a river that runs rust‑coloured after every rain. The Lethai‑ar guard the bridges. Their web‑law is unwritten. Cross a bridge without a witness, and the silk will remember your weight – and your name.”}

\noindent
The Crimson Valley is a deep jungle rift between Sekogo and Oshiira – a scar of red stone, carved by a river that runs rust‑coloured after every rain. Here, the Lethai‑ar guard their borders with fierce neutrality, their chestnut‑brown skin blending with the canyon walls. Silk bridges span the gorges, woven from spider‑sap and oath‑hair. The web‑law is unwritten, carried in the memory of each warden; it keeps both empires from using the valley as a shortcut to war. Every promise is a strand. Every broken oath is a tremor through the whole valley.

\vspace{0.5em}
\noindent\textit{“Never step onto a silk bridge without first naming a witness. The bridge will remember your footsteps. If you fall, it will ask why no one was watching.”} \\
— Silk‑Warden Anara

\vspace{1em}
\begin{almanac}
\textbf{What do people eat for breakfast?} Rice porridge and fermented tea – the rice from terraces cut into canyon walls, the tea aged in clay pots buried beneath the bridges. The tea is bitter. The bitterness is for memory.

\textbf{What does a child learn before they can walk?} To read a Mark. Promise, Debt, Witness, Apology, Interest, Cut, Forgive, and the ninth – the empty space where a Mark was erased. The erased Mark is a broken promise. The cut is a scar.

\textbf{What is the first thing you notice about a stranger?} Their Marks. Fresh Marks mean new promises. Faded Marks mean unpaid debts. Cut Marks mean exile. The exiled are not spoken to. They are not remembered.

\textbf{What is the most common crime?} Bridge‑cutting – which is war. A bridge is a promise between two shores. Cut it, and the valley tightens its web.

\textbf{What does no one talk about but everyone knows?} The valley is not a place. It is a creature – a sleeping serpent whose spine is the canyon walls, whose scales are the silk bridges, whose heart is the rust‑coloured river. The Lethai‑ar are its parasites. They do not know it is asleep.

\textbf{Where do people go when they die?} To the river. Bodies wrapped in silk and cast into the rust‑coloured water. The river carries them downstream, into the Ukwe’s territory. The Ukwe drinks. The Ukwe remembers.

\textbf{How do you apologize for a serious wrong?} You offer a Mark – a promise cut from your own skin. It is pressed into the wronged party’s ledger. If they accept, the debt is transferred. If they refuse, the Mark is buried. Buried Marks are debts that never die.

\textbf{What is the most important festival?} The Weaving of the New Bridge (spring). The silk is harvested from the canyon’s spiders. The bridge is woven in a single night. It is named after a promise. The promise is always the same: the valley will hold.

\textbf{What is the secret that could destroy a powerful person?} The river’s rust colour is not iron. It is blood – the blood of the serpent that sleeps beneath the canyon. The serpent is wounded. The wound has been bleeding for nine centuries. The Lethai‑ar know. They do nothing. If the serpent wakes, the law ends.
\end{almanac}

\newpage
\section{Lethai−ar — ``The Crimson Valley''}
łabel{chap:lethai−ar}

\textbf{Key Demographic: } Lethai−ar
\textbf{Capital City: } Ananse

\subsection*{Quick Reference}
\textbf{Tagline:} \textit{The bridge remembers every step. The silk knows every promise.}

\begin{quote}
  \textbf{Silk−Warden (Elite):}  
  ``The web−law is not written on parchment. It is woven into every bridge, every strand, every knot. When a promise is made, the valley remembers. When a promise is broken, the silk tightens. We are not judges. We are the loom. The spider spins; we only follow the thread.''
\end{quote}

\begin{quote}
  \textbf{Mark−Cutter (Commoner):}  
  ``Each line on my skin is a promise I witnessed. Each line is a story I carry. When I die, they will cut these marks from my body and bury them at the Ukwe's roots. Then the tree will remember what I witnessed. I hope it is enough. The thread does not break; it only passes to the next hand.''
\end{quote}

\textbf{Genre / Mood:} High−fantasy legal drama, spidersilk noir, ancestral witness.

\textbf{Starting Location:} The central silk bridge, where the gorge is deepest and the wind carries the smell of rust. A Silk−Warden blocks your path, her mask painted with roots: ``The web−law demands a toll. Not in coin – in a promise. Swear that you will not use this valley to carry steel to Ashaan, or turn back now. The thread is watching. So is the spider.''

\vspace{2mm}
\begin{tcolorbox}[colback=black!3,colframe=blue!60!black,title={Theme & Atmosphere}]
The Crimson Valley is a deep jungle rift between Sekogo and Oshiira – a wound of red stone, carved by a river that runs rust‑coloured after every rain. Here, the Lethai−ar guard their borders with fierce neutrality, their chestnut‑brown skin blending with the canyon walls. Silk bridges span the gorges, woven from spider‑sap and oath‑hair, strong enough to hold a caravan and light enough to sway in the monsoon wind. The web−law is unwritten, carried in the memory of each warden; it keeps both empires from using the valley as a shortcut to war.

Unlike their northern kin who fell from grace, the Lethai−ar are adopted children of the old powers – not fallen, but chosen to carry memory. Each can name the specific mountain, river, or grove where their ancestor left the spirit world to become mortal. Their skin‑tattoos (Marks) are woven with red ochre and jungle sap, each line a promise witnessed, each curve a story recorded. When a Lethai−ar dies, their Marks are cut from the skin and buried at the Ukwe's roots, so the tree may remember what they witnessed.

The valley itself is alive. The silk remembers every footfall; the river whispers the names of those who have crossed; and in the oldest bridges, the spider – Inaea, the weaver of bonds – still spins. The Lethai−ar do not worship her, but they honour the thread. They say that every promise is a strand in her web, and that a broken oath sends a tremor through the whole valley.

\vspace{2mm}
\textbf{What you notice first:} The whisper of silk ropes in the wind. The smell of red ochre and damp stone. The way every Lethai−ar touches their left wrist when they speak a name – a gesture of witness.

\textbf{The one rule that kills newcomers:} Never step onto a silk bridge without first naming a witness. The bridge will remember your footsteps. If you fall, it will ask why no one was watching.
\end{tcolorbox}

\subsection*{Regional Mood: The Weight of Witness}

Power here flows from witnessed promises, silk strands, and the web−law's quiet authority. The valley is neutral ground – Oshiiran grain barges do not float here, and Ashaani armies dare not march. The Lethai−ar trade silk for honey with Sekogo druids and stonesong for poison with Ngomebe masons. Their Marks are not decoration; they are a public memory, written on skin, read by anyone who knows the phase of the crimson moon. A Lethai−ar who bears a Mark they cannot fulfil loses standing; a warden who cuts a bridge without cause is exiled. But the true authority is the web itself – the invisible threads that connect every oath ever sworn in the valley. The Lethai−ar do not create the web. They tend it.

\subsection*{Prominent NPCs of the Crimson Valley}
łabel{sec:lethaiar−npcs}

\begin{tabular}{|p{0.2\textwidth}|p{0.25\textwidth}|p{0.3\textwidth}|p{0.15\textwidth}|}
ℎline
\textbf{Name} & \textbf{Role/Title} & \textbf{Motivation} & \textbf{Rumoured Quirk} \\
ℎline
Warden Anara & Silk‑warden of the central span & Keep the Oshiiran grain barges from crossing & Her left arm is covered in Marks from wrist to shoulder. She will not say what they mean. \\
ℎline
Mark‑Cutter Tebe & Tattooist who records promises & Witness a treaty between Sekogo and Ngomebe & She uses a needle of silk‑fiber and ink made from crushed crimson stone. The ink glows under the full moon. \\
ℎline
Bridge‑Keeper Solan & Guardian of the eastern crossing & Prevent Ashaani scouts from using the valley & He has not slept in nine days. He says the bridge whispers when someone lies. \\
ℎline
Elder Mira & Memory‑keeper of the valley's spirits & Record the name of every ancestor who walked the spirit path & She carries a bundle of silk cords, each tied to a different knot. The knots are the names. \\
ℎline
The Spider's Child & Tulkani trader of oath‑silk (whisperer of Inaea) & Buy a bolt of crimson silk for a story – and to listen & She is not Lethai−ar but has lived in the valley for thirty years. The wardens tolerate her. She never closes her eyes. \\
ℎline
The Unmarked & A Lethai−ar who refused any Marks & An outcast who lives in the gorge below the bridges & Her skin is blank. She is invisible to the web‑law – and to the valley's spirits. She carries a single knotted thread. \\
ℎline
\end{tabular}

\subsection*{Names of the Crimson Valley}
łabel{sec:lethaiar−names}

Lethai−ar names are soft, with rolling consonants and vowel endings. Surnames often denote a bridge, a gorge, or a spirit‑ancestor. Epithets refer to Marks or witnessed deeds.

\begin{tabular}{|p{0.25\textwidth}|p{0.25\textwidth}|p{0.2\textwidth}|p{0.2\textwidth}|}
ℎline
\textbf{First Names (Male)} & \textbf{First Names (Female)} & \textbf{Surnames} & \textbf{Epithets} \\
ℎline
Solan & Anara & Silkweave & \textit{the Unbroken} \\
ℎline
Tebe & Mira & Gorge‑born & \textit{the Ninth Witness} \\
ℎline
Liam & Nira & Root‑memory & \textit{the Marked} \\
ℎline
Veth & Liana & Spider‑kin & \textit{the Bridge‑Heart} \\
ℎline
Renn & Sora & Crimson‑touch & \textit{the Unweaver} \\
ℎline
Thorn & Aria & Oath‑strand & \textit{the Still Watcher} \\
ℎline
\end{tabular}

\paragraph*{(Bridge/Gorge/Roots)} Silk bridge spanning a crimson gorge – ropes sing in the wind; a Lethai−ar village built into the canyon wall; the Ukwe's roots clutching the water below the central span.

\section*{Spades — Places}
łabel{sec:lethaiar−places}
\begin{enumerate}
\setcounter{enumi}{1}
ℹtem \textbf{Central Silk Bridge} – A wide span woven from crimson and grey strands. The bridge hums when a promise is made on it.
  ℎfill \textit{The pitch tells the warden if the promise is true. If the thread vibrates at the ninth frequency, the spider is listening.} \texttt{[WITNESS]}
ℹtem \textbf{Mark‑Cutter's Studio} – A cave with a stone table and pots of ochre.
  ℎfill \textit{The walls are carved with Marks that were never cut into skin. They are stories the cutter refused to witness. The stone is warm to the touch.} \texttt{[CRAFT]}
ℹtem \textbf{Ukwe's Riverside Root} – A massive root that dips into the crimson river.
  ℎfill \textit{The root is warm. Touch it, and the tree will show you one forgotten promise from a dead Lethai−ar. The water ripples in response.} \texttt{[MEMORY]}
ℹtem \textbf{Abandoned Silk Span} – A broken bridge hanging over the gorge. The cut ends are frayed.
  ℎfill \textit{The cut was not clean – someone broke an oath mid‑crossing. The strands still twitch.} \texttt{[DEATH]}
ℹtem \textbf{Spirit‑Anchor Stone} – A boulder where the first ancestor left the spirit world.
  ℎfill \textit{The stone is smooth. It is warm even in winter. Lethai−ar leave offerings of red ochre here – and sometimes a single knotted thread.} \texttt{[RITUAL]}
ℹtem \textbf{Warden's Watchtower} – A stone needle built into the canyon wall, with a signal mirror.
  ℎfill \textit{The mirror flashes in code. The code is the web‑law. No outsider has broken it. The code changes when the spider dreams.} \texttt{[AUTHORITY]}
ℹtem \textbf{Hollow of the Unmarked} – A cave below the eastern bridge, hidden by a waterfall.
  ℎfill \textit{The cave has no silk. The Unmarked lives here, weaving nets from her own hair. Her hair is knotted.} \texttt{[MYSTERY]}
ℹtem \textbf{Crimson River Ford} – A shallow crossing where the water is clear, not red.
  ℎfill \textit{The ford is sacred. No oath sworn here can be broken. The water would turn black – and the spider would feel it.} \texttt{[OATH]}
ℹtem \textbf{Elder's Cord‑House} – A longhouse where the memory‑knots are stored.
  ℎfill \textit{Each cord is a name. The cords are hung from the ceiling. They move in no wind. When a name is forgotten, a cord falls.} \texttt{[MEMORY]}
ℹtem[\card{♠}{J}] \textbf{The Spider's Stall} – A Tulkani merchant's tent at the valley's southern exit. She sells silk to anyone.
  ℎfill \textit{The price is a story. The story must be about a thread that held – or a thread that broke. She never blinks.} \texttt{[TRADE]}
ℹtem[\card{♠}{Q}] \textbf{Bridge of the Ninth Strand} – A legend: a bridge woven from a single hair.
  ℎfill \textit{It appears only at dawn on the first day of spring. Those who cross it gain a new name – and a new thread to their web.} \texttt{[LEGEND]}
ℹtem[\card{♠}{K}] \textbf{The Cut Stone} – A boulder split by an oath‑breaker's fall.
  ℎfill \textit{The split is shaped like a human body. The stone weeps water when the oath‑breaker's name is spoken. The water tastes of rust.} \texttt{[DEATH]}
ℹtem[\card{♠}{A}] \textbf{The First Bridge (Myth)} – Where the first Lethai−ar wove the first strand.
  ℎfill \textit{It is invisible. Only those who have witnessed nine true promises can see it. The spider sits at the far end, waiting.} \texttt{[MYTH]}
\end{enumerate}

\paragraph*{(Warden/Cutter/Trader)} Silk‑warden with a mask painted with the Ukwe's roots; Mark‑cutter with a needle of silk‑fiber; Tulkani spider‑trader with a bundle of oath‑cords.

\section*{Hearts — People & Factions}
łabel{sec:lethaiar−people}
\begin{enumerate}
\setcounter{enumi}{1}
ℹtem \textbf{Silk−Warden} – Mask carved from crimson wood, painted with oath‑names.
  ℎfill \textit{She never blinks when standing on the bridge. The bridge does not forgive distraction. She counts her breaths in eights.} \texttt{[AUTHORITY]}
ℹtem \textbf{Mark−Cutter} – Hands stained red with ochre, eyes that measure skin.
  ℎfill \textit{She asks each client: ``What promise do you want the world to remember?'' The answer becomes a Mark – and a thread in the web.} \texttt{[CRAFT]}
ℹtem \textbf{Bridge−Keeper} – Old, grey‑haired, missing three fingers from a snapped strand.
  ℎfill \textit{He carries a whistle made from a spider's fang. The whistle calls the bridge to tighten – or to loosen. He never uses it.} \texttt{[WITNESS]}
ℹtem \textbf{Elder (Memory‑Keeper)} – Bundles of silk cords hang from her belt.
  ℎfill \textit{She knows every name carved into every root of the Ukwe. She has forgotten her own. The cords remember it for her.} \texttt{[MEMORY]}
ℹtem \textbf{Tulkani Spider (Trader)} – Wears a shawl of crimson silk, speaks in riddles.
  ℎfill \textit{She pays for silk with stories. The stories become new strands. The bridge grows. She never closes her eyes – she says the spider does not sleep.} \texttt{[TRADE]}
ℹtem \textbf{The Unmarked} – A Lethai−ar without Marks, living in the hollow.
  ℎfill \textit{She whispers to the waterfall. The waterfall whispers back. No one knows what they say. Her thread is her own.} \texttt{[MYSTERY]}
ℹtem \textbf{Youth Seeking a Mark} – A young Lethai−ar on their first witness journey.
  ℎfill \textit{They carry a single strand of undyed silk. When they witness a promise, they will dye it crimson. The dye is their own blood.} \texttt{[COMMUNITY]}
ℹtem \textbf{Oath‑Breaker's Ghost} – A translucent figure pacing the broken bridge.
  ℎfill \textit{It repeats the words of the promise it broke. The words are getting quieter. When they stop, the thread will snap.} \texttt{[DEATH]}
ℹtem \textbf{Sekogo Honey Trader} – A druid with a jar of dark honey, trading for silk poison.
  ℎfill \textit{He offers the honey as a gift. The wardens test it by feeding a drop to the bridge. The bridge sways if it is poisoned – or if the gift is false.} \texttt{[TRADE]}
ℹtem[\card{♥}{J}] \textbf{Ngomebe Stone‑Singer} – A mason with a tuning fork, here to buy silk for sledge‑wax seals.
  ℎfill \textit{He taps the bridge with his fork. The note tells him how many promises it has witnessed. The note is always an octave below the spider's hum.} \texttt{[CRAFT]}
ℹtem[\card{♥}{Q}] \textbf{The First Ancestor's Echo} – A voice heard at the Spirit‑Anchor Stone at midnight.
  ℎfill \textit{It asks: ``Why did you leave the spirit world?'' The answer must be a promise. The spider listens.} \texttt{[RITUAL]}
ℹtem[\card{♥}{K}] \textbf{The Spider's Daughter} – A half‑Lethai−ar child raised by the Tkanki trader. She has no Marks.
  ℎfill \textit{Her skin is blank. She is learning to read knots – and to spin her own thread. She never asks twice.} \texttt{[COMMUNITY]}
ℹtem[\card{♥}{A}] \textbf{The Ninth Warden (Legend)} – A warden who cut her own bridge to stop an invasion.
  ℎfill \textit{She is said to walk the valley at dusk, carrying a coil of silk. She offers one free crossing – but the crossing leads to the spider's web.} \texttt{[LEGEND]}
\end{enumerate}

\paragraph*{(Strand/Challenge/Fall)} A silk strand snaps – lose 1 Supply to repair; a Lethai−ar challenges you to recite the web‑law; a traveler falls from the bridge – the party must witness their last words.

\section*{Clubs — Complications/Threats}
łabel{sec:lethaiar−complications}
\begin{enumerate}
\setcounter{enumi}{1}
ℹtem \textbf{Snapped Strand} – A rope on a bridge breaks. Lose 1 Supply to repair.
  ℎfill \textit{The break is clean. Someone cut it with a blade. The blade was from Ashaan – or from a broken promise.} \texttt{[DISASTER]}
ℹtem \textbf{Web−Law Challenge} – A warden blocks the bridge: ``Recite the ninth clause.''
  ℎfill \textit{The ninth clause is unwritten. Each warden carries a different version. You must guess. The spider knows the true one.} \texttt{[SOCIAL]} \texttt{[RIDDLE]}
ℹtem \textbf{Fall from the Bridge} – A traveler slips. The bridge asks: ``Who witnessed this?''
  ℎfill \textit{If no one answers, the traveler's name is erased from the Ukwe's roots – and the spider weaves a new strand.} \texttt{[DEATH]} \texttt{[WITNESS]}
ℹtem \textbf{Oshiiran Surveyors} – Grain auditors arrive with a map, claiming the valley's water rights.
  ℎfill \textit{They carry a treaty stamped with the Widow's seal. The seal is a forgery. The bridge hums in warning.} \texttt{[AUTHORITY]} \texttt{[POLITICS]}
ℹtem \textbf{Ashaani Scouts (Disguised)} – Veiled Sister agents posing as spice merchants.
  ℎfill \textit{They do not touch their left wrist when speaking a name. The wardens notice. The thread tightens.} \texttt{[STEALTH]} \texttt{[POLITICS]}
ℹtem \textbf{Mark−Cutter's Witness Called} – A Lethai−ar's Marks begin to glow red. A promise is overdue.
  ℎfill \textit{The witness is not present. The party is the closest witness. The duty transfers – as does the thread.} \texttt{[WITNESS]} \texttt{[COMMUNITY]}
ℹtem \textbf{Bridge Poisoning} – Someone smears a solvent on the silk strands. The bridge begins to fray.
  ℎfill \textit{The solvent is made from Sekogo honey and Ashaani dust. The Spider knows who sold it. The threads are weeping.} \texttt{[MAGIC]} \texttt{[DISASTER]}
ℹtem \textbf{Elder's Memory Fails} – The cord‑keeper forgets a name. The knot unravels.
  ℎfill \textit{The forgotten name is an old promise. The Ukwe's root trembles. Something is coming – or someone.} \texttt{[MEMORY]} \texttt{[DEATH]}
ℹtem \textbf{The Unmarked's Bargain} – The outcast offers to guide the party through the gorge. Her price is a promise.
  ℎfill \textit{The promise: ``Do not speak my name after I die.'' Refuse, and she leaves you to the spiders. Accept, and she ties a knot.} \texttt{[SOCIAL]} \texttt{[WITNESS]}
ℹtem[\card{♣}{J}] \textbf{Crimson River Floods} – Rain turns the ford to a torrent. The bridge is the only crossing.
  ℎfill \textit{The water is red. Faces appear in the foam. They are oath‑breakers, still trying to cross. The spider watches.} \texttt{[WEATHER]} \texttt{[DEATH]}
ℹtem[\card{♣}{Q}] \textbf{Warden Schism} – Two wardens disagree on the web‑law. They cut each other's bridges.
  ℎfill \textit{The valley is now divided. The party must choose which side to cross – or stay stuck forever. The spider waits.} \texttt{[AUTHORITY]} \texttt{[POLITICS]}
ℹtem[\card{♣}{K}] \textbf{The First Bridge Appears} – The mythical bridge manifests at dawn. Crossing it gives a new name.
  ℎfill \textit{The new name is a promise. The promise is to the First Ancestor. The ancestor will collect – and the spider will weave.} \texttt{[MYTH]} \texttt{[MAGIC]}
ℹtem[\card{♣}{A}] \textbf{The Ninth Warden Returns} – The legend appears, offering a free crossing.
  ℎfill \textit{The crossing leads to the spider's web. The Hollow asks: ``What promise did you break?'' Answer poorly, and you become a new strand.} \texttt{[DEATH]} \texttt{[LEGEND]}
\end{enumerate}

\paragraph*{(Cord/Mask/Shard)} Silk truth‑cord (enforce any oath); warden's mask (safe passage on any bridge for one crossing); Ukwe root shard (one question about a forgotten promise).

\section*{Diamonds — Rewards/Leverage}
łabel{sec:lethaiar−rewards}
\begin{enumerate}
\setcounter{enumi}{1}
ℹtem \textbf{Silk Truth‑Cord} – A braided strand dyed crimson. Use it to enforce any oath.
  ℎfill \textit{The cord tightens when the oath is broken. It will not stop until the promise is kept – or the spider releases it.} \texttt{[OATH]}
ℹtem \textbf{Warden's Mask} – A wooden mask painted with roots. Safe passage on any bridge for one crossing.
  ℎfill \textit{The mask has no eye holes. You must trust the bridge to see for you. The bridge sees through the spider's eyes.} \texttt{[TRAVEL]}
ℹtem \textbf{Ukwe Root Shard} – A piece of the tree's root that grows near the river.
  ℎfill \textit{Touch it, and the Ukwe will answer one question about a forgotten promise. The answer is always true – and always a thread.} \texttt{[MEMORY]}
ℹtem \textbf{Mark‑Cutter's Needle} – A silk‑fiber needle. Use it to record a witnessed promise as a temporary Mark.
  ℎfill \textit{The Mark fades after one season. The story remains – and the spider remembers.} \texttt{[CRAFT]}
ℹtem \textbf{Bridge‑Keeper's Whistle} – A fang whistle. Blow it to tighten any silk bridge for one scene.
  ℎfill \textit{The whistle works once. The bridge will remember the favour – and the spider will know.} \texttt{[AUTHORITY]}
ℹtem \textbf{Elder's Memory Cord} – A knotted silk cord that contains a forgotten name.
  ℎfill \textit{Untie the knot to learn the name. The name is a promise. Fulfil it, or the cord will retie itself – and the spider will watch.} \texttt{[MEMORY]}
ℹtem \textbf{Spider's Story‑Strand} – A bolt of crimson silk bought with a story.
  ℎfill \textit{The silk is soft. It glows under moonlight. Trade it for safe passage through Tkanki camps – or give it to the spider.} \texttt{[TRADE]}
ℹtem \textbf{The Unmarked's Net} – A net woven from human hair.
  ℎfill \textit{Throw it to catch a spirit. The net can hold a bound spirit for one hour. The spirit will bargain for freedom – with a promise.} \texttt{[MAGIC]}
ℹtem \textbf{Crimson Ochre Pot} – A clay pot of red pigment. Use it to Mark a promise on stone or skin.
  ℎfill \textit{The ochre never dries. The promise remains witnessable forever – until the spider cuts the thread.} \texttt{[CRAFT]}
ℹtem[\card{♦}{J}] \textbf{Ninth Strand Fragment} – A piece of the legendary bridge. Once per campaign, cross any gap as if on solid ground.
  ℎfill \textit{The fragment crumbles after use. The first bridge will remember you – and the spider will ask for a toll.} \texttt{[LEGEND]}
ℹtem[\card{♦}{Q}] \textbf{The Cut Stone's Water} – A vial of water that seeps from the oath‑breaker's boulder.
  ℎfill \textit{Drink it to remember a promise you broke. The memory is painful. The pain is the price – and the thread.} \texttt{[DEATH]}
ℹtem[\card{♦}{K}] \textbf{First Ancestor's Blessing} – A whispered promise from the Spirit‑Anchor Stone.
  ℎfill \textit{Once, you may reroll any failed roll involving a witnessed promise. The ancestor watches – and the spider weaves.} \texttt{[RITUAL]}
ℹtem[\card{♦}{A}] \textbf{The Ninth Warden's Coil} – A coil of silk said to belong to the legendary warden.
  ℎfill \textit{Use it to cut a bridge without angering the web‑law. The coil then unravels – and the spider's shadow falls on you.} \texttt{[LEGEND]}
\end{enumerate}

\section*{Quick use notes}
łabel{sec:lethaiar−quick−use}
\begin{itemize}
ℹtem Draw until you have all four suits: Spade = place, Heart = actor, Club = pressure, Diamond = leverage. Highest rank sets the main Timer (2–5 \(→\) 4, 6–10 \(→\) 6, J/Q/K \(→\) 8, A \(→\) 10).
ℹtem Diamonds are codified outcomes (cords, masks, shards) that change position rather than call for a roll.
ℹtem If any A appears, echo \textbf{valley‑omens} – silk cords humming in still air, the crimson river running clear, a bridge that sways without wind. The spider is near.
\end{itemize}

\subsection*{Additional Features (Cultural Mechanics)}
łabel{sec:lethaiar−mechanics}

\begin{itemize}
ℹtem \textbf{The Web‑Law:} The unwritten legal code of the Crimson Valley. Each warden carries a different clause in memory; together, they form a complete web. To cross a bridge, a traveler must answer a question about the web‑law. The answer changes with each warden. On a successful roll (Wits + Lore, DV 3), you may bypass the challenge; on a failure, you must offer a promise instead – and the spider will witness it.
ℹtem \textbf{Marks (Skin‑Tattoos):} Each Lethai−ar's Marks are a public record of promises witnessed. The Marks are cut from the skin after death and buried at the Ukwe's roots. To read a living Lethai−ar's Marks, one must know the phase of the crimson moon and the position of the valley's spirits. A character who studies a Lethai−ar's Marks for a scene may ask the GM one question about a promise that Lethai−ar has witnessed; the GM must answer truthfully – but the spider may also listen.
ℹtem \textbf{Spirit Ancestry:} Every Lethai−ar can name the specific spirit (mountain, river, grove) where their ancestor left the spirit world to become mortal. That spirit is their witness; they owe it a promise at birth, and that promise is never fully paid. Once per journey, a Lethai−ar may invoke their spirit ancestor to gain +1 die on a roll related to that spirit's domain (e.g., a river ancestor aids a crossing). The spider notes the favour.
ℹtem \textbf{Witness Economy:} Instead of tracking debt, the Crimson Valley uses \textbf{Witness} as currency. A character who witnesses a promise gains 1 Witness token. Tokens can be spent to: gain +1 die on a roll involving that promise, force an oath‑breaker to reroll a successful deception, or demand a favour from someone who witnessed the same promise. Witness tokens are lost when the promise is fulfilled or broken – and the spider watches the exchange.
ℹtem \textbf{Local Patrons (Syncretic Names):}
  \begin{itemize}
    ℹtem \textbf{The First Weaver (Inaea):} The spider who taught the first knot; mercy is a strand, not a cut. Her web is the valley's memory.
    ℹtem \textbf{The Serpent Judge (Isoka):} The fang that decides; a promise kept is a skin shed. She watches from the river stones.
    ℹtem \textbf{The Crimson Witness (The Witness):} The flame that sees through silk; no oath hidden from its light. The bridge glows when it passes.
    ℹtem \textbf{The Bridge Father (The Traveler):} The road that crosses the gorge; every crossing is a promise. He walks the strands at dusk.
  \end{itemize}
\end{itemize}

\subsection*{Lore Echoes (Cross‑Expansion Hooks)}
łabel{sec:lethaiar−lore}

\begin{itemize}
  ℹtem \textbf{Crimson Silk and the Way of Silk:} The silk of the Crimson Valley is prized across the continent for vow‑cords, well‑contracts, and bridge‑building. The Tkanki spider‑traders are the primary merchants; a single bolt can buy safe passage from Fhara to Dhahara. A Tkanki story‑seller may accept crimson silk as payment for a truth – but the spider will ask for a copy.
  ℹtem \textbf{The Ukwe's Roots and Sekogo:} The Ukwe's roots extend into the Crimson Valley. The Lethai−ar bury their dead Marks at the riverbank root, creating a unique relationship with the Speaking Tree. The Ukwe knows every promise ever made in the valley. A Sekogo druid may ask the Ukwe to witness a new oath, binding it with the weight of the valley – and the spider's thread.
  ℹtem \textbf{Web‑Law and the Covenant (Mykkiel):} The Covenant judges respect the web‑law as a valid legal system. A Lethai−ar truth‑cord is admissible in a Covenant court. Some judges have learned to read Marks. A Covenant justiciar may offer to translate a web‑law dispute into imperial law – but the translation always favours the spider.
  ℹtem \textbf{The Unmarked and the Hollow (Eastern Realms):} Lethai−ar who refuse all Marks become invisible to the web‑law – and to the Hollow. The Hollow cannot track them. Some Aelaerem seek out the Unmarked to learn the secret of blank skin. The Unmarked charge a promise: ``Do not seek me again.'' The spider does not visit them.
\end{itemize}

\subsection*{Lethai−ar Superstitions (burn for +1 die once)}
\begin{itemize}
  ℹtem \textbf{Bread to the Bridge:} Crumble a crust onto the first strand. Ignore the first \emph{Snapped Strand} penalty. \texttt{[SUPERSTITION]}
  ℹtem \textbf{Knot the Cord:} Tie a single knot in a truth‑cord. Cancel \emph{Web−Law Challenge} or take –1 die on your next negotiation (your choice). \texttt{[SUPERSTITION]}
  ℹtem \textbf{Name to the Root:} Whisper a dead warden's name into the Ukwe's root. Gain +1 Effect when crossing a bridge this scene, but mark Obligation to the First Ancestor – and the spider will know. \texttt{[SUPERSTITION]}
\end{itemize}

\subsection*{Thematic SB Spend Table}
łabel{sec:lethaiar−sb}

\paragraph{Minor Complications (1 SB)}
\begin{itemize}
ℹtem \textbf{Exposure:} Your actions draw unwanted attention from Silk‑Wardens or Mark‑Cutters.
ℹtem \textbf{Noise:} Sounds of your actions alert nearby bridge‑keepers or Tkanki traders.
ℹtem \textbf{Trace:} Evidence of your passage marks your route for Elders or the Unmarked.
ℹtem \textbf{Delay:} A brief but meaningful setback costs you time or safe crossing.
ℹtem \textbf{Supply Strain:} Mark +1 segment on a relevant resource timer.
\end{itemize}

\paragraph{Moderate Setbacks (2 SB)}
\begin{itemize}
ℹtem \textbf{Alarm Raised:} Local Silk‑Warden or Bridge‑Keeper becomes aware and begins responding.
ℹtem \textbf{Position Lost:} You lose advantageous ground/cover/stealth due to snapped strand or rising river.
ℹtem \textbf{Foe Appears:} An Oshiiran surveyor or Ashaani scout arrives on scene.
ℹtem \textbf{Gear Trouble:} A piece of equipment becomes Compromised/Neglected.
ℹtem \textbf{Lock/Barrier:} A simple obstacle now requires a test to overcome.
\end{itemize}

\paragraph{Serious Trouble (3 SB)}
\begin{itemize}
ℹtem \textbf{Reinforcements:} Additional Wardens, Mark‑Cutters, or bridge guards arrive.
ℹtem \textbf{Key Gear Breaks:} A crucial tool/weapon becomes temporarily unusable.
ℹtem \textbf{Major Twist:} The situation fundamentally changes – bridge cut / web‑law broken / First Bridge appears.
ℹtem \textbf{Rail Tick:} Advance a relevant campaign/front timer by 1 segment.
ℹtem \textbf{Condition Applied:} Mark \textbf{Fatigue 1/Harm 1/Condition} appropriate to fiction.
\end{itemize}

\paragraph{Major Turns (4+ SB)}
\begin{itemize}
ℹtem \textbf{Trap Springs:} A prepared danger activates with full effect.
ℹtem \textbf{Authority Arrival:} \textbf{Elder, Ninth Warden, or First Ancestor's Echo} intervenes.
ℹtem \textbf{Scene Shift:} The environment changes dramatically – \textbf{bridge collapses / river turns black / Ukwe root rises}.
ℹtem \textbf{Patron Omen:} Divine/arcane forces take notice – \textbf{omen appears / blessing lost / curse manifests}.
ℹtem \textbf{Narrative Pivot:} The story takes an unexpected turn that reframes objectives.
\end{itemize}

\paragraph{Region−Specific SB Options}
\begin{itemize}
ℹtem \textbf{Crimson Valley (Web‑Law):} Silk strands glow, truth‑cords tighten, the Ukwe's root hums a name.
ℹtem \textbf{Crimson Valley (Marks):} Tattoos itch, ochre stains spread, a dead Mark appears on living skin.
ℹtem \textbf{Crimson Valley (Bridges):} Ropes whisper, the gorge echoes with footsteps that are not there, a bridge sways toward the Hollow.
\end{itemize}

\subsection*{Pursuit Ladder (Near / Mid / Far)}
\begin{itemize}
ℹtem \textbf{Advance/Withdraw (Wits + Survival):} On a hit, shift one band. On a strong hit, also impose \emph{Swaying Bridge}.
ℹtem \textbf{Cut the Strand (Presence + Sway):} On a hit, freeze range for one exchange; if a truth‑cord is in your possession, +1 die.
ℹtem \textbf{Weather the Crossing (Resolve + Spirit):} On a hit in \emph{Fall from the Bridge}, you pick who catches the fallen; on a miss, the GM does.
\end{itemize}

\subsection*{Boss Seeds (Silk & Promise)}
\begin{itemize}
ℹtem \textbf{The Unweaver (Nine Broken Strands)} – \emph{Phases}: Snapped Cords \(→\) Orchestrated Falls \(→\) Bridge Coup. \emph{Cracks}: repair a broken strand, expose a false witness, survive a crossing without a promise.
ℹtem \textbf{The Silent Warden (Oath of the Unmarked)} – \emph{Phases}: Forgotten Promises \(→\) Procession of the Blank \(→\) Skin‑Baptism of an Elder. \emph{Cracks}: relic truth‑cord, survivor testimony, the Marks turning against them.
\end{itemize}

\begin{tcolorbox}[colback=black!3,colframe=black!40!white,title={Patronage & Power}]
Among the Lethai−ar, power flows through witnessed promises, silk strands, and the authority of the wardens. The web‑law is unwritten, but every warden carries a piece of it. Marks are public, stories are visible, and the Ukwe remembers everything. The spider – Inaea – does not rule, but she watches. Her threads are the bonds between families, the knots that hold the valley together, and the silent witness to every oath.

\textbf{For the GM:} Power in the Crimson Valley is measured in witnessed oaths and the trust of the bridge. Patronage often takes the form of truth‑cords, masks, or root shards – each tied to a specific promise. To emphasise this:
\begin{itemize}
ℹtem Tie rewards to objects that record or enforce promises (cords, masks, shards) that can be challenged, broken, or called in.
ℹtem Let rival wardens issue conflicting interpretations of the web‑law, forcing players to choose whose bridge to cross.
ℹtem Use the bridges, Mark‑cutter's studio, and the Ukwe's root as arenas for social contests, where a witnessed promise is a weapon and a broken strand is a death sentence.
ℹtem Remember that the spider is not a god – she is a presence. She does not speak, but the valley hums with her attention. When a promise is broken, the threads tremble.
\end{itemize}
In the Crimson Valley, the bridge remembers every step – and the spider forgets no promise.
\end{tcolorbox}

% Index entries
∈dex{Lethai−ar|see{Crimson Valley}}
∈dex{Theme generator}
∈dex{Crimson Valley generator}
∈dex{Places!Lethai−ar}
∈dex{People!Lethai−ar}
∈dex{Complications!Lethai−ar}
∈dex{Rewards!Lethai−ar}
∈dex{Silk−Warden}
∈dex{Mark−Cutter}
∈dex{Web−law}
∈dex{Truth−cord}
∈dex{Inaea (spider)}
∈dex{Ukwe (root)}
∈dex{Local Patrons!Lethai−ar}

\newpage

% ----------------------------------------------------------------------
% APPENDIX: TRAVEL & TRADE IN THE SOUTH
% ----------------------------------------------------------------------
\appendix
\chapter{Travel \& Trade in the South}

This appendix summarises the unique travel modes and trade goods of the Southern Realms. For core travel procedures (legs, clocks, roles), see the \textit{Fate’s Edge SRD}.

\section*{Travel Modes}
\begin{itemize}
  \item \textbf{River Legs (Oshiira, Crimson Valley)}: Canals, weigh‑houses, and canal mothers. Travel clocks advance when barges queue at locks.
  \item \textbf{Desert Legs (Ashaan)}: Water is the primary resource. A dry well hazard can end a journey.
  \item \textbf{Jungle Legs (Sekogo, Crimson Valley)}: Navigation by smell and sound; hidden Leopard Society patrols.
  \item \textbf{Highland Legs (Taharka)}: Thin air; breath‑counting; passes that close for storms.
  \item \textbf{Stone Legs (Ngomebe)}: Travel with a walking city; roller integrity and ancestor favour clocks.
\end{itemize}

\section*{Trade Goods \& Currencies}
\begin{itemize}
  \item \textbf{Oshiira}: Grain ledger receipts, Vermilion Corps stamps. Accepted in all Covenant ports.
  \item \textbf{Ashaan}: Spirit‑lamps (bound names), brass seals. Illegal in most Covenant jurisdictions.
  \item \textbf{Ngomebe}: Keystone shards, sledge‑wax. Bartered for iron and Aeler engineering.
  \item \textbf{Taharka}: Pass‑stones, horse‑bonds, coffee measures. Honour currency in the highlands.
  \item \textbf{Sekogo}: Honey‑jars, Ukwe leaves (witness tokens). Spirits accept them as offerings.
  \item \textbf{Crimson Valley}: Witness cords, warden’s masks. Accepted as proof of oath in Covenant courts.
\end{itemize}

\section*{Regional Hazards (d6)}
When travelling through the Southern Realms, the GM may roll on this table for a regional hazard:
\begin{enumerate}
  \item \textbf{Dust Storm (Ashaan)}: Navigation DV +1; mark Fatigue.
  \item \textbf{Canal Lock Sabotage (Oshiira)}: Delay 2 segments; bribes required.
  \item \textbf{Keystone Crack (Ngomebe)}: City stops; ancestor favour clock advances.
  \item \textbf{Landslide (Taharka)}: Pass blocked; detour adds 3 segments.
  \item \textbf{Spirit Demands Toll (Sekogo)}: Offer honey or a story; refusal causes spirit attack.
  \item \textbf{Bridge Web‑Law Challenge (Crimson Valley)}: Recite a clause or pay a promise.
\end{enumerate}

% ----------------------------------------------------------------------
% COLOPHON
% ----------------------------------------------------------------------
\chapter*{Colophon}
\addcontentsline{toc}{chapter}{Colophon}

This book was written on paper that had been a cargo manifest, a love letter, and a bandage. The ink is lampblack cut with salt water and the ash of a burnt confession. The binding is grey cord tied with nine knots. The ninth knot is not cut.

Kofi of the Salt Syndicate composed the preface under a borrowed lamp in a Mbaro weigh‑house, though he admits he may have embellished a few tariffs. The almanacs were gathered from traders, spies, and the occasional honest farmer. The generator tables were tested in the field – meaning we ran them while a Oath‑Keeper counted our breaths. They work.

If you find yourself in a market and a salt trader offers you a discount, check his seals. The ledger is open. The scales are balanced. The ninth page is waiting.

\vspace{1em}
\noindent\textit{“The road is a ledger. Every oasis is a credit, every broken wheel a debt.”} – Dusana of the Raven Road

\end{document}
